% Lesson 54 — Vocabulary: Words and expressions related to using ICTs to enhance work or learning. # Semaine / Période Chapitre Titre Date effective Visa de l'enseignant — Anglais Tle A — Cours signé OnBuch+
% Programme officiel MINESEC. Compiler avec : tectonic EN54-vocabulary-words-and-expressions-related-to-using-icts.tex
\newcommand{\DOCMATIERE}{Anglais}
\newcommand{\DOCNIVEAU}{Tle A}
\newcommand{\DOCMODULE}{Chapter 5 : Media and Communication}
\newcommand{\DOCLECON}{EN54}
\newcommand{\DOCTITRE}{Lesson 54 — Vocabulary: Words and expressions related to using ICTs to enhance work or learning. \# Semaine / Période Chapitre Titre Date effective Visa de l'enseignant}
% OnBuch+ — préambule « cours signés » ENRICHI (compilé avec tectonic / XeLaTeX)
% Système visuel « Pop » d'OnBuch+ : fond crème (jamais blanc), bordures encre
% épaisses, ombres « dures » décalées, titres Archivo Black en capitales.
% Version MATHÉMATIQUES : graphes (pgfplots/tikz), boîtes pédagogiques
% (définition, propriété, méthode, attention, le savais-tu, exercice résolu,
% exercice, TP, bilan), page de garde et signature OnBuch+ ancrée en bas.
%
% Macros attendues dans main.tex AVANT \input{preamble} :
%   \DOCMATIERE  \DOCNIVEAU  \DOCMODULE  \DOCLECON (numéro)  \DOCTITRE
\documentclass[11pt]{article}

\usepackage{fontspec}
\usepackage[english,french]{babel} % français = langue principale ; l'anglais via \eng{..} / textebox
\usepackage[a4paper,margin=2cm,top=3.4cm,bottom=2.8cm,headheight=1.4cm,footskip=1.3cm]{geometry}
\usepackage{amsmath,amssymb,amsfonts}
\usepackage{mathtools} % \xmapsto, \coloneqq... souvent utilisés par les modèles
\usepackage{cancel} % simplifications barrées (bilans de forces)
% \abs{z} (module, valeur absolue) et \norm{u} (norme), utilisés par les modèles
\providecommand{\abs}{}\let\abs\relax\DeclarePairedDelimiter{\abs}{\lvert}{\rvert}
\providecommand{\norm}{}\let\norm\relax\DeclarePairedDelimiter{\norm}{\lVert}{\rVert}
\usepackage[table]{xcolor} % [table] : fournit aussi \rowcolors (alternance de lignes)
\usepackage{pagecolor}
\usepackage{tikz}
\usetikzlibrary{arrows.meta,positioning,calc,shapes.geometric,shapes.misc,shapes.arrows,decorations.pathmorphing,decorations.pathreplacing,decorations.markings,patterns,shadows,mindmap,trees,fit,backgrounds,matrix,intersections,angles,quotes}
\usepackage{pgfplots}
\usepackage[version=4]{mhchem} % équations nucléaires \ce{^{235}_{92}U}
\usepackage[european,siunitx]{circuitikz} % schémas électriques
\pgfplotsset{compat=1.18}
\usepgfplotslibrary{fillbetween,groupplots} % aires entre courbes (calcul intégral)
\usepackage{enumitem}
\usepackage{fancyhdr}
% [most] charge toutes les bibliothèques tcolorbox (listings, minted, raster,
% documentation...) et gonfle inutilement la mémoire TeX sur un document long
% (~300 boîtes) : on ne charge que ce qui est réellement utilisé.
\usepackage[skins,breakable]{tcolorbox}
\usepackage{graphicx}
\usepackage{colortbl}
\usepackage{booktabs}
\usepackage{tabularx}
\usepackage{diagbox} % case d'en-tête diagonale des tableaux à double entrée
% Colonnes extensibles centrée / gauche / droite (C, L, R), souvent utilisées par les modèles
\newcolumntype{C}{>{\centering\arraybackslash}X}
\newcolumntype{L}{>{\raggedright\arraybackslash}X}
\newcolumntype{R}{>{\raggedleft\arraybackslash}X}
\usepackage{array}
\usepackage{multicol}
\usepackage{multirow}
\usepackage{siunitx}
% parse-numbers=false : accepte aussi \SI{2,5 \times 10^{-3}}{...}
\sisetup{locale=FR,output-decimal-marker={,},group-minimum-digits=4,parse-numbers=false}
\DeclareSIUnit\ohm{\ensuremath{\symup{\Omega}}} % U+2126 absent de la police
\DeclareSIUnit\division{div} % réglages d'oscilloscope (ms/div, V/div)
\DeclareSIUnit\tr{tr} % tours (vitesse de rotation, ex. \SI{1500}{\tr\per\minute})
\DeclareSIUnit\molar{M} % concentration molaire (ex. \SI{3}{\molar}, \SI{1}{\micro\molar})
\DeclareSIUnit\an{an} % année (le modèle confond parfois avec \year, un compteur TeX) — ex. \SI{5}{\kilo\meter\per\an}
\DeclareSIUnit\FCFA{FCFA} % franc CFA (ex. \SI{500}{\FCFA\per\kilo\gram})
\DeclareSIUnit\degreeBrix{\textdegree Brix} % teneur en sucre d'un sirop/jus (réfractométrie)
\DeclareSIUnit\month{mois} % le modèle utilise parfois \month, un compteur TeX, comme unité de durée
\usepackage{qrcode}
% \degree : commande fréquemment utilisée par les modèles pour un angle ou une
% température en dehors de \SI{}{} (non fournie par les paquets chargés).
\providecommand{\degree}{^{\circ}}
% \qed (fin de démonstration), utilisé par les modèles sans amsthm
\providecommand{\qed}{\hfill$\square$}
% \barre : double barre des valeurs interdites dans un tableau de variations (tkz-tab)
\providecommand{\barre}{\ensuremath{\|}}
% environnement proof (amsthm non chargé) utilisé par les modèles
\ifdefined\proof\else\newenvironment{proof}[1][Démonstration]{\par\noindent\textit{#1.}\ }{\qed\par}\fi
\usepackage{lastpage}
\usepackage{hyperref}
\hypersetup{hidelinks}

% ============================================================
% Palette « Pop » OnBuch+ — mêmes hex que la classe P (plus_ui.dart)
% ============================================================
\definecolor{poporange}{HTML}{F59321}
\definecolor{popdark}{HTML}{E07A0C}
\definecolor{popcream}{HTML}{FFF6E8}
\definecolor{popink}{HTML}{1C1714}
\definecolor{poppurple}{HTML}{7A5AE0}
\definecolor{popgreen}{HTML}{2E9E6A}
\definecolor{poppink}{HTML}{E0457B}
\definecolor{popblue}{HTML}{2F7DE1}
\definecolor{popgold}{HTML}{C98A12}
\definecolor{popmuted}{HTML}{8A7F76}
\definecolor{popline}{HTML}{EADFD2}
\definecolor{popgray}{HTML}{8A7F76}
\colorlet{popgrey}{popgray}
% Alias tolérants (le modèle nomme parfois les couleurs différemment)
\colorlet{popwhite}{white}
\colorlet{popblack}{popink}
\colorlet{popred}{poppink}
\colorlet{popyellow}{popgold}
% Teintes claires (fonds de tableaux/encadrés) — le modèle nomme parfois
% les couleurs avec un suffixe L (ex. popblueL) sans les avoir définies.
\colorlet{poporangeL}{poporange!20}
\colorlet{popdarkL}{popdark!20}
\colorlet{poppurpleL}{poppurple!20}
\colorlet{popgreenL}{popgreen!20}
\colorlet{poppinkL}{poppink!20}
\colorlet{popblueL}{popblue!20}
\colorlet{popgoldL}{popgold!20}
\colorlet{popredL}{popred!20}
\colorlet{popgrayL}{popgray!20}
% Teintes claires pour fonds de tableaux / zones
\colorlet{popblueL}{popblue!10!white}
\colorlet{poporangeL}{poporange!14!white}
\colorlet{popgreenL}{popgreen!12!white}
\colorlet{poppurpleL}{poppurple!12!white}
\colorlet{poppinkL}{poppink!12!white}
\colorlet{popgoldL}{popgold!14!white}

\pagecolor{popcream}

% ============================================================
% Typographie
% ============================================================
\defaultfontfeatures{Ligatures=TeX}
\setmainfont{PlusJakartaSans-Medium.ttf}[Path=../../../fonts/,BoldFont=PlusJakartaSans-Bold.ttf]
% Symboles, API et emojis absents de la police principale : repli sur DejaVu Sans (caractères actifs)
\newfontface\symfont{DejaVuSans.ttf}[Path=../../../fonts/]
\makeatletter
\newcommand\symactive[1]{\catcode#1=13 \begingroup\lccode`\~=#1 \lowercase{\endgroup\def~}{{\symfont\char#1}}}
\newcommand\symblank[1]{\catcode#1=13 \begingroup\lccode`\~=#1 \lowercase{\endgroup\def~}{}}
\newcommand\symhyph[1]{\catcode#1=13 \begingroup\lccode`\~=#1 \lowercase{\endgroup\def~}{\mbox{-}}}
\newcount\symc
\newcommand\symrange[2]{\symc=#1 \@whilenum\symc<#2 \do{\expandafter\symactive\expandafter{\the\symc}\advance\symc by 1 }}
\newcommand\symblankrange[2]{\symc=#1 \@whilenum\symc<#2 \do{\expandafter\symblank\expandafter{\the\symc}\advance\symc by 1 }}
\makeatother
\symrange{"0250}{"02FF}\symactive{"2713}\symactive{"2714}\symactive{"2717}\symactive{"2718}\symactive{"2610}\symactive{"2611}\symactive{"2612}
\symhyph{"2011}\symhyph{"2E1A}\symblankrange{"1F300}{"1FAFF}\symblank{"FE0F}
% Maths en sans-serif (Fira Math) avec chiffres et lettres droites de Plus
% Jakarta, pour que les formules se fondent dans le texte.
\usepackage[math-style=ISO,bold-style=ISO]{unicode-math}
\setmathfont{FiraMath-Regular.otf}[Path=../../../fonts/]
\setmathfont{PlusJakartaSans-Medium.ttf}[Path=../../../fonts/,range={up/num,up/{Latin,latin},\mathup}]
\usepackage{icomma} % virgule décimale sans espace en mode math
% Caractères mathématiques Unicode absents des polices de texte (Plus Jakarta,
% Archivo Black) : rendus par leur équivalent mathématique, en texte comme en
% formule — sinon ils s'affichent en carré vide (ex. « Arithmétique dans ℤ »).
\usepackage{newunicodechar}
\newunicodechar{ℝ}{\ensuremath{\mathbb{R}}}
\newunicodechar{ℤ}{\ensuremath{\mathbb{Z}}}
\newunicodechar{ℂ}{\ensuremath{\mathbb{C}}}
\newunicodechar{ℕ}{\ensuremath{\mathbb{N}}}
\newunicodechar{ℚ}{\ensuremath{\mathbb{Q}}}
\newunicodechar{𝔻}{\ensuremath{\mathbb{D}}}
\newunicodechar{α}{\ensuremath{\alpha}}
\newunicodechar{β}{\ensuremath{\beta}}
\newunicodechar{γ}{\ensuremath{\gamma}}
\newunicodechar{δ}{\ensuremath{\delta}}
\newunicodechar{ε}{\ensuremath{\varepsilon}}
\newunicodechar{θ}{\ensuremath{\theta}}
\newunicodechar{λ}{\ensuremath{\lambda}}
\newunicodechar{μ}{\ensuremath{\mu}}
\newunicodechar{σ}{\ensuremath{\sigma}}
\newunicodechar{τ}{\ensuremath{\tau}}
\newunicodechar{φ}{\ensuremath{\varphi}}
\newunicodechar{ψ}{\ensuremath{\psi}}
\newunicodechar{ω}{\ensuremath{\omega}}
\newunicodechar{Σ}{\ensuremath{\Sigma}}
% majuscules produites par \MakeUppercase dans les titres (α-aminés -> Α)
\newunicodechar{Α}{\ensuremath{\alpha}}
\newunicodechar{Β}{\ensuremath{\beta}}
\newunicodechar{Γ}{\ensuremath{\Gamma}}
\newunicodechar{Δ}{\ensuremath{\Delta}}
\newunicodechar{Ε}{\ensuremath{\varepsilon}}
\newunicodechar{Θ}{\ensuremath{\Theta}}
\newunicodechar{Λ}{\ensuremath{\Lambda}}
\newunicodechar{Μ}{\ensuremath{\mu}}
\newunicodechar{Τ}{\ensuremath{\tau}}
\newunicodechar{Φ}{\ensuremath{\Phi}}
\newunicodechar{Ψ}{\ensuremath{\Psi}}
\newunicodechar{Ω}{\ensuremath{\Omega}}
\newunicodechar{↦}{\ensuremath{\mapsto}}
\newunicodechar{∈}{\ensuremath{\in}}
\newunicodechar{∉}{\ensuremath{\notin}}
\newunicodechar{∧}{\ensuremath{\wedge}}
\newunicodechar{⇔}{\ensuremath{\Leftrightarrow}}
\newunicodechar{⟺}{\ensuremath{\Longleftrightarrow}}
\newunicodechar{⇒}{\ensuremath{\Rightarrow}}
\newunicodechar{⟹}{\ensuremath{\Longrightarrow}}
\newunicodechar{∘}{\ensuremath{\circ}}
\newunicodechar{∪}{\ensuremath{\cup}}
\newunicodechar{∩}{\ensuremath{\cap}}
\newunicodechar{≡}{\ensuremath{\equiv}}
\newunicodechar{⊂}{\ensuremath{\subset}}
\newunicodechar{⊥}{\ensuremath{\perp}}
\newunicodechar{∃}{\ensuremath{\exists}}
\newunicodechar{∀}{\ensuremath{\forall}}
\newunicodechar{∛}{\ensuremath{\sqrt[3]{\,}}}
\newunicodechar{∜}{\ensuremath{\sqrt[4]{\,}}}
\newunicodechar{⃗}{\ensuremath{{}^{\to}}}
\newunicodechar{ⁿ}{\ifmmode^{n}\else\textsuperscript{\ensuremath{n}}\fi}
\newunicodechar{ˣ}{\ifmmode^{x}\else\textsuperscript{\ensuremath{x}}\fi}
\newunicodechar{ᵇ}{\ifmmode^{b}\else\textsuperscript{\ensuremath{b}}\fi}
\newunicodechar{ʳ}{\ifmmode^{r}\else\textsuperscript{\ensuremath{r}}\fi}
\newunicodechar{ᵘ}{\ifmmode^{u}\else\textsuperscript{\ensuremath{u}}\fi}
\newunicodechar{ʸ}{\ifmmode^{y}\else\textsuperscript{\ensuremath{y}}\fi}
\newunicodechar{ᵃ}{\ifmmode^{a}\else\textsuperscript{\ensuremath{a}}\fi}
\newunicodechar{ᵉ}{\ifmmode^{e}\else\textsuperscript{\ensuremath{e}}\fi}
\newunicodechar{ᵏ}{\ifmmode^{k}\else\textsuperscript{\ensuremath{k}}\fi}
\newunicodechar{ᵖ}{\ifmmode^{p}\else\textsuperscript{\ensuremath{p}}\fi}
\newunicodechar{ᴺ}{\ifmmode^{N}\else\textsuperscript{\ensuremath{N}}\fi}
\newunicodechar{ᶜ}{\ifmmode^{c}\else\textsuperscript{\ensuremath{c}}\fi}
\newunicodechar{ʰ}{\ifmmode^{h}\else\textsuperscript{\ensuremath{h}}\fi}
\newunicodechar{ᵠ}{\ifmmode^{q}\else\textsuperscript{\ensuremath{q}}\fi}
\newunicodechar{ₙ}{\ifmmode_{n}\else\textsubscript{\ensuremath{n}}\fi}
\newunicodechar{ₐ}{\ifmmode_{a}\else\textsubscript{\ensuremath{a}}\fi}
\newunicodechar{ᵢ}{\ifmmode_{i}\else\textsubscript{\ensuremath{i}}\fi}
\newunicodechar{ᵣ}{\ifmmode_{r}\else\textsubscript{\ensuremath{r}}\fi}
\newunicodechar{ₖ}{\ifmmode_{k}\else\textsubscript{\ensuremath{k}}\fi}
\newunicodechar{ᵦ}{\ifmmode_{\beta}\else\textsubscript{\ensuremath{\beta}}\fi}
\newunicodechar{ₓ}{\ifmmode_{x}\else\textsubscript{\ensuremath{x}}\fi}
\newfontfamily\popdisplay{ArchivoBlack-Regular.ttf}[Path=../../../fonts/]
\newfontfamily\poplogo{SpaceGrotesk-Bold.ttf}[Path=../../../fonts/]
\newfontfamily\popsemi{PlusJakartaSans-SemiBold.ttf}[Path=../../../fonts/]
\newfontfamily\popxbold{PlusJakartaSans-ExtraBold.ttf}[Path=../../../fonts/]
\newfontfamily\popxboldls{PlusJakartaSans-ExtraBold.ttf}[Path=../../../fonts/,LetterSpace=3.0]

\setlist[enumerate]{leftmargin=1.7em,itemsep=5pt,topsep=4pt}
\setlist[itemize]{leftmargin=1.7em,itemsep=4pt,topsep=4pt,label={\color{poporange}\textbullet}}
\setlength{\parindent}{0pt}
\setlength{\parskip}{5pt}
\renewcommand{\arraystretch}{1.25}

% pgfplots : style Pop par défaut
\pgfplotsset{
  every axis/.append style={axis line style={popink,line width=1pt},tick style={popink},
    grid style={popline,line width=0.5pt},label style={font=\small},tick label style={font=\footnotesize}},
  popcurve/.style={poporange,line width=1.8pt,no markers,smooth},
}
% Même style disponible dans un \draw TikZ simple (hors pgfplots)
\tikzset{popcurve/.style={poporange,line width=1.8pt}}
\tikzset{popgrid/.style={popline,very thin}}
\colorlet{popcurve}{poporange} % le nom du style est parfois utilisé comme couleur

% ============================================================
% Pill / squiggle
% ============================================================
\newcommand{\poppill}[3]{%
  \tikz[baseline=-0.55ex]\node[draw=popink,line width=1.2pt,rounded corners=2.6pt,%
    fill=#2,inner xsep=6.5pt,inner ysep=3pt]{%
    \popxboldls\fontsize{7}{7}\selectfont\color{#3}\MakeUppercase{#1}};%
}
\newcommand{\popsquiggle}[1]{%
  \tikz[baseline=-0.4ex]{
    \draw[#1,line width=2.6pt,line cap=round]
      plot [smooth] coordinates {(0,0.09) (0.34,0.01) (0.68,0.21) (1.02,0.01) (1.36,0.10)};
  }%
}
% Mot-clé surligné (marqueur orange sous le texte)
\newcommand{\cle}[1]{{\popxbold\color{popdark}#1}}

% ============================================================
% En-tête / pied
% ============================================================
\pagestyle{fancy}
\fancyhf{}
\renewcommand{\headrulewidth}{0pt}
\renewcommand{\footrulewidth}{0pt}
\lhead{\raisebox{-0.15ex}{\poplogo\fontsize{13}{13}\selectfont\color{popink}OnBuch\color{poporange}+}}
\rhead{\poppill{\DOCMATIERE\ \textbullet\ \DOCNIVEAU\ \textbullet\ Leçon \DOCLECON}{white}{popink}}
\lfoot{\popxbold\fontsize{7.6}{7.6}\selectfont\color{popmuted}\MakeUppercase{Cours signé OnBuch+}\ \textbullet\ {\popsemi\DOCTITRE}}
\rfoot{\tikz[baseline=-0.6ex]\node[rounded corners=8.5pt,fill=popink,minimum height=17pt,minimum width=17pt,inner xsep=5pt,inner sep=0]{\popxbold\fontsize{8}{8}\selectfont\color{popcream}\thepage\,/\,\pageref*{LastPage}};}

% ============================================================
% Titres
% ============================================================
\newcommand{\chaptitre}[2]{%
  \begin{tcolorbox}[enhanced,colback=poporange,colframe=popink,boxrule=2.6pt,arc=11pt,
      drop shadow=popink,left=15pt,right=15pt,top=12pt,bottom=11pt,before skip=3pt,after skip=18pt]
    \poppill{#2}{popink}{popcream}\par\vspace{9pt}
    {\raggedright\hyphenpenalty=10000\exhyphenpenalty=10000\popdisplay\fontsize{22}{24}\selectfont\color{popcream}\MakeUppercase{#1}\par}
  \end{tcolorbox}
}
% Section numérotée : pastille orange avec le numéro + titre Archivo
\newcounter{popsec}
\newcommand{\coursec}[1]{%
  \stepcounter{popsec}\setcounter{popsub}{0}%
  \par\vspace{16pt}\noindent
  \begin{minipage}{\linewidth}
  \tikz[baseline=-0.7ex]\node[draw=popink,line width=1.6pt,fill=poporange,rounded corners=4pt,minimum size=22pt,inner sep=0]{\popdisplay\fontsize{12}{12}\selectfont\color{popcream}\thepopsec};%
  \hspace{8pt}{\popdisplay\fontsize{15}{17}\selectfont\color{popink}\MakeUppercase{#1}}\par
  \vspace{4pt}\noindent{\color{poporange}\rule{36pt}{3.6pt}}
  \end{minipage}\par\nopagebreak\vspace{8pt}
}
\newcounter{popsub}
\newcommand{\courssub}[1]{%
  \stepcounter{popsub}%
  \par\vspace{9pt}\noindent{\popxbold\fontsize{11.5}{13}\selectfont\color{popink}{\color{poporange}\thepopsec.\thepopsub}\hspace{6pt}#1}\par\nopagebreak\vspace{4pt}
}

% ============================================================
% Boîtes pédagogiques — même recette : fond blanc, bordure épaisse colorée,
% ombre dure encre, badge plein. Titre optionnel [#1].
% ============================================================
\tcbset{popbase/.style={breakable,enhanced,colback=white,boxrule=2.3pt,arc=9pt,
  shadow={3pt}{-3pt}{0pt}{popink},left=14pt,right=14pt,top=11pt,bottom=11pt,before skip=11pt,after skip=15pt,
  fontupper=\popsemi\fontsize{10.6}{14.5}\selectfont\color{popink},
  fontlower=\popsemi\fontsize{10.6}{14.5}\selectfont\color{popink}}}
% Coche dessinée (le glyphe ✓ n'existe pas dans les polices du document)
\newcommand{\popcheck}[1]{\tikz[baseline=-0.5ex]\draw[#1,line width=1.6pt,line cap=round,line join=round](-0.13,0.0)--(-0.04,-0.09)--(0.13,0.1);}
\providecommand{\coche}{\popcheck{popgreen}}
\providecommand{\resultat}[1]{\par\smallskip\noindent\fcolorbox{popgreen}{popgreenL}{\parbox{\dimexpr\linewidth-2\fboxsep-2\fboxrule\relax}{\textbf{Résultat.} #1}}\par\smallskip}
\newcommand{\popboxhead}[3]{\poppill{#1}{#2}{white}\ifx\relax#3\relax\else\hspace{6pt}{\popxbold\fontsize{10.6}{12}\selectfont\color{popink}#3}\fi\par\vspace{7pt}}

\newtcolorbox{aretenir}[1][]{popbase,colframe=popgreen,before upper={\popboxhead{À retenir}{popgreen}{#1}}}
% Texte en anglais (document de lecture, dialogue, extrait) : cadre
% d'étude. Un titre optionnel (source, auteur) s'affiche dans l'en-tête.
\newtcolorbox{textebox}[1][]{popbase,colframe=popblue,before upper={\popboxhead{Text}{popblue}{#1}\begin{otherlanguage}{english}},after upper={\end{otherlanguage}}}
% Marques ✓ / ✗ (forme correcte / fautive) souvent utilisées par les modèles
\providecommand{\cross}{\textcolor{poppink}{\ensuremath{\times}}}
\providecommand{\xmark}{\cross}\providecommand{\crossmark}{\cross}\providecommand{\cmark}{\ensuremath{\checkmark}}
% Ligne de réponse / justification à compléter (inventée par les modèles)
\providecommand{\justification}[1]{\par\noindent\textit{Justification~:} \dotfill\par}
% \textipa (package tipa non chargé) : la police ne couvre pas l'API -> simple passage
\providecommand{\textipa}[1]{#1}
% Soulignement (ulem non chargé) : \uline, \ul
\providecommand{\uline}[1]{\underline{#1}}\providecommand{\ul}[1]{\underline{#1}}
% \glossaire{\item ...} inventée par les modèles : liste de glossaire
\providecommand{\glossaire}[1]{\begin{itemize}#1\end{itemize}}
% \alert (beamer) inventée par les modèles : texte mis en évidence
\providecommand{\alert}[1]{\textbf{\textcolor{poppink}{#1}}}
% variantes ASCII d'ordinaux écrites en macro
\providecommand{\iere}{\textsuperscript{re}}\providecommand{\ieme}{\textsuperscript{e}}\providecommand{\ere}{\textsuperscript{re}}\providecommand{\eme}{\textsuperscript{e}}\providecommand{\er}{\textsuperscript{er}}\providecommand{\ie}{\textsuperscript{e}}
% Ordinaux français écrits en macro par les modèles : 1\ière, 3\ième
\newcommand{\ière}{\textsuperscript{re}}\newcommand{\ième}{\textsuperscript{e}}
% \exemple (inventée par les modèles) : étiquette « Exemple : »
\providecommand{\exemple}{\textbf{Exemple~:}\ }
% \square utilisable aussi hors mode math (cases à cocher)
\AtBeginDocument{\let\squareMath\square\renewcommand{\square}{\ensuremath{\squareMath}}}
% Mot ou phrase en anglais à l'intérieur d'un paragraphe français
\newcommand{\eng}[1]{\ifmmode\text{\textit{#1}}\else\begingroup\selectlanguage{english}\itshape #1\endgroup\fi}
\let\esp\eng % alias toléré
% flèches d'intonation inventées par les modèles
\providecommand{\textsearrow}{}\renewcommand{\textsearrow}{\ensuremath{\searrow}}\providecommand{\textnearrow}{}\renewcommand{\textnearrow}{\ensuremath{\nearrow}}
% \fr{..} : mot français (inventé par les modèles)
\providecommand{\fr}[1]{\textit{#1}}
\newtcolorbox{exemplebox}[1][]{popbase,colframe=poporange,before upper={\popboxhead{Exemple}{poporange}{#1}}}
\newtcolorbox{definition}[1][]{popbase,colframe=popblue,before upper={\popboxhead{Définition}{popblue}{#1}}}
\newtcolorbox{propriete}[1][]{popbase,colframe=poppurple,before upper={\popboxhead{Propriété}{poppurple}{#1}}}
\newtcolorbox{methode}[1][]{popbase,colframe=popgold,before upper={\popboxhead{Méthode}{popgold}{#1}}}
\newtcolorbox{attention}[1][]{popbase,colframe=poppink,before upper={\popboxhead{Attention}{poppink}{#1}}}
\newtcolorbox{experience}[1][]{popbase,colframe=popblue,colback=popblueL,before upper={\popboxhead{Expérience}{popblue}{#1}}}
% « Le savais-tu ? » : carte violette pleine (contexte, culture, Cameroun)
\newtcolorbox{savaistu}[1][]{popbase,colback=poppurple,colframe=popink,
  fontupper=\popsemi\fontsize{10.4}{14.2}\selectfont\color{white},
  before upper={\poppill{Le savais-tu\,?}{white}{poppurple}\ifx\relax#1\relax\else\hspace{6pt}{\popxbold\color{white}#1}\fi\par\vspace{7pt}}}
% Exercice résolu : énoncé en haut, \tcblower puis la solution sur fond crème
\newcounter{popexo}\newcounter{popexr}
\newtcolorbox{exoresolu}[1][]{popbase,colframe=poporange,colbacklower=popcream,
  segmentation style={popink,line width=1.2pt,dashed},
  before upper={\stepcounter{popexr}\popboxhead{Exercice résolu \thepopexr}{poporange}{#1}},
  before lower={\poppill{Solution}{popink}{popcream}\par\vspace{6pt}}}
% Exercice (non résolu en place) — la correction est dans la partie Corrigés
\newtcolorbox{exercice}[1][]{popbase,colframe=popink,
  before upper={\stepcounter{popexo}\popboxhead{Exercice \thepopexo}{popink}{#1}}}
% Corrigé
\newtcolorbox{corrige}[1][]{popbase,colframe=popgreen,colback=popgreenL,
  before upper={\popboxhead{Corrigé}{popgreen}{#1}}}
% Objectifs / prérequis / bilan
\newtcolorbox{objectifs}{popbase,colframe=popink,colback=poporangeL,
  before upper={\popboxhead{Objectifs de la leçon}{popink}{}}}
\newtcolorbox{prerequis}{popbase,colframe=popink,colback=popblueL,
  before upper={\popboxhead{Prérequis}{popblue}{}}}
\newtcolorbox{bilan}{popbase,colframe=popink,colback=white,boxrule=2.8pt,
  before upper={\popboxhead{Fiche bilan}{poporange}{L'essentiel à connaître pour l'examen}}}
% Figure encadrée avec légende
\newtcolorbox{popfigure}[1][]{enhanced,breakable=false,colback=white,colframe=popline,boxrule=1.4pt,arc=8pt,
  left=10pt,right=10pt,top=10pt,bottom=8pt,before skip=10pt,after skip=12pt,halign=center,
  lower separated=false,halign lower=center,
  fontlower=\popsemi\fontsize{9}{11.5}\selectfont\color{popmuted},
  #1}
\newcommand{\legende}[1]{\par\vspace{6pt}{\popsemi\fontsize{9}{11.5}\selectfont\color{popmuted}#1}}

% ============================================================
% Signature OnBuch+
% ============================================================
\newcommand{\poplogoinline}[1]{{\poplogo\fontsize{#1}{#1}\selectfont\color{popink}OnBuch\color{poporange}+}}
\newcommand{\signatureonbuch}{%
  \begin{tcolorbox}[enhanced,colback=popink,colframe=popink,boxrule=2.2pt,arc=10pt,
      drop shadow=poporange,left=14pt,right=14pt,top=10pt,bottom=10pt,before skip=0pt,after skip=0pt]
    \tikz[baseline=-0.7ex]\node[circle,fill=popgreen,draw=popcream,line width=1.3pt,minimum size=22pt,inner sep=0]{\popcheck{white}};%
    \hspace{9pt}\begin{minipage}[c]{0.62\linewidth}
      {\popxbold\fontsize{12}{13}\selectfont\color{popcream}Signé\ \textbullet\ {\poplogo OnBuch\color{poporange}+}}\par\vspace{2pt}
      {\raggedright\popsemi\fontsize{8}{10.5}\selectfont\color{popline}\MakeUppercase{Équipe pédagogique OnBuch+}\\\MakeUppercase{Conforme au programme officiel MINESEC}\par}
    \end{minipage}\hfill
    {\poplogo\fontsize{22}{22}\selectfont\color{popcream}OnBuch\color{poporange}+}
  \end{tcolorbox}%
}

% ============================================================
% Page de garde
% #1 titre · #2 matière · #3 niveau · #4 module · #5 n° leçon · #6 durée ·
% #7 description / objectifs (texte ou itemize)
% ============================================================
\newcommand{\pagegarde}[7]{%
  \thispagestyle{empty}
  \newgeometry{a4paper,margin=1.7cm,top=1.6cm,bottom=1.4cm}%
  \begin{tikzpicture}[remember picture,overlay]
    \fill[poporange!18] ([xshift=-3.2cm,yshift=-3.6cm]current page.north east) circle (4.6cm);
    \fill[poppurple!14] ([xshift=2.4cm,yshift=7.6cm]current page.south west) circle (3.8cm);
  \end{tikzpicture}%
  {\poplogo\fontsize{30}{30}\selectfont\color{popink}OnBuch\color{poporange}+}\hfill
  \raisebox{0.6ex}{\poppill{Cours signé \textbullet\ Premium}{popink}{popcream}}\par
  \vspace{1pt}\popsquiggle{poppurple}\par
  \vspace{8pt}
  \begin{tcolorbox}[enhanced,colback=poporange,colframe=popink,boxrule=2.8pt,arc=13pt,
      drop shadow=popink,left=18pt,right=18pt,top=12pt,bottom=13pt,after skip=10pt]
    \poppill{#2 \textbullet\ #3}{popink}{popcream}\hspace{5pt}\poppill{#4}{white}{popink}\par\vspace{8pt}
    {\popdisplay\fontsize{14}{14}\selectfont\color{popink}LEÇON #5}\par\vspace{4pt}
    {\raggedright\hyphenpenalty=10000\exhyphenpenalty=10000\popdisplay\fontsize{25}{27}\selectfont\color{popcream}\MakeUppercase{#1}\par}
    \vspace{8pt}
    {\popxbold\fontsize{9.5}{9.5}\selectfont\color{popink}\MakeUppercase{Durée conseillée : #6}}
  \end{tcolorbox}
  \begin{tcolorbox}[enhanced,colback=white,colframe=popink,boxrule=2.2pt,arc=10pt,
      drop shadow=popink,left=16pt,right=16pt,top=10pt,bottom=10pt,after skip=8pt]
    \poppill{Dans cette leçon}{poppurple}{white}\par\vspace{5pt}
    \popsemi\fontsize{9.3}{12.6}\selectfont\color{popink}#7
  \end{tcolorbox}
  \noindent\begin{minipage}[t]{0.487\linewidth}
    \begin{tcolorbox}[enhanced,colback=popgreen,colframe=popink,boxrule=2.2pt,arc=10pt,
        drop shadow=popink,left=11pt,right=11pt,top=10pt,bottom=10pt,height=3.1cm,valign=center]
      \tcbox[colback=white,colframe=popink,boxrule=1.3pt,arc=5pt,boxsep=0pt,nobeforeafter,
        left=3pt,right=3pt,top=3pt,bottom=3pt,valign=center]{\qrcode[height=1.9cm]{https://play.google.com/store/apps/details?id=com.luvvix.onbuch&hl=fr}}%
      \hspace{8pt}\begin{minipage}[c]{0.5\linewidth}
        {\popdisplay\fontsize{11}{12.5}\selectfont\color{white}\MakeUppercase{Télécharge OnBuch}}\par\vspace{4pt}
        {\raggedright\popsemi\fontsize{8.4}{11}\selectfont\color{white}Gratuit \textbullet\ Google Play\\Tuteur IA, annales, concours, résultats\par}
      \end{minipage}
    \end{tcolorbox}
  \end{minipage}\hfill
  \begin{minipage}[t]{0.487\linewidth}
    \begin{tcolorbox}[enhanced,colback=poppurple,colframe=popink,boxrule=2.2pt,arc=10pt,
        drop shadow=popink,left=11pt,right=11pt,top=10pt,bottom=10pt,height=3.1cm,valign=center]
      \tikz[baseline=-0.7ex]\node[circle,fill=white,draw=popink,line width=1.3pt,minimum size=1.6cm,inner sep=0]{\popdisplay\fontsize{24}{24}\selectfont\color{poppurple}+};%
      \hspace{8pt}\begin{minipage}[c]{0.6\linewidth}
        {\popdisplay\fontsize{11}{12.5}\selectfont\color{white}\MakeUppercase{Passe à OnBuch+}}\par\vspace{4pt}
        {\raggedright\popsemi\fontsize{8.4}{11}\selectfont\color{white}Tous les cours signés, Tuteur IA illimité, fiches PDF — onglet OnBuch+ de l'app\par}
      \end{minipage}
    \end{tcolorbox}
  \end{minipage}
  \vspace{8pt}
  \popcontenu
  \vfill
  \signatureonbuch
  \restoregeometry
  \newpage
}

% Rangée « Ce cours contient » (page de garde)
\newcommand{\poptile}[3]{% #1 couleur #2 gros texte #3 légende
  \begin{tcolorbox}[enhanced,colback=#1,colframe=popink,boxrule=2pt,arc=9pt,shadow={2.5pt}{-2.5pt}{0pt}{popink},
      left=6pt,right=6pt,top=9pt,bottom=9pt,halign=center,valign=center,height=1.75cm,width=\linewidth,nobeforeafter]
    {\popdisplay\fontsize{12.5}{13}\selectfont\color{white}\MakeUppercase{#2}}\par\vspace{3pt}
    {\centering\popsemi\fontsize{7.8}{9.5}\selectfont\color{white}#3\par}
  \end{tcolorbox}}
\newcommand{\popcontenu}{%
  \par\noindent{\popxboldls\fontsize{7.5}{7.5}\selectfont\color{popmuted}CE COURS SIGNÉ CONTIENT}\par\vspace{7pt}
  \noindent\begin{minipage}[t]{0.235\linewidth}\poptile{popblue}{Cours}{complet, illustré, pas à pas}\end{minipage}\hfill
  \begin{minipage}[t]{0.235\linewidth}\poptile{poporange}{Méthodes}{et exercices résolus}\end{minipage}\hfill
  \begin{minipage}[t]{0.235\linewidth}\poptile{poppink}{Exercices}{type examen + corrigés}\end{minipage}\hfill
  \begin{minipage}[t]{0.235\linewidth}\poptile{popgreen}{Fiche bilan}{l'essentiel à retenir}\end{minipage}\par}

% Sommaire
\newtcolorbox{sommaire}{enhanced,colback=white,colframe=popink,boxrule=2.2pt,arc=10pt,
  drop shadow=popink,left=18pt,right=18pt,top=13pt,bottom=13pt,before skip=6pt,after skip=20pt,
  before upper={\poppill{Sommaire}{poppurple}{white}\par\vspace{9pt}\popsemi\fontsize{10.6}{14.5}\selectfont\color{popink}}}

% Signature de fin de cours — poussée tout en bas de la dernière page
\newcommand{\findecours}{%
  \par\vfill\nopagebreak
  \begin{center}\popsquiggle{poporange}\end{center}\vspace{4pt}
  \signatureonbuch
}
\AtBeginDocument{\def\llbracket{\mathopen{[\mkern-3.5mu[}}\def\rrbracket{\mathclose{]\mkern-3.5mu]}}} % intervalles d'entiers (glyphe ⟦ absent de la police)
% Glyphes absents de Fira Math : redéfinis à partir de glyphes disponibles
\AtBeginDocument{%
  \def\setminus{\mathbin{\backslash}}\let\smallsetminus\setminus
  \def\vdots{\mathord{\vbox{\baselineskip3.2pt\lineskiplimit0pt\kern4pt\hbox{.}\hbox{.}\hbox{.}}}}%
  \def\ddots{\mathinner{\mkern1mu\raise6.5pt\hbox{.}\mkern2mu\raise3.6pt\hbox{.}\mkern2mu\raise0.7pt\hbox{.}\mkern1mu}}%
  \def\triangle{\mathord{\scalebox{0.9}{$\symup{\Delta}$}}}%
  \def\nabla{\mathord{\raisebox{\depth}{\scalebox{1}[-1]{$\symup{\Delta}$}}}}%
  \def\checkmark{\popcheck{popdark}}%
  \def\star{\mathbin{\ast}}\def\bigstar{\mathord{\ast}}%
  \def\diamond{\mathbin{\scalebox{0.6}{$\Diamond$}}}%
  \def\bigcup{\mathop{\mathchoice{\raisebox{-0.3em}{\scalebox{1.7}{$\cup$}}}{\scalebox{1.25}{$\cup$}}{\cup}{\cup}}\displaylimits}%
  \def\bigcap{\mathop{\mathchoice{\raisebox{-0.3em}{\scalebox{1.7}{$\cap$}}}{\scalebox{1.25}{$\cap$}}{\cap}{\cap}}\displaylimits}%
  \def\lesseqqgtr{\mathrel{\lessgtr}}\def\bigoplus{\mathop{\oplus}\displaylimits}%
}
\providecommand{\ssi}{si et seulement si} % abréviation fréquente
\AtBeginDocument{\providecommand{\wideparen}{\overparen}} % arc de cercle

%% FIGFIX-BEGIN v5 — correctifs globaux des figures (docs/onbuch-generation/tools/figfix)
\usepackage{adjustbox}\usepackage{etoolbox}\tcbuselibrary{raster}
% toute figure TikZ/pgfplots plus large que la ligne est réduite à la largeur disponible
\BeforeBeginEnvironment{tikzpicture}{\begin{adjustbox}{max width=\linewidth}}
\AfterEndEnvironment{tikzpicture}{\end{adjustbox}}
% légendes pgfplots lisibles par-dessus les courbes
\pgfplotsset{every axis/.append style={legend style={fill=white,fill opacity=0.92,text opacity=1,draw=black!15,inner sep=3pt}}}
% étiquettes d'arêtes (syntaxe edge["..."]) avec fond blanc
\tikzset{every edge quotes/.append style={fill=white,inner sep=1.2pt}}
%% FIGFIX-END
\begin{document}

\pagegarde{Lesson 54 — Vocabulary: Words and expressions related to using ICTs to enhance work or learning. \# Semaine / Période Chapitre Titre Date effective Visa de l'enseignant}{Anglais}{Tle A}{Chapter 5 : Media and Communication}{EN54}{2 h}{This vocabulary lesson equips Terminale A students with the essential words, collocations and expressions used to talk about Information and Communication Technologies (ICTs) in work and learning contexts. Through thematic lexical fields, word families, false friends and contextual practice, learners will be able to speak and write accurately about digital tools and their uses.
\vspace{4pt}
\begin{multicols}{2}\small
\begin{itemize}
\item À la fin de cette leçon, je sais nommer les principaux outils et services numériques (hardware, software, applications, plateformes) en anglais
\item À la fin de cette leçon, je sais employer les verbes et collocations typiques du domaine des TIC (browse the web, download a file, log in, back up data...)
\item À la fin de cette leçon, je sais utiliser le vocabulaire de l'apprentissage en ligne (e-learning, distance learning, online course, tutorial...)
\item À la fin de cette leçon, je sais décrire comment les TIC améliorent le travail (teleworking, video conferencing, cloud storage, productivity...)
\item À la fin de cette leçon, je sais éviter les faux amis et les calques du français dans ce champ lexical (to assist ≠ assister à, digital ≠ numérique dans tous les contextes...)
\item À la fin de cette leçon, je sais former des mots de la même famille à partir d'une racine (compute → computer → computing → computerize)
\end{itemize}\end{multicols}}

\begin{sommaire}
\begin{multicols}{2}
\begin{enumerate}
\item Discovering the theme: ICTs around us
\item Lexical field 1: Devices, tools and the Internet
\item Lexical field 2: ICTs for learning
\item Lexical field 3: ICTs for work
\item Collocations, word families and false friends
\item Using the vocabulary in context
\item Activité d'intégration
\item Méthodes et exercices résolus
\item Exercices
\item Corrigés des exercices
\item Fiche bilan
\end{enumerate}
\end{multicols}
\end{sommaire}

\begin{objectifs}
\begin{itemize}
\item À la fin de cette leçon, je sais nommer les principaux outils et services numériques (hardware, software, applications, plateformes) en anglais
\item À la fin de cette leçon, je sais employer les verbes et collocations typiques du domaine des TIC (browse the web, download a file, log in, back up data...)
\item À la fin de cette leçon, je sais utiliser le vocabulaire de l'apprentissage en ligne (e-learning, distance learning, online course, tutorial...)
\item À la fin de cette leçon, je sais décrire comment les TIC améliorent le travail (teleworking, video conferencing, cloud storage, productivity...)
\item À la fin de cette leçon, je sais éviter les faux amis et les calques du français dans ce champ lexical (to assist ≠ assister à, digital ≠ numérique dans tous les contextes...)
\item À la fin de cette leçon, je sais former des mots de la même famille à partir d'une racine (compute → computer → computing → computerize)
\item À la fin de cette leçon, je sais prononcer correctement le vocabulaire technique courant (data, device, website, software...)
\item À la fin de cette leçon, je sais rédiger un court paragraphe expliquant comment j'utilise les TIC pour apprendre ou travailler
\end{itemize}
\end{objectifs}

\begin{prerequis}
\begin{itemize}
\item Vocabulaire général des médias et de la communication vu en début de chapitre (media, press, social networks)
\item Structures de base : present simple et present continuous pour décrire des habitudes et des actions
\item Notion de champ lexical et de collocation vue en Première
\item Verbes courants de la vie quotidienne (use, work, learn, send, receive)
\item Expression de l'opinion simple (I think, in my opinion, it is useful because...)
\end{itemize}
\end{prerequis}

\newpage

\coursec{Discovering the theme: ICTs around us}

\courssub{Warm-up: a world connected}

Imagine this: during the holidays, your cousin in Bamenda prepares her GCE revision using online past papers, a classmate in Douala follows free coding tutorials on YouTube, and your uncle in Buea runs his whole business from his smartphone with WhatsApp and mobile money. ICTs are no longer a luxury in Cameroon — they are everyday tools for learning and working. But do you have the right English words to talk about them? Today, we build the vocabulary you need to describe how technology enhances work and learning.

\medskip
\textbf{Problématique de la leçon :} How can we name, classify and use the English words that describe the digital tools which improve our studies and our jobs?

\courssub{Brainstorming: which digital tools do you use every day?}

Avant toute définition, partons de votre expérience. Répondez mentalement à ces trois questions, en anglais :

\begin{itemize}
\item Which devices do you use every day? (\eng{a smartphone, a laptop, a tablet, a desktop computer...})
\item Which applications or websites do you open most often? (\eng{WhatsApp, Google Classroom, YouTube, Facebook, Zoom...})
\item What do you use them for? (\eng{to chat, to revise, to download documents, to send money, to search for information...})
\end{itemize}

\begin{experience}[Classroom brainstorming]
En groupes de quatre, listez en deux minutes tous les outils numériques que vous utilisez chaque semaine, en anglais. Classez-les ensuite en trois colonnes : \emph{things I can touch} (que je peux toucher), \emph{programs I use} (programmes que j'utilise), \emph{connections I need} (connexions dont j'ai besoin). Comparez ensuite votre liste avec celle d'un autre groupe. Gardez votre tableau : il servira à vérifier la distinction \emph{hardware / software / services} à la fin de la section.
\end{experience}

\courssub{Reading: How ICTs Changed My School Year}

Lisons maintenant le récit de Nfor, un élève de Terminale qui raconte comment les outils numériques ont transformé sa préparation du baccalauréat.

\begin{textebox}[How ICTs Changed My School Year]
Last term, I could not attend extra classes in town because the transport fares were too high. I felt worried about my GCE preparation. Then my elder brother showed me a website where I could download past GCE papers for free. I saved them on my phone and printed a few at the cybercafe near our quarter.

I also started watching video tutorials on my phone every evening. Some teachers explain mathematics better on screen than in a crowded classroom! Our English teacher created a WhatsApp group where she uploads notes and audio explanations after every lesson. When I do not understand something, I simply replay the recording.

I even joined an online study group with students from Buea and Yaoundé. We share summaries, quiz each other and correct our essays together. At first, my father thought I was wasting time on my phone. But when my mock exam results came out, he changed his mind.

Thanks to these tools, I revised more efficiently and improved my grades in three subjects. ICTs did not replace my teachers or my textbooks — they helped me use them better. Now I know that a smartphone is not only for games and music: it can be a real classroom in your pocket.
\end{textebox}

\textbf{Glossaire (mots difficiles du texte) :}

\begin{center}
\begin{tabularx}{\linewidth}{|l|l|X|}
\hline
\rowcolor{popblueL} English word & Nature & Français \\
\hline
fare & nom & prix du transport, tarif \\
\hline
quarter & nom & quartier (anglais camerounais) \\
\hline
cybercafe & nom & cybercafé \\
\hline
to download & verbe & télécharger (vers son appareil) \\
\hline
to upload & verbe & téléverser, mettre en ligne \\
\hline
tutorial & nom & tutoriel, leçon guidée \\
\hline
recording & nom & enregistrement \\
\hline
mock exam & nom & examen blanc \\
\hline
grade & nom & note, résultat scolaire \\
\hline
efficiently & adverbe & efficacement \\
\hline
\end{tabularx}
\end{center}

\textbf{Comprehension questions :}
\begin{enumerate}
\item Why could Nfor not attend extra classes in town?
\item Name three digital tools or services mentioned in the text.
\item What does the teacher upload on the WhatsApp group?
\item How did Nfor's father react at first, and why did he change his mind?
\item In your opinion, what is the main message of the last paragraph?
\end{enumerate}

\courssub{Defining ICTs}

À partir du texte, observons : Nfor utilise un \emph{phone} (un objet), un \emph{website} et des \emph{tutorials} (des contenus et programmes), et un \emph{WhatsApp group} (un service en réseau). Tous ces éléments appartiennent à une même famille : les \cle{ICTs}.

\begin{definition}[What are ICTs?]
\cle{ICTs} is the abbreviation of \emph{Information and Communication Technologies}. It refers to all the digital tools — devices, software and networks — used to \emph{create, store, share and exchange information}. In French: les \emph{technologies de l'information et de la communication} (TIC).
\par\medskip
\eng{ICTs help students learn and workers communicate faster.} « Les TIC aident les élèves à apprendre et les travailleurs à communiquer plus vite. »
\end{definition}

\begin{attention}[ICT ou ICTs ?]
En anglais, on dit souvent \eng{ICT} (sans \emph{s}) comme nom collectif ou adjectif : \eng{an ICT lesson}, \eng{ICT skills}. Avec un \emph{s}, \eng{ICTs} insiste sur la pluralité des technologies : \eng{Modern ICTs include smartphones and social media.} Les deux formes sont correctes ; ne traduisez jamais par « ICT's » avec une apostrophe, qui marquerait le génitif.
\end{attention}

\courssub{Guessing meaning from context}

Dans le texte, vous avez rencontré des mots nouveaux : \eng{uploads}, \eng{tutorials}, \eng{mock exam}. Bonne nouvelle : un bon lecteur n'a pas besoin du dictionnaire pour chaque mot. Voici la démarche professionnelle.

\begin{methode}[How to guess the meaning of an unknown word]
\begin{enumerate}
\item \textbf{Read the whole sentence} — lisez la phrase entière, pas seulement le mot inconnu.
\item \textbf{Identify the word class} — identifiez la nature du mot : nom, verbe, adjectif, adverbe ? (Sa position et ses terminaisons aident : \emph{-tion}, \emph{-ment} = nom ; \emph{-ly} = adverbe ; \emph{-ed} = verbe au passé.)
\item \textbf{Look for context clues} — cherchez les indices du contexte : exemples, synonymes, oppositions, conséquences logiques.
\item \textbf{Make a hypothesis} — formulez une hypothèse de sens en français ou en anglais simple.
\item \textbf{Check in the glossary} — vérifiez enfin dans le glossaire ou le dictionnaire.
\end{enumerate}
\end{methode}

\begin{exemplebox}[Applying the method: “uploads”]
Phrase du texte : \eng{Our English teacher created a WhatsApp group where she uploads notes and audio explanations.}
\begin{itemize}
\item Nature : \eng{she uploads} → verbe conjugué (3\textsuperscript{e} personne du singulier).
\item Indices : un groupe WhatsApp, des notes et des audios → l'action de \emph{mettre} des fichiers à disposition des élèves.
\item Hypothèse : « mettre en ligne, envoyer sur le groupe ».
\item Vérification : glossaire → \emph{to upload} = téléverser, mettre en ligne. Hypothèse confirmée.
\end{itemize}
\end{exemplebox}

\courssub{Hardware, software or services?}

Trois grandes familles organisent tout le vocabulaire des TIC. Cette distinction est la clé de toute la leçon.

\begin{definition}[Hardware, software, services]
\begin{itemize}
\item \cle{Hardware} (mot invariable, sans pluriel) : the physical equipment you can touch — le \emph{matériel} : \eng{a computer, a smartphone, a printer, a keyboard, a router}.
\item \cle{Software} (invariable aussi) : the programs and applications that run on the hardware — les \emph{logiciels} : \eng{a word processor, a browser, an app, an operating system}.
\item \cle{Services} : the connections and online platforms that link users — les \emph{services} : \eng{the Internet, a mobile network, e-mail, social media, cloud storage}.
\end{itemize}
\end{definition}

\begin{attention}[Deux mots invariables à ne jamais pluraliser]
\eng{Hardware} et \eng{software} n'ont \textbf{jamais} de pluriel en anglais. On dit \eng{This software is expensive} (jamais \emph{softwares are}). Pour compter, utilisez \eng{a piece of software} ou \eng{an application}. De même : \eng{information} est invariable — jamais \emph{informations} comme en français !
\end{attention}

\begin{popfigure}
\begin{center}
\begin{center}\tcbox[colback=popblue,colframe=popblue,coltext=white,arc=9pt,boxsep=4pt,left=6pt,right=6pt]{\bfseries\small ICTs}\end{center}
\vspace{-2pt}
\begin{tcbraster}[raster columns=2,raster equal height=rows,raster column skip=6pt,raster row skip=6pt,enhanced,blankest]
\begin{tcolorbox}[enhanced,colback=white,colframe=poporange,boxrule=0.9pt,arc=3pt,title={Hardware \emph{matériel}},fonttitle=\bfseries\scriptsize,coltitle=white,colbacktitle=poporange,left=4pt,right=4pt,top=2pt,bottom=2pt,boxsep=1pt,toptitle=1.5pt,bottomtitle=1.5pt,halign=left]
\begin{itemize}[nosep,leftmargin=*,itemsep=1pt,font=\scriptsize]
\item smartphone
\item laptop
\item printer
\end{itemize}
\end{tcolorbox}
\begin{tcolorbox}[enhanced,colback=white,colframe=popgreen,boxrule=0.9pt,arc=3pt,title={Software \emph{logiciels}},fonttitle=\bfseries\scriptsize,coltitle=white,colbacktitle=popgreen,left=4pt,right=4pt,top=2pt,bottom=2pt,boxsep=1pt,toptitle=1.5pt,bottomtitle=1.5pt,halign=left]
\begin{itemize}[nosep,leftmargin=*,itemsep=1pt,font=\scriptsize]
\item apps
\item browser
\item word processor
\end{itemize}
\end{tcolorbox}
\begin{tcolorbox}[enhanced,colback=white,colframe=poppurple,boxrule=0.9pt,arc=3pt,title={Services \emph{services}},fonttitle=\bfseries\scriptsize,coltitle=white,colbacktitle=poppurple,left=4pt,right=4pt,top=2pt,bottom=2pt,boxsep=1pt,toptitle=1.5pt,bottomtitle=1.5pt,halign=left]
\begin{itemize}[nosep,leftmargin=*,itemsep=1pt,font=\scriptsize]
\item Internet
\item e-mail
\item cloud storage
\end{itemize}
\end{tcolorbox}
\end{tcbraster}

\end{center}
\legende{Carte mentale : les trois familles du vocabulaire des TIC}
\end{popfigure}

\begin{popfigure}
\begin{center}
\begin{tikzpicture}[
  phone/.style={draw=popdark, fill=popblueL, rounded corners=4pt, minimum width=1.6cm, minimum height=2.8cm, align=center, font=\bfseries},
  icon/.style={draw=popdark, rounded corners=3pt, minimum width=2.2cm, minimum height=1.1cm, align=center, font=\small},
  link/.style={-{Stealth[length=2.5mm]}, thick, popmuted}]
\node[phone, align=center] (p) at (0,0){smart-\\phone};
\node[icon, fill=popgreenL, align=center] (c) at (4.2,2.4){cloud\\ \emph{storage}};
\node[icon, fill=poporangeL, align=center] (b) at (4.2,0.8){e-book\\ \emph{reading}};
\node[icon, fill=poppinkL, align=center] (w) at (4.2,-0.8){briefcase\\ \emph{work}};
\node[icon, fill=popgoldL, align=center] (g) at (4.2,-2.4){graduation cap\\ \emph{studies}};
\draw[link] (p.east) -- (c.west);
\draw[link] (p.east) -- (b.west);
\draw[link] (p.east) -- (w.west);
\draw[link] (p.east) -- (g.west);
\end{tikzpicture}
\end{center}
\legende{Un seul appareil, quatre usages : stocker, lire, travailler, étudier}
\end{popfigure}

\courssub{First vocabulary table: English and French}

\begin{center}
\begin{tabularx}{\linewidth}{|l|X|X|}
\hline
\rowcolor{popblueL} English & Français & Exemple en contexte \\
\hline
a device & un appareil & \eng{A smartphone is a useful device.} \\
\hline
hardware & le matériel & \eng{Our school bought new hardware.} \\
\hline
software & le logiciel & \eng{This software corrects your spelling.} \\
\hline
an app / an application & une application & \eng{She uses a dictionary app.} \\
\hline
to download & télécharger (recevoir) & \eng{I downloaded past papers.} \\
\hline
to upload & téléverser (envoyer) & \eng{The teacher uploads notes.} \\
\hline
a network & un réseau & \eng{The network is slow today.} \\
\hline
online / offline & en ligne / hors ligne & \eng{I study online every evening.} \\
\hline
\end{tabularx}
\end{center}

\begin{attention}[Faux amis et calques à éviter]
\begin{itemize}
\item \eng{an application} ne signifie pas « l'application » au sens de « le soin » ; c'est bien le programme informatique — mais attention, « postuler » se dit \eng{to apply}, pas « to application ».
\item Ne dites pas \emph{« I connected my phone on internet »} : dites \eng{I connected my phone to the Internet} (préposition \emph{to}, article \emph{the}).
\item « Télécharger » a deux directions en anglais : \eng{download} (vers vous) et \eng{upload} (vers le réseau). Les francophones confondent souvent les deux.
\end{itemize}
\end{attention}

\begin{exoresolu}[Classify the words]
Énoncé — Classez les mots suivants dans la bonne colonne : \emph{hardware}, \emph{software} ou \emph{service} : \eng{a keyboard, WhatsApp, the Internet, a printer, a browser, e-mail, a laptop, a word processor, a mobile network}.
\tcblower
Solution détaillée —
\begin{itemize}
\item \textbf{Hardware} (objets physiques) : \eng{a keyboard, a printer, a laptop} — on peut les toucher.
\item \textbf{Software} (programmes) : \eng{WhatsApp} (une application), \eng{a browser} (un navigateur), \eng{a word processor} (un traitement de texte).
\item \textbf{Services} (connexions et plateformes) : \eng{the Internet, e-mail, a mobile network} — ils relient les utilisateurs entre eux.
\end{itemize}
Astuce de vérification : posez la question « Can I touch it? » Si oui → hardware. Si c'est un programme qui tourne sur l'appareil → software. Si cela connecte des personnes ou stocke à distance → service.
\end{exoresolu}

\begin{exercice}[Your turn: from context to meaning]
Relevez dans le texte \emph{How ICTs Changed My School Year} cinq mots liés aux TIC que vous ne connaissiez pas. Pour chacun, appliquez la méthode en cinq étapes (nature, indices, hypothèse, vérification), puis utilisez le mot dans une phrase personnelle qui parle de votre vie scolaire à Yaoundé, Douala ou ailleurs au Cameroun.
\end{exercice}

\begin{corrige}[Exercice : From context to meaning]
Exemple de production attendue avec le mot \eng{tutorial} : nature = nom (précédé de \emph{video}) ; indices = \eng{watching... on my phone}, \eng{teachers explain mathematics} → hypothèse : « une leçon en vidéo » ; vérification : glossaire = tutoriel. Phrase personnelle : \eng{Last week I watched a tutorial about essay writing before my English test.} « La semaine dernière, j'ai regardé un tutoriel sur la rédaction avant mon test d'anglais. » Toute phrase grammaticalement correcte réemployant un mot du texte dans un contexte personnel est acceptée.
\end{corrige}

\begin{savaistu}[ICTs in Cameroon]
Le Cameroun compte plus de dix millions d'utilisateurs d'internet mobile : le smartphone est devenu le premier « ordinateur » du pays. Le sigle \eng{ICT} apparaît partout dans les documents officiels du MINESEC : au lycée, \eng{ICT} désigne tout simplement l'informatique scolaire, et l'épreuve pratique sur ordinateur s'appelle souvent \eng{the ICT practical}. Maîtriser ce vocabulaire en anglais, c'est donc aussi mieux comprendre votre propre programme scolaire !
\end{savaistu}

\begin{aretenir}[What to remember]
\begin{itemize}
\item \cle{ICTs} = Information and Communication Technologies : digital tools used to create, store, share and exchange information.
\item Trois familles : \cle{hardware} (matériel), \cle{software} (logiciels), \cle{services} (réseaux et plateformes).
\item \eng{Hardware}, \eng{software} et \eng{information} sont invariables : jamais de \emph{s}.
\item \eng{Download} (recevoir) et \eng{upload} (envoyer) ne se confondent pas.
\item Pour deviner un mot inconnu : phrase entière → nature → indices du contexte → hypothèse → vérification.
\end{itemize}
\end{aretenir}

\coursec{Lexical field 1: Devices, tools and the Internet}

\courssub{Transition : des usages aux outils}
Dans la section précédente, nous avons découvert comment les TIC (Technologies de l'Information et de la Communication) transforment notre quotidien au Cameroun : paiement mobile (\eng{Mobile Money}), cours en ligne, visioconférences entre Yaoundé et Douala, ou encore recherche d'informations pour un devoir. Mais derrière chaque usage se cache un \textbf{outil} (matériel ou logiciel) et une \textbf{action} précise. Pour parler de ces outils avec aisance — à l'oral comme à l'écrit — il faut maîtriser le vocabulaire de base : le \textbf{matériel} (\eng{hardware}), les \textbf{logiciels} (\eng{software}), le \textbf{réseau} et les \textbf{verbes d'action} qui vont avec. C'est l'objet de cette section.

\courssub{Le matériel (Hardware) : ce qu'on peut toucher}

\begin{textebox}[Dialogue au Cyber Café « Digital Hub » à Buea]
\textbf{Customer:} Good morning. I need to type and print my CV. Do you have a \textbf{computer} available?
\textbf{Attendant:} Yes, we have \textbf{desktops} and \textbf{laptops}. The \textbf{desktops} have large \textbf{screens} and comfortable \textbf{keyboards}. The \textbf{laptops} are portable but the \textbf{trackpad} can be tricky if you are used to a \textbf{mouse}.
\textbf{Customer:} I'll take a desktop. I also need to scan my diploma. Is the \textbf{scanner} working?
\textbf{Attendant:} Yes, it's connected to the \textbf{printer} — it's an all-in-one device. Don't forget to save your file on a \textbf{USB flash drive} or your \textbf{hard disk} before you leave.
\textbf{Customer:} Thanks. Oh, my phone battery is dead. Do you have a \textbf{charger} for a \textbf{smartphone}?
\textbf{Attendant:} We have universal chargers at the counter. You can also use the \textbf{projector} in the meeting room if you need to present something later.
\end{textebox}

\begin{center}
\begin{tabularx}{\linewidth}{|X|X|X|}
\hline
\rowcolor{popblueL} \textbf{English} & \textbf{Français} & \textbf{Remarque / Prononciation} \\
\hline
a \textbf{computer} / a \textbf{desktop computer} & un ordinateur (fixe) & \eng{desktop} = ordinateur de bureau (ne bouge pas) \\
a \textbf{laptop} & un ordinateur portable & aussi \eng{notebook} ; \eng{laptop} = « sur les genoux » \\
a \textbf{tablet} & une tablette & tactile, pas de clavier physique intégré \\
a \textbf{smartphone} & un smartphone & téléphone intelligent ; \eng{phone} ou \eng{mobile} (UK) / \eng{cell phone} (US) \\
a \textbf{keyboard} & un clavier & \eng{to type on the keyboard} (taper au clavier) \\
a \textbf{mouse} (pl. \textbf{mice}) & une souris & \eng{trackpad} / \eng{touchpad} = pavé tactile (sur laptop) \\
a \textbf{screen} / a \textbf{monitor} & un écran & \eng{screen} plus général (téléphone, tablette aussi) \\
a \textbf{printer} & une imprimante & \eng{to print} (imprimer) \\
a \textbf{scanner} & un scanner & \eng{to scan} (numériser) \\
a \textbf{projector} & un vidéoprojecteur & \eng{to project} (projeter) \\
a \textbf{USB flash drive} / a \textbf{flash drive} & une clé USB & \textbf{pas} \eng{a USB key} (calque fr.) \\
a \textbf{hard disk} / a \textbf{hard drive} & un disque dur & \eng{external hard drive} = disque dur externe \\
a \textbf{charger} & un chargeur & \eng{to charge} (recharger) ; \eng{battery} = batterie \\
a \textbf{workstation} & un poste de travail & ordinateur puissant en réseau (ex. salle informatique) \\
\hline
\end{tabularx}
\end{center}

\begin{attention}[Pièges classiques « Hardware »]
\begin{itemize}
\item \textbf{Computer vs Ordinator} : « Un ordinateur » se dit \eng{a computer} (ou \eng{a desktop computer}), \textbf{jamais} \eng{an ordinator} (mot inexistant).
\item \textbf{USB key vs Flash drive} : « Une clé USB » se dit \eng{a USB flash drive}, \eng{a flash drive} ou \eng{a thumb drive} (US). \eng{A USB key} est un calque du français à éviter.
\item \textbf{Mouse → Mice} : Le pluriel de \eng{mouse} est \eng{mice} (irrégulier), pas \eng{mouses}.
\item \textbf{Smartphone} : Un seul mot, pas \eng{smart phone} (séparé).
\item \textbf{Charger} : Le verbe est \eng{to charge} (recharger la batterie), pas \eng{to load} (charger un camion / un logiciel).
\end{itemize}
\end{attention}

\begin{savaistu}[Vocabulaire UK vs US : la clé USB]
Au \textbf{Royaume-Uni}, on entend très souvent \eng{memory stick} pour désigner une clé USB. Aux \textbf{États-Unis}, les termes les plus courants sont \eng{flash drive} ou \eng{thumb drive} (litt. « clé pouce » car de la taille d'un pouce). Au Cameroun, \eng{flash drive} est le plus compris dans les cybercafés et les magasins d'informatique.
\end{savaistu}

\courssub{Les logiciels et applications (Software) : ce qu'on utilise}

\begin{definition}[Software vs Hardware]
Le \textbf{hardware} (matériel) désigne les composants \textbf{physiques}, tangibles (clavier, écran, disque dur). Le \textbf{software} (logiciel) désigne les \textbf{programmes}, instructions et données qui font fonctionner le matériel. Sans software, le hardware est inerte ; sans hardware, le software ne peut s'exécuter.
\end{definition}

\begin{textebox}[Extrait d'un guide de configuration pour nouveaux élèves (Lycée de Buea)]
Welcome to the school computer lab. Each \textbf{workstation} runs on \textbf{Windows 11} as its \textbf{operating system}. You will find the \textbf{Microsoft Office} suite installed: \textbf{Word} (a \textbf{word processor}) for your essays, \textbf{Excel} (a \textbf{spreadsheet}) for data analysis in Geography or Economics, and \textbf{PowerPoint} for presentations. To access the internet, open a \textbf{browser} like \textbf{Chrome} or \textbf{Firefox}. The default \textbf{search engine} is Google. An \textbf{antivirus} runs in the background to protect the \textbf{network}. Remember to log out of your \textbf{account} before leaving.
\end{textebox}

\begin{center}
\begin{tabularx}{\linewidth}{|X|X|X|}
\hline
\rowcolor{poppurpleL} \textbf{English} & \textbf{Français} & \textbf{Collocation / Note} \\
\hline
a \textbf{program} / an \textbf{application} (\textbf{app}) & un programme / une application & \eng{app} = abréviation courante (mobile surtout) \\
a \textbf{word processor} & un traitement de texte & ex. \eng{Microsoft Word}, \eng{Google Docs}, \eng{LibreOffice Writer} \\
a \textbf{spreadsheet} & un tableur & ex. \eng{Excel}, \eng{Google Sheets} ; \eng{spreadsheet software} \\
a \textbf{browser} & un navigateur (web) & \textbf{ne pas confondre avec le verbe \eng{to browse}} \\
a \textbf{search engine} & un moteur de recherche & ex. \eng{Google}, \eng{Bing}, \eng{DuckDuckGo} \\
an \textbf{antivirus} & un antivirus & \eng{antivirus software} ; \eng{to run a scan} \\
an \textbf{operating system} (\textbf{OS}) & un système d'exploitation & ex. \eng{Windows}, \eng{macOS}, \eng{Linux}, \eng{Android}, \eng{iOS} \\
\hline
\end{tabularx}
\end{center}

\begin{definition}[Browser (Navigateur)]
Un \textbf{browser} (ou \textbf{web browser}) est \textbf{a program used to visit websites} (un programme utilisé pour visiter des sites web). Exemples : \eng{Google Chrome}, \eng{Mozilla Firefox}, \eng{Microsoft Edge}, \eng{Safari}.
\par\medskip
\textbf{Attention :} Ne pas confondre le nom \eng{browser} (le logiciel) et le verbe \eng{to browse} (feuilleter, naviguer sur le web, regarder sans but précis).
\par\medskip
\eng{I use Chrome as my default browser.} = J'utilise Chrome comme navigateur par défaut. \\
\eng{She browsed the internet for hours.} = Elle a navigué sur internet pendant des heures.
\end{definition}

\begin{exemplebox}[Exemples traduits — Software]
\eng{Open your word processor and type the report.} = Ouvre ton traitement de texte et tape le rapport.
\eng{The spreadsheet calculates the average automatically.} = Le tableur calcule la moyenne automatiquement.
\eng{Which browser do you prefer: Chrome or Firefox?} = Quel navigateur préfères-tu : Chrome ou Firefox ?
\eng{The antivirus detected a threat and removed it.} = L'antivirus a détecté une menace et l'a supprimée.
\eng{You need to update your operating system.} = Tu dois mettre à jour ton système d'exploitation.
\end{exemplebox}

\courssub{Internet et réseaux : se connecter et naviguer}

\begin{textebox}[Échange WhatsApp entre deux élèves (Yaoundé / Douala)]
\textbf{Alice (Yaoundé):} Hey! The \textbf{Wi-Fi} at home is down. I can't join the Zoom class. No \textbf{connection} at all.
\textbf{Paul (Douala):} Use your phone's \textbf{hotspot}. Buy a \textbf{data bundle} from MTN or Orange. It's faster.
\textbf{Alice:} Good idea. I'll \textbf{connect} to 4G. What's the \textbf{link} for the class?
\textbf{Paul:} Check the \textbf{website} of the school. \textbf{Username}: student2025. \textbf{Password}: Buea2025! (case sensitive).
\textbf{Alice:} Thanks. I'll \textbf{log in} now. My \textbf{account} was locked yesterday.
\textbf{Paul:} Click \textbf{"Forgot password"} on the \textbf{web page} if needed. See you online!
\end{textebox}

\begin{center}
\begin{tabularx}{\linewidth}{|X|X|X|}
\hline
\rowcolor{popgreenL} \textbf{English} & \textbf{Français} & \textbf{Contexte / Collocation} \\
\hline
the \textbf{web} / the \textbf{Internet} & le web / internet & \eng{on the web}, \eng{access the Internet} \\
a \textbf{website} & un site web & \eng{visit a website}, \eng{the school website} \\
a \textbf{web page} & une page web & une page \textbf{au sein} d'un site \\
a \textbf{link} / a \textbf{URL} & un lien / une URL & \eng{click on the link}, \eng{copy the link} \\
a \textbf{username} & un nom d'utilisateur & \eng{enter your username} \\
a \textbf{password} & un mot de passe & \eng{strong password} (mot de passe robuste) \\
an \textbf{account} & un compte & \eng{create an account}, \eng{log into your account} \\
\textbf{Wi-Fi} & le Wi-Fi & \eng{connect to Wi-Fi}, \eng{Wi-Fi password} \\
a \textbf{network} & un réseau & \eng{local network}, \eng{mobile network} \\
a \textbf{connection} & une connexion & \eng{internet connection}, \eng{slow/fast connection} \\
a \textbf{data bundle} / a \textbf{data plan} & un forfait data / un paquet data & \eng{buy a data bundle}, \eng{run out of data} \\
a \textbf{hotspot} & un point d'accès (partage de connexion) & \eng{personal hotspot}, \eng{mobile hotspot} \\
the \textbf{cloud} & le cloud (stockage en ligne) & \eng{save to the cloud}, \eng{cloud storage} \\
\hline
\end{tabularx}
\end{center}

\begin{exemplebox}[Exemples traduits — Internet \& Réseau]
\eng{I typed my password and clicked on the link.} = J'ai tapé mon mot de passe et j'ai cliqué sur le lien.
\eng{The Wi-Fi connection is very slow today.} = La connexion Wi-Fi est très lente aujourd'hui.
\eng{He forgot his username, so he couldn't log in.} = Il a oublié son nom d'utilisateur, donc il n'a pas pu se connecter.
\eng{She bought a 5~GB data bundle for the week.} = Elle a acheté un forfait data de 5~Go pour la semaine.
\eng{The website is down; the server is not responding.} = Le site est en panne ; le serveur ne répond pas.
\end{exemplebox}

\courssub{Actions de base (Verbes) : interagir avec la machine}

\begin{propriete}[Verbes d'action ICT — Forme et Emploi]
Ces verbes sont au cœur de toute consigne informatique. La plupart sont \textbf{réguliers} au prétérit (\eng{-ed}) sauf \textbf{to type} → \eng{typed}.
\begin{center}
\begin{tabular}{|l|l|l|}
\hline
\rowcolor{poporangeL} \textbf{Verbe (Base)} & \textbf{Prétérit / Participe} & \textbf{Sens / Contexte typique} \\
\hline
to \textbf{switch on} / \textbf{turn on} & switched on / turned on & allumer (appareil) — \eng{switch off/turn off} = éteindre \\
to \textbf{plug in} & plugged in & brancher (prise, câble, clé USB) — \eng{unplug} = débrancher \\
to \textbf{type} & typed & taper au clavier — \eng{type your password} \\
to \textbf{click} & clicked & cliquer (souris) — \eng{click on the icon/link/button} \\
to \textbf{double-click} & double-clicked & double-cliquer — \eng{double-click the icon} \\
to \textbf{scroll} & scrolled & faire défiler (page, liste) — \eng{scroll down/up} \\
to \textbf{tap} & tapped & toucher/tapper (écran tactile) — \eng{tap the screen} \\
to \textbf{connect} & connected & connecter (au Wi-Fi, à un serveur) — \eng{connect to} \\
to \textbf{disconnect} & disconnected & déconnecter — \eng{disconnect from} \\
to \textbf{log in} / \textbf{sign in} & logged in / signed in & ouvrir une session — \eng{log in to your account} \\
to \textbf{log out} / \textbf{sign out} & logged out / signed out & fermer une session — \eng{log out of the website} \\
to \textbf{download} & downloaded & télécharger (vers son appareil) — \eng{download a file} \\
to \textbf{upload} & uploaded & téléverser / envoyer (vers le cloud, un site) — \eng{upload a photo} \\
to \textbf{install} & installed & installer (logiciel, app, mise à jour) — \eng{install an update} \\
to \textbf{update} & updated & mettre à jour — \eng{update the OS / the app} \\
\hline
\end{tabular}
\end{center}
\end{propriete}

\begin{attention}[Faux amis et calques fréquents — Verbes]
\begin{itemize}
\item \textbf{Download vs Upload} : \eng{Download} = télécharger \textbf{vers le bas} (depuis internet vers ton appareil). \eng{Upload} = téléverser \textbf{vers le haut} (depuis ton appareil vers internet/cloud). En français « télécharger » couvre souvent les deux ; en anglais, la distinction est stricte.
\item \textbf{Log in vs Log on} : \eng{Log in} (ou \eng{sign in}) est le standard pour une session utilisateur (compte, site). \eng{Log on} s'emploie plutôt pour un réseau ou un système (ex. \eng{log on to the domain}).
\item \textbf{Click on} : Préposition \textbf{on} obligatoire après \eng{click} quand on désigne la cible (\eng{click on the link}). Pas \eng{click the link} (même si entendu à l'oral, \eng{click on} est plus correct en contexte formel).
\item \textbf{Type} : Pas de préposition après \eng{type} pour le contenu (\eng{type your name}), mais \eng{type on the keyboard}.
\item \textbf{Switch on/off} : Verbe à particule (phrasal verb). Objet pronom \textbf{entre} les deux : \eng{switch it on}, \eng{turn it off}. Pas \eng{switch on it}.
\end{itemize}
\end{attention}

\courssub{Prononciation : les pièges sonores}

\begin{aretenir}[Prononciation clé — Approximation française]
Ne pas lire l'anglais comme du français. Entraîne-toi à dire ces mots correctement :
\begin{center}
\begin{tabular}{|l|l|l|}
\hline
\rowcolor{poppinkL} \textbf{Mot} & \textbf{API (réf.)} & \textbf{Approximation française} \\
\hline
\textbf{data} & /ˈdeɪtə/ & « \textbf{déï-ta} » (diphtongue \eng{ei} comme dans \eng{day}, pas « da-ta ») \\
\textbf{device} & /dɪˈvaɪs/ & « \textbf{di-vaïss} » (accent sur 2\textsuperscript{e} syllabe, \eng{ice} se prononce \eng{aïss}) \\
\textbf{software} & /ˈsɒftwɛə(r)/ & « \textbf{sof-touèr} » (deux mots joints, \eng{ware} = \eng{touèr}) \\
\textbf{website} & /ˈwɛbsaɪt/ & « \textbf{ouèb-saït} » (pas « ouèb-site » ; \eng{web} = \eng{ouèb}) \\
\textbf{Wi-Fi} & /ˈwaɪfaɪ/ & « \textbf{ouaï-faï} » (deux diphtongues \eng{ai} = \eng{aï}) \\
\textbf{router} & /ˈruːtə(r)/ (US) /ˈraʊtə(r)/ (UK) & US : « \textbf{routeur} » (ou « raou-teur ») ; UK : « \textbf{raou-teur} » \\
\textbf{USB} & /ˌjuː es ˈbiː/ & « \textbf{iou-esse-bi} » (épeler les lettres) \\
\hline
\end{tabular}
\end{center}
\end{aretenir}

\begin{methode}[Comment travailler la prononciation ICT]
\begin{enumerate}
\item \textbf{Écoute} : Trouve la prononciation sur un dictionnaire en ligne (Oxford, Cambridge, Forvo) pour chaque mot du tableau ci-dessus.
\item \textbf{Répète} : Fais 3 répétitions à voix haute en exagérant les diphtongues (\eng{data}, \textbf{device}, \textbf{Wi-Fi}).
\item \textbf{Enregistre} : Enregistre-toi sur ton smartphone (\eng{voice recorder}) et compare.
\item \textbf{Contextualise} : Dis une phrase complète par mot (ex. \eng{The data is on the device.} → « Le déï-ta is on the di-vaïss. »).
\item \textbf{Révise} : Refais l'exercice 2 jours plus tard (répétition espacée).
\end{enumerate}
\end{methode}

\courssub{Synthèse visuelle : Hardware vs Software}

\begin{popfigure}
\begin{tikzpicture}[
  node distance=1.2cm and 2cm,
  main/.style={rectangle, rounded corners, draw=popdark, thick, fill=popcream, minimum width=3.5cm, minimum height=1.2cm, align=center, font=\bfseries},
  cat/.style={rectangle, rounded corners, draw=popblue, fill=popblueL, minimum width=3cm, minimum height=1cm, align=center, font=\small},
  ex/.style={rectangle, draw=popgreen, fill=popgreenL, minimum width=2.8cm, minimum height=0.8cm, align=center, font=\small\itshape},
  arrow/.style={-Stealth, thick, popdark}
]
% Nœuds principaux
\node[main, align=center] (hw){HARDWARE\\(Matériel)};
\node[main, right=of hw, align=center] (sw){SOFTWARE\\(Logiciels)};

% Catégories Hardware
\node[cat, below=of hw, xshift=-1.5cm, align=center] (hw1){Input Devices\\(Périph. entrée)};
\node[cat, below=of hw, xshift=1.5cm, align=center] (hw2){Output Devices\\(Périph. sortie)};
\node[cat, below=of hw1, yshift=-0.5cm, align=center] (hw3){Storage\\(Stockage)};
\node[cat, below=of hw2, yshift=-0.5cm, align=center] (hw4){Network/Periph.\\(Réseau / Autre)};

% Exemples Hardware
\node[ex, left=of hw1, align=center] (ex1){Keyboard, Mouse,\\Scanner, Microphone};
\node[ex, right=of hw2, align=center] (ex2){Screen, Printer,\\Projector, Speakers};
\node[ex, left=of hw3, align=center] (ex3){Hard Disk, USB\\Flash Drive, SSD};
\node[ex, right=of hw4, align=center] (ex4){Router, Modem,\\Charger, Webcam};

% Catégories Software
\node[cat, below=of sw, xshift=-1.5cm, align=center] (sw1){System Software\\(Système)};
\node[cat, below=of sw, xshift=1.5cm, align=center] (sw2){Application Software\\(Applications)};

% Exemples Software
\node[ex, left=of sw1, align=center] (ex5){Windows, macOS,\\Linux, Android, iOS,\\Drivers, Antivirus};
\node[ex, right=of sw2, align=center] (ex6){Word, Excel,\\Chrome, Firefox,\\Zoom, WhatsApp,\\Games, Apps};

% Flèches
\draw[arrow] (hw) -- (hw1);
\draw[arrow] (hw) -- (hw2);
\draw[arrow] (hw) -- (hw3);
\draw[arrow] (hw) -- (hw4);
\draw[arrow] (sw) -- (sw1);
\draw[arrow] (sw) -- (sw2);

\end{tikzpicture}
\legende{Carte mentale : Hardware (objets tangibles) vs Software (programmes intangibles) — Flèches vers catégories et exemples.}
\end{popfigure}

\courssub{Exercice résolu : Mettre en pratique}

\begin{exoresolu}[Compléter un dialogue technique]
\textbf{Énoncé :} Complète le dialogue ci-dessous avec les mots de la liste. Attention : certains mots peuvent servir deux fois. Traduis ensuite le dialogue en français.
\par\medskip
\noindent\textbf{Liste :} \eng{password / link / browser / click / connect / Wi-Fi / username / log in / website / type}
\par\medskip
\noindent\textbf{Alice:} I can't access the school \_\_\_\_\_\_\_\_. The page won't load.
\textbf{Paul:} Check your \_\_\_\_\_\_\_\_ signal. Are you \_\_\_\_\_\_\_\_ to the internet?
\textbf{Alice:} Yes, I'm connected. But when I \_\_\_\_\_\_\_\_ on the \_\_\_\_\_\_\_\_ you sent me, nothing happens.
\textbf{Paul:} Which \_\_\_\_\_\_\_\_ are you using? Try Chrome or Firefox.
\textbf{Alice:} I'm on Edge. Okay, I'll open Chrome. Now I need to \_\_\_\_\_\_\_\_ my \_\_\_\_\_\_\_\_ and \_\_\_\_\_\_\_\_.
\textbf{Paul:} Exactly. \_\_\_\_\_\_\_\_ your \_\_\_\_\_\_\_\_ first, then your \_\_\_\_\_\_\_\_. Then \_\_\_\_\_\_\_\_.
\par\medskip
\tcblower
\textbf{Solution détaillée :}
\begin{enumerate}
\item \textbf{website} (contexte : « page won't load » → site web). \textit{Traduction : site web}
\item \textbf{Wi-Fi} (signal sans fil). \textit{Traduction : Wi-Fi}
\item \textbf{connected} (verbe \eng{connect} au participe passé adjectival après \eng{be}). \textit{Traduction : connecté(e)}
\item \textbf{click} (verbe à l'impératif/présent narration : \eng{when I click}). \textit{Traduction : clique}
\item \textbf{link} (nom : le lien envoyé). \textit{Traduction : lien}
\item \textbf{browser} (nom : logiciel pour naviguer). \textit{Traduction : navigateur}
\item \textbf{type} (verbe : action de saisir au clavier). \textit{Traduction : tape / saisis}
\item \textbf{username} (nom : identifiant). \textit{Traduction : nom d'utilisateur}
\item \textbf{password} (nom : mot de passe). \textit{Traduction : mot de passe}
\item \textbf{log in} (verbe phrasal : ouvrir la session). \textit{Traduction : connecte-toi / ouvre ta session}
\end{enumerate}
\par\medskip
\textbf{Traduction du dialogue :}
\par\medskip
\noindent\textbf{Alice :} Je n'arrive pas à accéder au site web de l'école. La page ne se charge pas.
\textbf{Paul :} Vérifie ton signal Wi‑Fi. Es-tu connecté à internet ?
\textbf{Alice :} Oui, je suis connectée. Mais quand je clique sur le lien que tu m'as envoyé, il ne se passe rien.
\textbf{Paul :} Quel navigateur utilises-tu ? Essaie Chrome ou Firefox.
\textbf{Alice :} Je suis sur Edge. D'accord, j'ouvre Chrome. Maintenant je dois taper mon nom d'utilisateur et mon mot de passe.
\textbf{Paul :} Exactement. Tape ton nom d'utilisateur d'abord, puis ton mot de passe. Puis connecte-toi.
\end{exoresolu}

\begin{exercice}[Entraînement personnel — À faire dans ton cahier]
\begin{enumerate}
\item \textbf{Traduction :} Traduis en anglais.
\begin{itemize}
\item[a)] J'ai branché ma clé USB dans l'ordinateur portable.
\item[b)] Le vidéoprojecteur ne fonctionne pas ; il faut l'allumer.
\item[c)] Elle a fait défiler la page web vers le bas pour lire l'article.
\item[d)] Nous devons mettre à jour l'antivirus et le système d'exploitation.
\item[e)] Il a oublié son mot de passe, alors il a cliqué sur « Mot de passe oublié ».
\end{itemize}
\item \textbf{Choix multiple :} Choisis la bonne réponse.
\begin{itemize}
\item[1.] To save a file on a \_\_\_\_\_\_, you plug it into the USB port. \quad (a) hard disk \quad (b) USB flash drive \quad (c) scanner \quad (d) projector
\item[2.] The \_\_\_\_\_\_ manages the computer's hardware and runs applications. \quad (a) browser \quad (b) word processor \quad (c) operating system \quad (d) search engine
\item[3.] \_\_\_\_\_\_ the icon twice to open the program. \quad (a) Type \quad (b) Scroll \quad (c) Tap \quad (d) Double-click
\item[4.] I need to \_\_\_\_\_\_ this photo to the cloud so my colleague can see it. \quad (a) download \quad (b) upload \quad (c) install \quad (d) scan
\end{itemize}
\item \textbf{Prononciation :} Écris l'approximation française pour : \eng{router}, \eng{antivirus}, \eng{spreadsheet}, \eng{username}, \eng{hotspot}.
\end{enumerate}
\end{exercice}

\begin{corrige}[Exercice entraînement personnel]
\textbf{1. Traduction :}
\begin{itemize}
\item[a)] I plugged my USB flash drive into the laptop. (\eng{plugged in} aussi possible)
\item[b)] The projector isn't working; we need to switch it on / turn it on.
\item[c)] She scrolled down the web page to read the article.
\item[d)] We need to update the antivirus and the operating system.
\item[e)] He forgot his password, so he clicked on "Forgot password".
\end{itemize}
\textbf{2. Choix multiple :}
\begin{itemize}
\item[1.] (b) USB flash drive — \eng{hard disk} = disque dur (interne/externe), pas clé USB.
\item[2.] (c) operating system — Le \eng{browser} navigue, le \eng{word processor} traite le texte, le \eng{search engine} cherche.
\item[3.] (d) Double-click — Action standard pour ouvrir (souris). \eng{Tap} sur tactile, \eng{click} simple = sélectionner.
\item[4.] (b) upload — Envoyer vers le cloud/serveur. \eng{Download} = recevoir.
\end{itemize}
\textbf{3. Prononciation (approximation) :}
\begin{itemize}
\item \eng{router} : « raou-teur » (UK) / « routeur » (US)
\item \eng{antivirus} : « an-ti-vaï-russe »
\item \eng{spreadsheet} : « spred-chite » (\eng{ea} = \eng{è}, \eng{sh} = \eng{ch})
\item \eng{username} : « you-zer-neim »
\item \eng{hotspot} : « hote-spote » (\eng{h} aspiré, \eng{o} court)
\end{itemize}
\end{corrige}

\courssub{Bilan de la section 1}
Tu disposes maintenant du socle lexical \textbf{Hardware / Software / Internet / Actions} pour décrire un environnement de travail ou d'apprentissage numérique. Les points de vigilance (faux amis \eng{computer/ordinator}, \eng{USB key/flash drive}, \eng{download/upload}, prononciation \eng{data/device}) sont les marqueurs qui distinguent un apprenant approximatif d'un utilisateur compétent. Dans la section suivante, nous verrons comment ce vocabulaire s'applique spécifiquement aux \textbf{usages pédagogiques} (apprentissage en ligne, plateformes, ressources éducatives).

\coursec{Lexical field 2: ICTs for learning}

\courssub{Transition : des outils aux usages pédagogiques}

Dans la section précédente, nous avons identifié le \textbf{matériel} (devices) et l'\textbf{infrastructure} (Internet, Wi-Fi, cloud) qui rendent possibles les technologies de l'information et de la communication. Nous savons désormais nommer un \eng{smartphone}, un \eng{laptop}, un \eng{router} ou une \eng{USB drive}. Mais posséder les outils ne suffit pas : encore faut-il savoir les \textbf{utiliser pour apprendre}. C'est tout l'objet de ce champ lexical 2 : passer du \eng{hardware} au \eng{software} pédagogique, de la machine à la méthode. Nous allons explorer le vocabulaire de l'\eng{e-learning}, des ressources numériques, des actions d'apprentissage autonomes ou collaboratives, et des acteurs de cette nouvelle école sans murs.

\begin{textebox}{Mini-dialogue : Preparing for the Baccalaureate}
\textbf{A:} How do you prepare for the Bac? I feel overwhelmed by the syllabus.\\
\textbf{B:} I mainly use \textbf{e-learning} platforms. I \textbf{watch video tutorials} on YouTube and \textbf{download past papers} from the school website.\\
\textbf{A:} That's smart! Does it really help?\\
\textbf{B:} Absolutely. I can \textbf{learn at my own pace}, pause the video, and rewind. Plus, I \textbf{join an online study group} on WhatsApp where we \textbf{share notes} and \textbf{discuss difficult topics}.\\
\textbf{A:} I prefer studying alone with my textbooks. Maybe I should try a \textbf{webinar} or an \textbf{educational app}.\\
\textbf{B:} You should! There are great \textbf{digital libraries} and \textbf{online dictionaries} too. It \textbf{saves time} and you \textbf{access information easily}.
\end{textebox}

\noindent
\emph{Glossaire :} \textbf{overwhelmed} (adj.) = submergé, dépassé ; \textbf{syllabus} (n.) = programme ; \textbf{platform} (n.) = plateforme ; \textbf{pause} (v.) = mettre en pause ; \textbf{rewind} (v.) = rembobiner, revenir en arrière ; \textbf{discuss} (v.) = discuter ; \textbf{textbook} (n.) = manuel scolaire.

\begin{aretenir}{Observation guidée}
Lisez le dialogue ci-dessus et relevez :
\begin{enumerate}
    \item Trois noms composés désignant des \textbf{modalités d'apprentissage} (ex. \eng{e-learning}).
    \item Quatre verbes d'\textbf{action} suivis d'un complément d'objet (ex. \eng{watch video tutorials}).
    \item Trois expressions d'\textbf{avantage} introduites par \eng{to} + infinitif (ex. \eng{learn at my own pace}).
    \item Deux noms désignant des \textbf{personnes} ou des \textbf{rôles} (ex. \eng{content creator} — absent ici, à anticiper).
\end{enumerate}
\end{aretenir}

\courssub{Les modalités de l'apprentissage en ligne (E-learning modalities)}

L'anglais utilise souvent des \textbf{noms composés} (formés de deux noms ou préfixe + nom) pour désigner les nouvelles réalités pédagogiques. Le préfixe \eng{e-} (pour \eng{electronic}) et \eng{online} / \eng{virtual} / \eng{digital} sont très productifs.

\begin{definition}[E-learning et termes apparentés]
\textbf{E-learning} (ou \eng{electronic learning}) : apprentissage réalisé via des supports électroniques, généralement sur Internet.\\
\textbf{Distance learning} / \eng{Distance education} : enseignement à distance (l'élève et le professeur sont géographiquement séparés).\\
\textbf{Online course} : cours en ligne (souvent structuré, avec évaluation).\\
\textbf{Virtual classroom} : classe virtuelle (espace de visioconférence pour cours en direct).\\
\textbf{Video tutorial} (familier : \eng{tuto}) : tutoriel vidéo expliquant une procédure ou une leçon.\\
\textbf{Webinar} : contraction de \eng{web} + \eng{seminar}. Séminaire ou cours donné en direct sur Internet, souvent interactif (chat, sondages).\\
\textbf{Podcast} : fichier audio (ou vidéo) épisodique, téléchargeable, à écouter quand on veut.\\
\textbf{Educational app} : application pédagogique (sur smartphone/tablette).
\end{definition}

\begin{exemplebox}{Exemples en contexte}
\eng{E-learning allows students in remote areas to access quality education.}\\
(L'apprentissage en ligne permet aux élèves des zones reculées d'accéder à une éducation de qualité.)

\eng{The university offers a wide range of online courses for working professionals.}\\
(L'université propose une large gamme de cours en ligne pour les professionnels en activité.)

\eng{During the lockdown, all lessons took place in a virtual classroom on Zoom.}\\
(Pendant le confinement, tous les cours ont eu lieu dans une classe virtuelle sur Zoom.)

\eng{I learned how to code by watching video tutorials on YouTube.}\\
(J'ai appris à coder en regardant des tutoriels vidéo sur YouTube.)

\eng{She registered for a webinar on digital marketing next Tuesday.}\\
(Elle s'est inscrite à un webinaire sur le marketing digital mardi prochain.)

\eng{He listens to a history podcast every morning on the bus.}\\
(Il écoute un podcast d'histoire chaque matin dans le bus.)

\eng{Duolingo is a popular educational app for learning languages.}\\
(Duolingo est une application pédagogique populaire pour l'apprentissage des langues.)
\end{exemplebox}

\begin{center}
\begin{tabularx}{\linewidth}{|X|X|X|}
\hline
\rowcolor{popblueL} \textbf{English Term} & \textbf{Traduction française} & \textbf{Note / Collocation fréquente} \\
\hline
e-learning & apprentissage en ligne & \eng{e-learning platform}, \eng{e-learning module} \\
distance learning & enseignement à distance & \eng{distance learning programme} \\
online course & cours en ligne & \eng{take an online course}, \eng{enrol in an online course} \\
virtual classroom & classe virtuelle & \eng{enter the virtual classroom}, \eng{join a session} \\
video tutorial & tutoriel vidéo & \eng{watch a video tutorial}, \eng{follow a tutorial} \\
webinar & webinaire, séminaire en ligne & \eng{attend a webinar}, \eng{host a webinar} \\
podcast & podcast (balado) & \eng{listen to a podcast}, \eng{download a podcast} \\
educational app & application éducative & \eng{download an educational app} \\
\hline
\end{tabularx}
\end{center}

\courssub{Les ressources numériques pour l'apprenant (Digital learning resources)}

L'apprenant autonome dispose aujourd'hui d'une bibliothèque mondiale dans sa poche. Le vocabulaire suivant est indispensable pour décrire \emph{ce qu'on consulte}.

\begin{propriete}{Ressources clés : noms et collocations}
\begin{itemize}
    \item \textbf{Online dictionary} : dictionnaire en ligne (ex. \eng{Oxford Learner's Dictionaries}, \eng{Cambridge Dictionary}). \textit{Colloc.} : \eng{look up a word in an online dictionary}.
    \item \textbf{Encyclopaedia} (UK) / \textbf{Encyclopedia} (US) : encyclopédie. \eng{Wikipedia} est une \eng{online encyclopaedia}. \textit{Attention à l'orthographe UK/US : au Cameroun, on privilégie l'orthographe britannique \textbf{encyclopaedia}}.
    \item \textbf{Past papers} (UK) / \textbf{Past exams} (US) : annales, sujets d'examens des années précédentes. \textbf{Toujours au pluriel}. \textit{Colloc.} : \eng{download past papers}, \eng{practise with past papers}, \eng{revise using past papers}.
    \item \textbf{E-book} (ou \eng{eBook}) : livre numérique. \textit{Colloc.} : \eng{read an e-book on a tablet}, \eng{buy an e-book}.
    \item \textbf{Digital library} / \textbf{Online library} : bibliothèque numérique. \textit{Colloc.} : \eng{access the digital library}, \eng{borrow an e-book from the digital library}.
    \item \textbf{Notes} (pluriel) : notes de cours. \eng{Lecture notes} = notes de cours magistral. \eng{Revision notes} = fiches de révision. \textit{Ne pas confondre avec \eng{note} (note de musique, billet, ou note/grade)}.
    \item \textbf{Audio explanations} / \textbf{Video explanations} : explications audio/vidéo.
\end{itemize}
\end{propriete}

\begin{exemplebox}{Utilisation des ressources}
\eng{Before the exam, I always \textbf{download past papers} from the Ministry of Education website.}\\
(Avant l'examen, je télécharge toujours les annales depuis le site du Ministère de l'Éducation.)

\eng{She uses an \textbf{online dictionary} to check the pronunciation of difficult words.}\\
(Elle utilise un dictionnaire en ligne pour vérifier la prononciation des mots difficiles.)

\eng{Our school \textbf{digital library} gives free access to thousands of \textbf{e-books}.}\\
(La bibliothèque numérique de notre école donne un accès gratuit à des milliers de livres numériques.)

\eng{He shared his \textbf{revision notes} with the \textbf{online study group}.}\\
(Il a partagé ses fiches de révision avec le groupe d'étude en ligne.)
\end{exemplebox}

\courssub{Les verbes d'action : l'apprenant actif (Action verbs for learning)}

L'apprentissage numérique n'est pas passif. Il se décrit par des verbes d'action précis, souvent suivis de particules (verbes à particule / \eng{phrasal verbs}) ou de prépositions spécifiques.

\begin{propriete}{Verbes et locutions verbales essentiels}
\begin{center}
\begin{tabularx}{\linewidth}{|l|X|X|}
\hline
\rowcolor{popblueL} \textbf{Verbe / Locution} & \textbf{Sens / Emploi} & \textbf{Exemple traduit} \\
\hline
\textbf{to download} & Télécharger (vers son appareil) & \eng{I \textbf{downloaded} the PDF file.} (J'ai téléchargé le fichier PDF.) \\
\textbf{to upload} & Téléverser, mettre en ligne (vers un serveur/cloud) & \eng{She \textbf{uploaded} her assignment to the platform.} (Elle a déposé son devoir sur la plateforme.) \\
\textbf{to watch a tutorial} & Regarder un tutoriel & \eng{We \textbf{watched a tutorial} on Excel.} (Nous avons regardé un tuto sur Excel.) \\
\textbf{to revise online} & Réviser en ligne & \eng{He \textbf{revises online} using an app.} (Il révise en ligne avec une appli.) \\
\textbf{to do research on the Internet} & Faire des recherches sur Internet & \eng{They \textbf{did research on the Internet} for their project.} (Ils ont fait des recherches sur le Net pour leur projet.) \\
\textbf{to take notes} & Prendre des notes & \eng{Don't forget to \textbf{take notes} during the webinar.} (N'oublie pas de prendre des notes pendant le webinaire.) \\
\textbf{to share documents} & Partager des documents & \eng{We \textbf{share documents} via Google Drive.} (Nous partageons des docs via Google Drive.) \\
\textbf{to join a study group} & Rejoindre un groupe d'étude & \eng{I \textbf{joined a study group} on WhatsApp.} (J'ai rejoint un groupe d'étude sur WhatsApp.) \\
\textbf{to sign up for} / \textbf{to enrol in} & S'inscrire à (un cours) & \eng{She \textbf{signed up for} an online course.} (Elle s'est inscrite à un cours en ligne.) \\
\textbf{to log in / log out} & Se connecter / se déconnecter & \eng{Please \textbf{log in} to access the content.} (Connectez-vous pour accéder au contenu.) \\
\hline
\end{tabularx}
\end{center}
\end{propriete}

\begin{attention}{Piège majeur : \textbf{Research} est indénombrable !}
En français, on dit « faire \textbf{des} recherches » (pluriel). En anglais, \textbf{research} est un nom \textbf{indénombrable (uncountable)} dans le sens « travail de recherche, investigation ».
\begin{itemize}
    \item \textbf{Correct} : \eng{to do research}, \eng{to carry out research}, \eng{some research}, \eng{much research}, \eng{a piece of research}.
    \item \textbf{Incorrect} : \eng{to make researches} \textcolor{red}{✗}, \eng{to do a research} \textcolor{red}{✗}, \eng{many researches} \textcolor{red}{✗}.
\end{itemize}
\emph{Exception :} Au sens scientifique très formel, \eng{researches} (pluriel) peut exister pour désigner « plusieurs programmes d'études distincts » (ex. \eng{His researches into malaria}). Au lycée, \textbf{jamais} au pluriel.
\end{attention}

\begin{exemplebox}{Download vs Upload : la direction compte}
\eng{Download} = \textbf{Down} (vers le bas, vers soi) + \textbf{Load}. On récupère un fichier \textbf{depuis} un serveur \textbf{vers} son ordinateur/téléphone.
\eng{Upload} = \textbf{Up} (vers le haut, vers le nuage/serveur) + \textbf{Load}. On envoie un fichier \textbf{depuis} son appareil \textbf{vers} un serveur/cloud.

\eng{The teacher \textbf{uploaded} the lesson; the students \textbf{downloaded} it.}\\
(Le prof a mis la leçon en ligne ; les élèves l'ont téléchargée.)
\end{exemplebox}

\courssub{Les acteurs de l'apprentissage numérique (People in digital learning)}

Qui fait quoi dans cet écosystème ? Les noms d'agent en \textbf{-er} / \textbf{-or} ou les noms composés sont la norme.

\begin{definition}{Acteurs et rôles}
\begin{itemize}
    \item \textbf{Online tutor} / \textbf{E-tutor} : tuteur en ligne, enseignant qui accompagne à distance.
    \item \textbf{Learner} : apprenant (terme préféré en pédagogie moderne à \eng{student} ou \eng{pupil} pour souligner l'autonomie).
    \item \textbf{User} : utilisateur (terme technique, centré sur l'interface).
    \item \textbf{Blogger} : blogueur (qui tient un blog éducatif par ex.).
    \item \textbf{Content creator} : créateur de contenu (vidéos, PDF, quiz, podcasts éducatifs).
    \item \textbf{Moderator} / \textbf{Admin} : modérateur / administrateur (d'un forum, groupe d'étude, classe virtuelle).
\end{itemize}
\end{definition}

\begin{exemplebox}{Acteurs en phrases}
\eng{The \textbf{online tutor} provides feedback on every assignment.}\\
(Le tuteur en ligne donne un retour sur chaque devoir.)

\eng{As a \textbf{learner}, you are responsible for your own progress.}\\
(En tant qu'apprenant, tu es responsable de ta propre progression.)

\eng{Many \textbf{content creators} offer free lessons on YouTube.}\\
(De nombreux créateurs de contenu proposent des leçons gratuites sur YouTube.)

\eng{The \textbf{moderator} removed the off-topic messages from the forum.}\\
(Le modérateur a supprimé les messages hors sujet du forum.)
\end{exemplebox}

\courssub{Exprimer les avantages : structures à réemployer (Expressing benefits)}

Pour argumenter en faveur des TICE en éducation (à l'oral comme à l'écrit), maîtriser ces structures infinitives de finalité/avantage est crucial.

\begin{propriete}{Structures \eng{to} + infinitif pour exprimer l'avantage}
\begin{center}
\begin{tabularx}{\linewidth}{|l|X|X|}
\hline
\rowcolor{popgreenL} \textbf{Expression anglaise} & \textbf{Traduction} & \textbf{Structure grammaticale} \\
\hline
\textbf{to learn at your own pace} & apprendre à son propre rythme & \eng{to} + verbe + \eng{at} + possessif + \eng{pace} \\
\textbf{to access information easily} & accéder facilement à l'information & \eng{to} + verbe + objet + adverbe \\
\textbf{to save time} & gagner du temps / économiser du temps & \eng{to} + \eng{save} + \eng{time} (collocation forte) \\
\textbf{to improve one's grades} & améliorer ses notes / résultats & \eng{to} + \eng{improve} + possessif + \eng{grades/results} \\
\textbf{to study anytime, anywhere} & étudier n'importe quand, n'importe où & \eng{to} + verbe + adverbes de temps/lieu \\
\textbf{to develop digital skills} & développer des compétences numériques & \eng{to} + verbe + objet \\
\hline
\end{tabularx}
\end{center}
\end{propriete}

\begin{exemplebox}{Argumenter avec les avantages}
\eng{Using \textbf{educational apps} allows students \textbf{to learn at their own pace}.}\\
(L'utilisation d'applications éducatives permet aux élèves d'apprendre à leur rythme.)

\eng{\textbf{Digital libraries} help researchers \textbf{to access information easily}.}\\
(Les bibliothèques numériques aident les chercheurs à accéder facilement à l'information.)

\eng{\textbf{Downloading past papers} helps you \textbf{save time} during revision.}\\
(Télécharger les annales t'aide à gagner du temps pendant la révision.)

\eng{Regular practice on \textbf{e-learning platforms} can \textbf{improve your grades}.}\\
(Une pratique régulière sur les plateformes d'e-learning peut améliorer tes notes.)
\end{exemplebox}

\begin{aretenir}{Règle de construction : \textbf{Allow / Enable / Help} + Objet + \textbf{to} + Infinitif}
\begin{itemize}
    \item \eng{Allow someone \textbf{to} do something} (Permettre à qqn de faire qqch).
    \item \eng{Enable someone \textbf{to} do something} (Rendre qqn capable de...).
    \item \eng{Help someone (\textbf{to}) do something} (Aider qqn à faire qqch — \eng{to} facultatif après \eng{help}).
\end{itemize}
\textbf{Ne dites jamais :} \eng{Allow someone doing} \textcolor{red}{✗} ou \eng{Allow to someone to do} \textcolor{red}{✗}.
\end{aretenir}

\courssub{Familles de mots, orthographe et prononciation (Word families, spelling \& pronunciation)}

Travailler en famille lexicale permet de multiplier son vocabulaire par 3 ou 4 sans effort supplémentaire.

\begin{center}
\begin{tabular}{|l|l|l|l|}
\hline
\rowcolor{poppurpleL} \textbf{Verbe} & \textbf{Nom (action/état)} & \textbf{Nom (agent)} & \textbf{Adjectif / Participe} \\
\hline
\eng{learn} & \eng{learning} & \eng{learner} & \eng{learned} (US) / \eng{learnt} (UK) (adj. = savant) \\
\eng{teach} & \eng{teaching} & \eng{teacher} / \eng{tutor} & \eng{teachable} \\
\eng{study} & \eng{study} / \eng{studies} & \eng{student} & \eng{studious} \\
\eng{revise} (UK) / \eng{review} (US) & \eng{revision} (UK) / \eng{review} (US) & — & \eng{revised} / \eng{reviewed} \\
\eng{search} / \eng{research} & \eng{search} / \eng{research} & \eng{researcher} & \eng{searchable} / \eng{research-based} \\
\eng{download} & \eng{download} (n.) & \eng{downloader} (rare) & \eng{downloadable} \\
\eng{upload} & \eng{upload} (n.) & \eng{uploader} (rare) & \eng{uploadable} \\
\eng{connect} & \eng{connection} & \eng{user} & \eng{connected} / \eng{connectable} \\
\eng{interact} & \eng{interaction} & — & \eng{interactive} \\
\eng{collaborate} & \eng{collaboration} & \eng{collaborator} & \eng{collaborative} \\
\hline
\end{tabular}
\end{center}

\begin{attention}{Faux-amis et pièges d'orthographe/prononciation}
\begin{itemize}
    \item \textbf{Library} $\neq$ Librairie. \eng{Library} = \textbf{bibliothèque} (prêt gratuit). \eng{Bookshop} (UK) / \eng{Bookstore} (US) = \textbf{librairie} (vente). \textit{Prononciation :} « laï-bré-ri » (3 syllabes), pas « li-bra-ri ».
    \item \textbf{Course} $\neq$ Cour. \eng{Course} = \textbf{cours} (série de leçons), \textbf{parcours} (golf), \textbf{plat} (repas). \eng{Court} = \textbf{tribunal} ou \textbf{terrain} (tennis). \textit{Prononciation :} \eng{course} /kɔːs/ (long), \eng{court} /kɔːt/ (avec t final entendu).
    \item \textbf{Notes} (pluriel) = \textbf{notes de cours}. \eng{Note} (singulier) = \textbf{note} (musique, billet, annotation) ou \textbf{note/grade} (résultat US). \eng{Grades} (US) / \eng{Marks} (UK) = \textbf{notes (résultats scolaires)}.
    \item \textbf{Eventually} $\neq$ Éventuellement. \eng{Eventually} = \textbf{finalement, au bout du compte}. \eng{Possibly} / \eng{Potentially} = \textbf{éventuellement}.
    \item \textbf{Assist} $\neq$ Assister (à un spectacle). \eng{Assist} = \textbf{aider}. \eng{Attend} = \textbf{assister à} (un cours, une réunion). \eng{Attendee} = participant.
    \item \textbf{Download / Upload} : ne pas confondre la direction (voir encadré plus haut).
\end{itemize}
\end{attention}

\begin{popfigure}
\begin{center}\tcbox[colback=popink,colframe=popink,coltext=white,arc=9pt,boxsep=4pt,left=6pt,right=6pt]{\bfseries\small E-learning}\end{center}
\vspace{-2pt}
\begin{tcbraster}[raster columns=2,raster equal height=rows,raster column skip=6pt,raster row skip=6pt,enhanced,blankest]
\begin{tcolorbox}[enhanced,colback=white,colframe=poporange,boxrule=0.9pt,arc=3pt,title={Resources},fonttitle=\bfseries\scriptsize,coltitle=white,colbacktitle=poporange,left=4pt,right=4pt,top=2pt,bottom=2pt,boxsep=1pt,toptitle=1.5pt,bottomtitle=1.5pt,halign=left]
\begin{itemize}[nosep,leftmargin=*,itemsep=1pt,font=\scriptsize]
\item E-books
\item Past papers
\item Video tutorials
\item Online dictionaries
\item Digital libraries
\item Audio explanations
\end{itemize}
\end{tcolorbox}
\begin{tcolorbox}[enhanced,colback=white,colframe=poppink,boxrule=0.9pt,arc=3pt,title={Activities},fonttitle=\bfseries\scriptsize,coltitle=white,colbacktitle=poppink,left=4pt,right=4pt,top=2pt,bottom=2pt,boxsep=1pt,toptitle=1.5pt,bottomtitle=1.5pt,halign=left]
\begin{itemize}[nosep,leftmargin=*,itemsep=1pt,font=\scriptsize]
\item Watch tutorials
\item Download / Upload
\item Share documents
\item Join study groups
\item Do research
\item Take notes
\item Revise online
\end{itemize}
\end{tcolorbox}
\begin{tcolorbox}[enhanced,colback=white,colframe=poppurple,boxrule=0.9pt,arc=3pt,title={Actors},fonttitle=\bfseries\scriptsize,coltitle=white,colbacktitle=poppurple,left=4pt,right=4pt,top=2pt,bottom=2pt,boxsep=1pt,toptitle=1.5pt,bottomtitle=1.5pt,halign=left]
\begin{itemize}[nosep,leftmargin=*,itemsep=1pt,font=\scriptsize]
\item Online tutor
\item Learner / User
\item Content creator
\item Blogger
\item Moderator
\end{itemize}
\end{tcolorbox}
\begin{tcolorbox}[enhanced,colback=white,colframe=popgold,boxrule=0.9pt,arc=3pt,title={Benefits},fonttitle=\bfseries\scriptsize,coltitle=white,colbacktitle=popgold,left=4pt,right=4pt,top=2pt,bottom=2pt,boxsep=1pt,toptitle=1.5pt,bottomtitle=1.5pt,halign=left]
\begin{itemize}[nosep,leftmargin=*,itemsep=1pt,font=\scriptsize]
\item Learn at own pace
\item Access info easily
\item Save time
\item Improve grades
\item Develop digital skills
\end{itemize}
\end{tcolorbox}
\end{tcbraster}

\legende{Carte mentale : L'écosystème de l'e-learning (Ressources, Activités, Acteurs, Avantages)}
\end{popfigure}

\courssub{Exercice résolu : Mettre en pratique}

\begin{exoresolu}{Compléter un texte argumentatif}
\textbf{Consigne :} Complétez le paragraphe ci-dessous avec les mots/expressions de la liste. Attention à la forme verbale (infinitif, gérondif, conjugaison) et aux prépositions.

\textit{Liste :} \textbf{access information easily, download, e-learning, educational app, improve one's grades, join, online dictionary, past papers, research, upload, virtual classroom, watch video tutorials.}

\eng{During the long holidays, many students in Yaoundé and Douala turn to \_\_\_\_\_\_\_\_\_\_\_\_ to keep up with their studies. Platforms like Khan Academy or Coursera offer free \_\_\_\_\_\_\_\_\_\_\_\_. Students can also \_\_\_\_\_\_\_\_\_\_\_\_ \_\_\_\_\_\_\_\_\_\_\_\_ from official websites to practise for the Baccalaureate. When they encounter a difficult word, they use an \_\_\_\_\_\_\_\_\_\_\_\_ to check the definition and pronunciation. Some prefer to \_\_\_\_\_\_\_\_\_\_\_\_ a WhatsApp \_\_\_\_\_\_\_\_\_\_\_\_ to discuss difficult topics with peers. Teachers often \_\_\_\_\_\_\_\_\_\_\_\_ corrected exercises to a shared drive so that students can \_\_\_\_\_\_\_\_\_\_\_\_ them. The main advantage is that learners can \_\_\_\_\_\_\_\_\_\_\_\_ and \_\_\_\_\_\_\_\_\_\_\_\_ at their own pace, which helps them \_\_\_\_\_\_\_\_\_\_\_\_.}

\tcblower
\textbf{Solution détaillée :}
\begin{enumerate}
    \item \textbf{e-learning} (nom, après \eng{turn to} $\rightarrow$ nom/gérondif. \eng{Turn to e-learning} fonctionne comme nom non comptable désignant le domaine).
    \item \textbf{watch video tutorials} (après \eng{offer free} $\rightarrow$ gérondif ou groupe nominal. \eng{Watch video tutorials} est une activité).
    \item \textbf{download past papers} (verbe \eng{download} + objet \eng{past papers}). \textit{Note :} \eng{download} à l'infinitif après \eng{can} (modal) $\rightarrow$ \eng{can download}.
    \item \textbf{online dictionary} (nom singulier après \eng{an}).
    \item \textbf{join a study group} (verbe \eng{join} + objet. Après \eng{prefer to} $\rightarrow$ infinitif \eng{to join}).
    \item \textbf{upload} (verbe à la 3e personne du pluriel \eng{Teachers often upload}).
    \item \textbf{download} (verbe à l'infinitif après \eng{can} $\rightarrow$ \eng{can download}).
    \item \textbf{access information easily} (infinitif de but/avantage après \eng{can}).
    \item \textbf{learn at their own pace} (infinitif après \eng{can} ; \eng{their} car \eng{learners} pluriel).
    \item \textbf{improve their grades} (infinitif après \eng{helps them} ; \eng{help someone (to) do}).
\end{enumerate}

\textbf{Texte reconstitué :}
\eng{During the long holidays, many students in Yaoundé and Douala turn to \textbf{e-learning} to keep up with their studies. Platforms like Khan Academy or Coursera offer free \textbf{video tutorials}. Students can also \textbf{download past papers} from official websites to practise for the Baccalaureate. When they encounter a difficult word, they use an \textbf{online dictionary} to check the definition and pronunciation. Some prefer to \textbf{join a study group} on WhatsApp to discuss difficult topics with peers. Teachers often \textbf{upload} corrected exercises to a shared drive so that students can \textbf{download} them. The main advantage is that learners can \textbf{access information easily} and \textbf{learn at their own pace}, which helps them \textbf{improve their grades}.}
\end{exoresolu}

\courssub{Entraînement personnel (Exercices)}

\begin{exercice}{Exercice 1 : Le bon verbe d'action}
Choisissez le verbe approprié parmi la liste (à conjuguer si nécessaire) pour compléter chaque phrase.
\textit{Liste : download, upload, watch, revise, join, share, do, take, sign up for, log in.}
\begin{enumerate}
    \item Before the exam, I always \_\_\_\_\_\_\_\_\_\_\_\_ my notes on my tablet.
    \item The teacher asked us to \_\_\_\_\_\_\_\_\_\_\_\_ our homework on the platform by Friday.
    \item If you don't understand the lesson, \_\_\_\_\_\_\_\_\_\_\_\_ a video tutorial on YouTube.
    \item We need to \_\_\_\_\_\_\_\_\_\_\_\_ research on the Internet for our geography project.
    \item Don't forget to \_\_\_\_\_\_\_\_\_\_\_\_ notes during the webinar.
    \item My friend \_\_\_\_\_\_\_\_\_\_\_\_ for an online course last week to learn coding.
    \item Please \_\_\_\_\_\_\_\_\_\_\_\_ to the platform using your student ID.
\end{enumerate}
\end{exercice}

\begin{exercice}{Exercice 2 : Faux-amis et indénombrables}
Corrigez les erreurs dans les phrases suivantes (vocabulaire, nombre, préposition).
\begin{enumerate}
    \item I made many researches on the Internet for my essay.
    \item She went to the librairie to borrow an e-book.
    \item He assisted the webinar yesterday afternoon.
    \item The course is very interesting, I will eventually pass the exam.
    \item We need to download the files on the server.
    \item She has good notes in mathematics (average 16/20).
\end{enumerate}
\end{exercice}

\begin{exercice}{Exercice 3 : Production écrite guidée (Writing)}
\textbf{Sujet :} \eng{"How do ICTs help you learn better? Write a short paragraph (80-100 words) explaining your personal experience. Use at least 8 words/expressions from this lesson (Lexical Field 2). Underline them."}
\textbf{Indications :} Structurez votre paragraphe : 1. Outil principal utilisé. 2. Ressources consultées. 3. Actions effectuées. 4. Avantages ressentis. Utilisez les connecteurs logiques (\eng{Firstly, Moreover, Consequently, Finally}).
\end{exercice}

\begin{corrige}{Exercice 1}
\begin{enumerate}
    \item \textbf{revise} (ou \textbf{review} US) — \eng{revise my notes}.
    \item \textbf{upload} — \eng{upload our homework} (mettre en ligne vers le prof).
    \item \textbf{watch a video tutorial} (ou \textbf{watch video tutorials}).
    \item \textbf{do research} (indénombrable ! pas \eng{researches}).
    \item \textbf{take notes}.
    \item \textbf{signed up for} (phrasal verb, prétérit).
    \item \textbf{log in}.
\end{enumerate}
\end{corrige}

\begin{corrige}{Exercice 2}
\begin{enumerate}
    \item \eng{I \textbf{did some research} (ou \textbf{carried out research}) on the Internet for my essay.} (\eng{research} indénombrable, \eng{do research}).
    \item \eng{She went to the \textbf{digital library} (ou \textbf{online library}) to borrow an e-book.} (\eng{Library} = bibliothèque, pas librairie).
    \item \eng{He \textbf{attended} the webinar yesterday afternoon.} (\eng{Attend} = assister à un événement ; \eng{assist} = aider).
    \item \eng{The course is very interesting, I \textbf{might} (ou \textbf{possibly}) pass the exam.} (\eng{Eventually} = finalement ; pour « éventuellement » = \eng{possibly} / \eng{might}).
    \item \eng{We need to \textbf{upload} the files \textbf{to} the server.} (\eng{Upload} = vers le serveur ; préposition \eng{to} après verbes de transfert).
    \item \eng{She has good \textbf{grades} (US) / \textbf{marks} (UK) in mathematics.} (\eng{Notes} = notes de cours ; \eng{Grades/Marks} = résultats scolaires).
\end{enumerate}
\end{corrige}

\begin{savaistu}{Le saviez-vous ? L'origine du mot "Webinar"}
Le mot \textbf{webinar} est un \textbf{mot-valise} (portmanteau word), contraction de \eng{Web} (toile, World Wide Web) + \eng{Seminar} (séminaire, du latin \emph{seminarium} « pépinière », de \emph{semen} « graine » $\rightarrow$ lieu où l'on sème le savoir). Il est apparu dans les années 1990 avec les premières conférences audio/vidéo sur Internet (ex. \eng{NetMeeting}). Aujourd'hui, des plateformes comme \eng{Zoom}, \eng{Teams}, \eng{Google Meet} ou \eng{Webex} en font un outil standard de la formation continue et de l'enseignement hybride. Au Cameroun, l'Université de Yaoundé I et l'Université de Buea ont massivement adopté le \eng{webinar} pendant la crise du Covid-19 pour assurer la continuité pédagogique.
\end{savaistu}

\begin{methode}{Comment réviser ce vocabulaire efficacement ? (Méthode "3R")}
\begin{enumerate}
    \item \textbf{Record (Enregistrer) :} Créez une \eng{mindmap} (carte mentale) sur une feuille A4 ou via une \eng{app} (ex. \eng{MindMeister}, \eng{XMind}) avec 4 branches : \eng{Modalities}, \eng{Resources}, \eng{Actions}, \eng{Actors/Benefits}. Écrivez \textbf{uniquement en anglais}.
    \item \textbf{Recycle (Recycler) :} Écrivez 5 phrases vraies sur VOS habitudes d'étude numérique en utilisant \textbf{au moins 3 mots par phrase} de la leçon. Ex. \eng{I usually download past papers and watch video tutorials to revise online.}
    \item \textbf{Retrieve (Récupérer) :} Dans 48h, sans regarder le cours, refaites l'Exercice 1 à l'oral (vous vous enregistrez sur votre téléphone) et corrigez-vous avec le corrigé. La récupération active (active recall) fixe la mémoire à long terme.
\end{enumerate}
\end{methode}

\coursec{Lexical field 3: ICTs for work}

Dans la section précédente, nous avons vu comment les TIC transforment l'apprentissage : \eng{e-learning}, \eng{online courses}, \eng{to do research online}. Mais l'école n'est qu'un début. Une fois diplômés, la plupart d'entre vous entreront dans un monde du travail où l'ordinateur, l'Internet et le téléphone portable sont devenus aussi indispensables que le stylo l'était autrefois. À Douala comme à Yaoundé, on paie ses factures par \eng{mobile money}, on envoie son CV par \eng{email} et on participe à des réunions sans quitter son salon. Cette section vous donne tout le vocabulaire nécessaire pour parler du \cle{travail à l'ère numérique} — un thème très fréquent à l'oral et à l'écrit du Baccalauréat.

\courssub{Discovering the vocabulary in context}

Avant toute liste de mots, lisons un texte authentique. Observez les mots en gras : ils forment le cœur du champ lexical de cette section.

\begin{textebox}[A Day in the Life of a Digital Worker — original text]
Mrs Etoundi runs a small clothing business in Bafoussam. Five years ago, she spent her days in a crowded shop; today, she manages most of her work from her laptop at home. \eng{“Working from home has changed my life,”} she says. \eng{“Twice a week, I attend video conferences with my suppliers in Douala, so I no longer waste hours in traffic jams.”}

Her morning routine is simple. First, she checks her \eng{emails} and replies to customers. She opens the \eng{attachments} they send — photos of new fabrics, price lists — and saves each \eng{file} in the right \eng{folder} on her computer. Then she \eng{backs up} all her data on a \eng{cloud storage} service, because she once lost everything when her old computer \eng{crashed}. \eng{“Now I can open my documents anywhere, even on my phone,”} she explains.

Most of her income comes from \eng{e-commerce}: customers order dresses on her website and pay by \eng{mobile money}. She also uses \eng{online banking} to pay her employees and \eng{digital marketing} — posts on social media — to attract new clients. When customers have a problem, her assistant provides \eng{customer service online}, answering questions within minutes.

Of course, technology is not perfect. Last month, a \eng{virus} infected her laptop and a \eng{bug} blocked her website for two days. Her Internet \eng{connection} is sometimes so \eng{slow} that she cannot join \eng{online meetings}, and when the network is completely \eng{offline}, she simply restarts her router and waits. \eng{“You need real computer skills to survive in business today,”} she laughs. \eng{“That is why I tell my children: become computer-literate, or you will be left behind.”}
\end{textebox}

\textbf{Glossaire.}

\begin{center}
\begin{tabularx}{\linewidth}{|l|l|X|}
\hline
\rowcolor{popblueL} English & Nature & Français \\ \hline
supplier & noun & fournisseur \\ \hline
traffic jam & noun & embouteillage \\ \hline
fabric & noun & tissu \\ \hline
income & noun & revenu \\ \hline
router & noun & routeur (boîtier Internet) \\ \hline
to be left behind & expression & être laissé de côté, dépassé \\ \hline
\end{tabularx}
\end{center}

\textbf{Comprehension questions.}
\begin{enumerate}
\item Where does Mrs Etoundi work, and how has her life changed?
\item Why does she back up her data on cloud storage?
\item How do her customers pay for their orders?
\item What technical problems did she face last month?
\item What advice does she give to her children?
\end{enumerate}

\courssub{Teleworking and the modern workplace}

\begin{definition}[Teleworking / working from home / remote work]
\eng{Teleworking} (or \eng{telecommuting}) means working from home or from another place outside the office, using a computer and the Internet. \eng{Remote work} is a synonym: \eng{a remote worker} = un travailleur à distance. In everyday English, people simply say \eng{to work from home}.
\end{definition}

\begin{exemplebox}[Observation]
\eng{My father works from home twice a week and attends video conferences.}\\
« Mon père travaille à domicile deux fois par semaine et participe à des visioconférences. »

\eng{Since 2020, remote work has become common in many companies.}\\
« Depuis 2020, le télétravail est devenu courant dans de nombreuses entreprises. »

\eng{Our manager schedules an online meeting every Monday at 9 a.m.}\\
« Notre chef programme une réunion en ligne chaque lundi à 9 h. »
\end{exemplebox}

Notez bien les verbes associés aux réunions : on \cle{attend} a meeting (y assister), on \cle{schedules} a meeting (la programmer), on \cle{joins} a meeting (la rejoindre en ligne), on \cle{cancels} a meeting (l'annuler).

\courssub{Emails, files and productivity tools}

Le courrier électronique est la colonne vertébrale du travail moderne. Voici le vocabulaire indispensable, avec les verbes qui l'accompagnent — ce sont ces \cle{collocations} qui font la différence entre une copie moyenne et une excellente copie.

\begin{center}
\begin{tabularx}{\linewidth}{|X|X|X|}
\hline
\rowcolor{popblueL} English & Français & Exemple \\ \hline
an email & un courriel & \eng{to send / to receive an email} \\ \hline
an attachment & une pièce jointe & \eng{to open / to send an attachment} \\ \hline
a file & un fichier & \eng{to save / to delete a file} \\ \hline
a folder & un dossier & \eng{to create a new folder} \\ \hline
to reply (to) & répondre (à) & \eng{Please reply to my email.} \\ \hline
to forward & transférer & \eng{Can you forward me the message?} \\ \hline
to print & imprimer & \eng{I need to print the report.} \\ \hline
to scan & numériser, scanner & \eng{She scanned her certificates.} \\ \hline
to back up data & sauvegarder des données & \eng{Back up your work every day.} \\ \hline
cloud storage & stockage en ligne & \eng{My photos are on cloud storage.} \\ \hline
\end{tabularx}
\end{center}

\begin{definition}[Cloud storage]
\eng{Cloud storage} means saving files on the Internet (not on your device) so you can open them anywhere, on any computer or phone — for example with Google Drive, Dropbox or OneDrive. « Le stockage en nuage » : vos fichiers vivent sur Internet, pas sur votre appareil.
\end{definition}

\begin{exemplebox}[Exemples traduits]
\eng{Don't forget to back up your files before the exam!}\\
« N'oublie pas de sauvegarder tes fichiers avant l'examen ! »

\eng{I saved the document in a folder called “Homework” and forwarded it to my teacher as an attachment.}\\
« J'ai enregistré le document dans un dossier appelé “Devoirs” et je l'ai transféré à mon professeur en pièce jointe. »

\eng{Could you scan this form and email it to the bank?}\\
« Pourrais-tu scanner ce formulaire et l'envoyer par courriel à la banque ? »
\end{exemplebox}

\begin{attention}[To reply, pas “to response”]
Le nom est \eng{a reply} ou \eng{a response}, mais le verbe est uniquement \eng{to reply (to)}. On écrit \eng{She replied to my email}, jamais \eng{She responsed my email}. De même, \eng{to forward} prend un complément direct : \eng{Forward the file to me} (« Transfère-moi le fichier »).
\end{attention}

\courssub{Digital business: e-commerce, mobile money and online services}

Le Cameroun est un terrain parfait pour observer ce vocabulaire : qui n'a jamais payé par \eng{mobile money} ou suivi une boutique sur les réseaux sociaux ?

\begin{definition}[Key business terms]
\begin{itemize}
\item \eng{E-commerce} (electronic commerce): buying and selling goods or services on the Internet. « Le commerce en ligne. »
\item \eng{Mobile money}: a service that lets you send, receive and keep money using a mobile phone, without a bank account. « L'argent mobile » (Orange Money, MTN Mobile Money).
\item \eng{Online banking}: managing your bank account on the Internet — checking your balance, transferring money, paying bills. « La banque en ligne. »
\item \eng{Digital marketing}: advertising products on the Internet and social media. « Le marketing numérique. »
\item \eng{Customer service online}: helping customers by email, chat or social media. « Le service client en ligne. »
\end{itemize}
\end{definition}

\begin{exemplebox}[Exemples traduits]
\eng{Many shops in Douala now accept mobile money payments.}\\
« Beaucoup de boutiques à Douala acceptent désormais les paiements par mobile money. »

\eng{With online banking, I can pay my electricity bill without leaving my house.}\\
« Avec la banque en ligne, je peux payer ma facture d'électricité sans quitter la maison. »

\eng{Their digital marketing campaign attracted thousands of new customers.}\\
« Leur campagne de marketing numérique a attiré des milliers de nouveaux clients. »
\end{exemplebox}

\courssub{Digital skills and technical problems}

\begin{definition}[Talking about skills]
\eng{Computer skills} / \eng{IT skills} (IT = Information Technology) are the abilities you need to use a computer well. \eng{Digital literacy} is the general ability to find, understand and use information online. Someone who has these skills \eng{is computer-literate} (« maîtrise l'outil informatique ») ; the opposite is \eng{computer-illiterate}.
\end{definition}

\begin{exemplebox}[Exemples traduits]
\eng{Good IT skills are essential for this job.}\\
« De bonnes compétences informatiques sont indispensables pour ce poste. »

\eng{Young people today are generally more computer-literate than their parents.}\\
« Les jeunes d'aujourd'hui maîtrisent généralement mieux l'informatique que leurs parents. »
\end{exemplebox}

Quand la technologie tombe en panne, il faut aussi savoir en parler :

\begin{center}
\begin{tabularx}{\linewidth}{|X|X|X|}
\hline
\rowcolor{popblueL} Problem & Français & Exemple \\ \hline
a breakdown & une panne & \eng{There was a network breakdown.} \\ \hline
a bug & un bogue, un bug & \eng{The app has a bug.} \\ \hline
a virus & un virus & \eng{A virus infected my laptop.} \\ \hline
to crash & planter & \eng{My computer crashed during the meeting.} \\ \hline
to be offline & être hors ligne & \eng{The server is offline.} \\ \hline
a slow connection & une connexion lente & \eng{The connection is too slow for video calls.} \\ \hline
to reset / to restart & réinitialiser / redémarrer & \eng{Try restarting your router.} \\ \hline
\end{tabularx}
\end{center}

\begin{popfigure}
\begin{tikzpicture}[node distance=0.9cm and 1.4cm,
box/.style={draw=popblue, fill=popblueL, rounded corners, align=center, minimum height=1cm, inner sep=5pt, font=\small},
arr/.style={-{Stealth[length=3mm]}, thick, poporange}]
\node[box, align=center] (create){create\\ a document};
\node[box, right=of create, align=center] (save){save it\\ in a folder};
\node[box, right=of save, align=center] (backup){back it up\\ (cloud storage)};
\node[box, below=of backup, align=center] (share){share it\\ by email (attachment)};
\node[box, left=of share, align=center] (feedback){receive\\ feedback};
\draw[arr] (create) -- (save);
\draw[arr] (save) -- (backup);
\draw[arr] (backup) -- (share);
\draw[arr] (share) -- (feedback);
\draw[arr, dashed, popgreen] (feedback.west) to[bend right=25] node[left, font=\small\itshape, popgreen]{revise} (create.west);
\end{tikzpicture}
\legende{The digital work cycle: create, save, back up, share, receive feedback — then revise and start again.}
\end{popfigure}

\courssub{False friends and common mistakes}

\begin{attention}[To assist ≠ assister à]
\eng{To assist} means “to help” (\eng{The nurse assisted the doctor} = « L'infirmière a aidé le médecin »). It does NOT mean « assister à ». For « assister à une réunion », say \eng{to attend a meeting}. This is a classic false friend at the Bac: \eng{I attended the video conference} (« J'ai assisté à la visioconférence »), never \eng{I assisted the meeting}. Similarly, \eng{an assistant} is « un(e) assistant(e) », someone who helps — not someone who attends.
\end{attention}

\begin{attention}[Autres pièges fréquents]
\begin{itemize}
\item \eng{Actually} signifie « en fait », pas « actuellement » : \eng{I currently work from home} (« Je travaille actuellement à domicile »).
\item \eng{A library} est une bibliothèque ; une librairie se dit \eng{a bookshop}.
\item On dit \eng{to do research online}, jamais \eng{to make researches}.
\item \eng{Data} est souvent suivi d'un verbe au singulier en anglais courant : \eng{My data is backed up}.
\end{itemize}
\end{attention}

\begin{methode}[Enrichir une copie : des mots simples aux collocations précises]
Pour obtenir une excellente note en expression écrite, remplacez les expressions vagues par des collocations exactes :
\begin{enumerate}
\item Repérez les expressions trop simples : \eng{use the computer}, \eng{do things on the Internet}.
\item Remplacez-les par des verbes précis : \eng{browse the web}, \eng{send attachments}, \eng{back up data}, \eng{schedule online meetings}, \eng{attend video conferences}.
\item Ajoutez un problème technique pour nuancer : \eng{although the connection was slow}, \eng{when the system crashed}.
\item Exemple de transformation : \eng{Workers use computers a lot.} → \eng{Remote workers attend video conferences, exchange attachments and back up their data on cloud storage.}
\end{enumerate}
\end{methode}

\begin{exoresolu}[Guided practice]
\textbf{Énoncé.} Complete the sentences with the correct word or expression from the lesson: \eng{attachment, back up, attend, crashed, mobile money, computer-literate, offline, schedule}.
\begin{enumerate}
\item My laptop \dots\ during the video conference, so I missed the end of the meeting.
\item Please \dots\ your files on cloud storage before you format the computer.
\item I cannot open the \dots\ you sent me; could you forward it again?
\item Many customers in Yaoundé prefer to pay by \dots.
\item The manager will \dots\ an online meeting for Friday at 10 a.m.
\item When the network is \dots, I cannot check my emails.
\item To get this job, you must be fully \dots.
\item Did you \dots\ the training session yesterday?
\end{enumerate}
\tcblower
\textbf{Solution détaillée.}
\begin{enumerate}
\item \eng{crashed} — « a planté » ; passé simple car l'action est terminée.
\item \eng{back up} — « sauvegarde » ; impératif avec \eng{please}.
\item \eng{attachment} — « la pièce jointe » ; on \eng{opens} ou \eng{forwards} une pièce jointe.
\item \eng{mobile money} — « l'argent mobile » ; collocation \eng{to pay by mobile money}.
\item \eng{schedule} — « programmera » ; futur avec \eng{will} + base verbale.
\item \eng{offline} — « hors ligne » ; attention, pas d'article : \eng{to be offline}.
\item \eng{computer-literate} — « maîtriser l'informatique » ; adjectif composé avec trait d'union.
\item \eng{attend} — « as-tu assisté à » ; jamais \eng{assist}, faux ami classique.
\end{enumerate}
\end{exoresolu}

\begin{exercice}[Your turn]
Write a short paragraph (80–100 words) entitled \eng{“How ICTs help my family at work”}. Use at least six words or expressions from this section, including one collocation with \eng{back up}, \eng{attend} or \eng{forward}, and one technical problem (\eng{crash, bug, slow connection, offline}). Underline the vocabulary items you use.
\end{exercice}

\begin{corrige}[Exercice : production écrite]
Modèle de réponse possible : \eng{My uncle runs a cybercafé in Bamenda, and ICTs are at the heart of his business. Every morning, he checks his emails and replies to customers who want to print or scan documents. He always backs up his data on cloud storage, because his old computer once crashed and he lost everything. His customers often pay by mobile money, which is fast and safe. He also attends online meetings with his suppliers. Of course, the connection is sometimes slow, and when the network goes offline, he simply restarts the router. Thanks to his computer skills, his business is growing every year.} — Vérifiez que votre paragraphe contient au moins six éléments du champ lexical, un problème technique et aucune occurrence du faux ami \eng{assist} pour « assister à ».
\end{corrige}

\begin{aretenir}[What to remember]
\begin{itemize}
\item Travail à distance : \eng{to work from home, remote work, to attend a video conference, to schedule an online meeting}.
\item Fichiers et emails : \eng{an attachment, a file, a folder, to save, to back up, to forward, to reply to, to print, to scan}.
\item \eng{Cloud storage} = sauvegarder ses fichiers sur Internet pour y accéder partout.
\item Commerce numérique : \eng{e-commerce, mobile money, online banking, digital marketing, customer service online}.
\item Compétences : \eng{computer skills, IT skills, digital literacy, to be computer-literate}.
\item Pannes : \eng{a breakdown, a bug, a virus, to crash, to be offline, a slow connection, to restart}.
\item Piège du Bac : \eng{to assist} = aider ; « assister à » = \eng{to attend}.
\end{itemize}
\end{aretenir}

\coursec{Collocations, word families and false friends}

Dans la section précédente, nous avons exploré le vocabulaire des TIC au travail : \eng{to attend a video conference}, \eng{to send an attachment}, \eng{to work remotely}. Vous avez remarqué que certains mots anglais « vont toujours ensemble » : on dit \eng{to attend a meeting} et non \eng{to assist a meeting}. C'est exactement l'objet de cette section : apprendre non pas des mots isolés, mais des \cle{combinaisons de mots} (les \emph{collocations}), des \cle{familles de mots} (comment un même radical produit un verbe, un nom, un adjectif) et surtout comment éviter les \cle{faux amis} qui coûtent cher aux examens.

\courssub{Observing the language: a short text}

Lisons d'abord ce court article original. Soulignez mentalement les groupes de mots qui se répètent autour des TIC.

\begin{textebox}[How ICTs changed my studies — an original article]
When I entered the Lower Sixth at a bilingual high school in Yaoundé, my father bought me a second-hand laptop and a smartphone. At first, I only used my phone to browse the web and chat with friends. Then our computer science teacher showed us how to use ICTs to enhance our learning. Now, every morning, I log in to the school platform with my password, I download the files our teachers have posted, and I watch recorded lessons when the connection is reliable. Before examinations, I surf the Internet to find past papers, and I post messages in our revision group on WhatsApp.

My phone is also my working tool. I always charge it at night so that the battery lasts the whole day. I have created an account on an e-learning site, and I update my apps regularly because old versions are slow. My favourite app is very user-friendly: even my grandmother can use it! She is not a digital native like me, but she is learning fast. Thanks to high-speed internet at the cybercafe near our house, she can even attend video calls with my uncle in Buea.

However, I have learnt to be careful. Not everything online is informative; some websites communicate false information. My teacher says a powerful computer and a good connectivity are useful only if the user is well informed. Computing is a skill, and like every skill, it must be practised every day.
\end{textebox}

\textbf{Glossaire.}

\begin{center}
\begin{tabularx}{\linewidth}{|l|l|X|}
\hline
\rowcolor{popblueL} \textbf{Mot anglais} & \textbf{Nature} & \textbf{Traduction française} \\
\hline
second-hand & adjectif & d'occasion \\
\hline
to enhance & verbe & améliorer, renforcer \\
\hline
past papers & nom pluriel & anciens sujets d'examen \\
\hline
user-friendly & adjectif & convivial, facile à utiliser \\
\hline
high-speed internet & nom & Internet à haut débit \\
\hline
connectivity & nom & connectivité (qualité de la connexion) \\
\hline
computing & nom & informatique (l'activité, la discipline) \\
\hline
\end{tabularx}
\end{center}

\textbf{Questions de compréhension.}

\begin{enumerate}
\item What does the narrator do every morning on the school platform?
\item Why does the narrator update his apps regularly?
\item Who is learning to use a smartphone in the narrator's family?
\item According to the teacher, what is necessary besides a powerful computer?
\end{enumerate}

\begin{corrige}[Questions de compréhension]
\begin{enumerate}
\item He logs in with his password and downloads the files the teachers have posted.
\item Because old versions are slow.
\item His grandmother (she is not a digital native, but she is learning fast).
\item The user must be well informed; computing is a skill that must be practised every day.
\end{enumerate}
\end{corrige}

\courssub{Collocations: verb + noun}

\begin{definition}[Collocation (combinaison de mots)]
Une \cle{collocation} est une combinaison habituelle de deux mots ou plus qui « sonnent juste » ensemble pour un locuteur natif. En anglais, on dit \eng{to browse the web} et non \eng{to navigate the web} dans le langage courant, de même qu'en français on dit « prendre une décision » et non « faire une décision ». Les collocations s'apprennent comme des blocs, jamais mot à mot.
\end{definition}

Voici les collocations \cle{verbe + nom} essentielles du domaine des TIC, toutes présentes dans le texte.

\begin{center}
\begin{tabularx}{\linewidth}{|l|X|X|}
\hline
\rowcolor{popblueL} \textbf{Collocation anglaise} & \textbf{Traduction} & \textbf{Exemple en contexte} \\
\hline
to browse the web & naviguer sur le web & \eng{I browse the web to find past papers.} \\
\hline
to surf the Internet & surfer sur Internet & \eng{Teenagers surf the Internet for hours.} \\
\hline
to download a file & télécharger un fichier & \eng{She downloaded the file before class.} \\
\hline
to log in / to log out & se connecter / se déconnecter & \eng{Log out when you finish your work.} \\
\hline
to create an account & créer un compte & \eng{He created an account on an e-learning site.} \\
\hline
to post a message & publier un message & \eng{They post messages in the revision group.} \\
\hline
to update an app & mettre à jour une application & \eng{Update your apps regularly.} \\
\hline
to charge your phone & charger ton téléphone & \eng{I charge my phone every night.} \\
\hline
\end{tabularx}
\end{center}

\begin{exemplebox}[Exemples traduits]
\eng{You must log in before you can download the file.} « Tu dois te connecter avant de pouvoir télécharger le fichier. »

\eng{Don't forget to log out when you use a public computer at the cybercafe.} « N'oublie pas de te déconnecter quand tu utilises un ordinateur public au cybercafé. »

\eng{My sister posted a message in our WhatsApp group to announce the good news.} « Ma sœur a publié un message dans notre groupe WhatsApp pour annoncer la bonne nouvelle. »
\end{exemplebox}

\begin{attention}[To download : un seul mot en anglais]
Le français distingue « télécharger » (dans les deux sens) ; l'anglais distingue \eng{to download} (télécharger depuis Internet vers ton appareil) et \eng{to upload} (envoyer un fichier de ton appareil vers Internet) : \eng{Students upload their homework to the platform.} « Les élèves envoient leurs devoirs sur la plateforme. » Attention aussi à \eng{to charge} : il signifie « charger (une batterie) » mais aussi « faire payer » : \eng{The cybercafe charges 200 francs per hour.} Ne confondez pas les deux sens.
\end{attention}

\courssub{Collocations: adjective + noun}

Les adjectifs aussi ont leurs partenaires privilégiés. Un ordinateur n'est pas \eng{strong} mais \eng{powerful} ; une connexion n'est pas \eng{sure} mais \eng{reliable}.

\begin{center}
\begin{tabularx}{\linewidth}{|l|X|X|}
\hline
\rowcolor{popblueL} \textbf{Collocation anglaise} & \textbf{Traduction} & \textbf{Exemple en contexte} \\
\hline
a reliable connection & une connexion fiable & \eng{We need a reliable connection for online exams.} \\
\hline
a powerful computer & un ordinateur puissant & \eng{A powerful computer processes data quickly.} \\
\hline
a user-friendly app & une application conviviale & \eng{This dictionary app is very user-friendly.} \\
\hline
a digital native & un natif du numérique & \eng{Most teenagers today are digital natives.} \\
\hline
high-speed internet & l'Internet à haut débit & \eng{High-speed internet is available in Douala.} \\
\hline
\end{tabularx}
\end{center}

\begin{exemplebox}[Exemple traduit]
\eng{This app is very user-friendly.} « Cette application est très facile à utiliser. »

\eng{With high-speed internet, you can download a film in a few minutes.} « Avec l'Internet à haut débit, tu peux télécharger un film en quelques minutes. »
\end{exemplebox}

\begin{attention}[Les calques de « fort » et « rapide »]
Ne traduisez jamais mot à mot : un ordinateur « fort » n'est pas \eng{a strong computer} mais \eng{a powerful computer} ; un Internet « rapide » n'est pas \eng{a fast internet} dans l'expression consacrée mais \eng{high-speed internet}. De même, « une application facile à utiliser » se dit élégamment \eng{a user-friendly app} — notez le trait d'union obligatoire dans cet adjectif composé placé devant le nom.
\end{attention}

\courssub{Word families: one root, many words}

L'anglais construit des familles entières à partir d'un même radical. Connaître ces familles multiplie votre vocabulaire sans effort et vous aide à choisir la bonne \cle{nature grammaticale} dans une phrase.

\begin{center}
\begin{tabular}{|l|l|l|l|}
\hline
\rowcolor{popblueL} \textbf{Verbe} & \textbf{Nom} & \textbf{Adjectif} & \textbf{Autre forme} \\
\hline
to compute (calculer) & a computer (ordinateur) & computerized (informatisé) & computing (l'informatique) \\
\hline
to connect (connecter) & a connection (connexion) & connected (connecté) & connectivity (connectivité) \\
\hline
to inform (informer) & information (information) & informative (instructif) & an informant (informateur) \\
\hline
to communicate (communiquer) & communication (communication) & communicative (communicatif) & telecommunications (télécommunications) \\
\hline
\end{tabular}
\end{center}

\begin{exemplebox}[La même famille dans quatre phrases]
\eng{Computers compute millions of operations per second.} « Les ordinateurs calculent des millions d'opérations par seconde. » (verbe)

\eng{Our school has a computerized library.} « Notre école a une bibliothèque informatisée. » (adjectif)

\eng{She studies computing at the university.} « Elle étudie l'informatique à l'université. » (nom d'activité, invariable)

\eng{This computer is very powerful.} « Cet ordinateur est très puissant. » (nom d'objet)
\end{exemplebox}

\begin{attention}[Information : toujours singulier]
\eng{Information} est \cle{indénombrable} en anglais : jamais de \emph{s}, jamais de \emph{a}. On dit \eng{some information}, \eng{a piece of information} « une information ». Erreur classique : \emph{I need informations} est faux ; dites \eng{I need some information.} Même piège avec \eng{news} (singulier : \eng{The news is good.}) et \eng{equipment}.
\end{attention}

La figure suivante présente la famille du verbe \eng{to connect} sous forme d'arbre lexical : observez comment les préfixes et suffixes transforment le sens et la nature du mot.

\begin{popfigure}
\begin{center}
\begin{tikzpicture}[
  grow=right,
  level distance=3.4cm,
  level 1/.style={sibling distance=2.6cm},
  every node/.style={font=\small},
  edge from parent/.style={draw=popblue, thick, -{Stealth[length=2mm]}}
]
\node[draw=popdark, fill=popblueL, rounded corners, thick, align=center] {\textbf{connect}\\ \emph{verbe : connecter}}
  child { node[draw=poppurple, fill=poppurpleL, rounded corners, align=center] {\textbf{disconnect}\\ \emph{verbe : déconnecter}\\ préfixe négatif \cle{dis-}} }
  child { node[draw=popgreen, fill=popgreenL, rounded corners, align=center] {\textbf{connectivity}\\ \emph{nom abstrait : connectivité}\\ suffixe \cle{-ivity}} }
  child { node[draw=poporange, fill=poporangeL, rounded corners, align=center] {\textbf{connected}\\ \emph{adjectif : connecté}\\ suffixe \cle{-ed}} }
  child { node[draw=poppink, fill=poppinkL, rounded corners, align=center] {\textbf{connection}\\ \emph{nom : connexion}\\ suffixe \cle{-ion}} };
\end{tikzpicture}
\end{center}
\legende{Arbre lexical de la racine \eng{connect} : préfixes et suffixes créent de nouveaux mots.}
\end{popfigure}

\courssub{Word formation: prefixes e- and cyber-, suffixes -er / -or}

\begin{propriete}[Les préfixes e- et cyber-]
De nombreux mots des TIC se forment avec deux préfixes :

\begin{itemize}
\item \cle{e-} = \emph{electronic} (électronique) : \eng{e-mail} « courriel », \eng{e-learning} « apprentissage en ligne », \eng{e-commerce} « commerce électronique », \eng{e-book} « livre numérique ». Notez le \cle{trait d'union} après \emph{e-}.
\item \cle{cyber-} = related to computers and the Internet : \eng{cybercafe} « cybercafé », \eng{cybersecurity} « cybersécurité », \eng{cyberspace} « cyberespace », \eng{cybercrime} « cybercriminalité ». Ce préfixe se soude généralement au mot, sans trait d'union.
\end{itemize}
\end{propriete}

\begin{propriete}[Les suffixes -er et -or pour les appareils et les métiers]
Beaucoup d'appareils et de métiers anglais se forment en ajoutant \cle{-er} ou \cle{-or} à un verbe :

\begin{itemize}
\item \eng{to print} → \eng{a printer} « une imprimante » ; \eng{to scan} → \eng{a scanner} « un scanner » ; \eng{to compute} → \eng{a computer} « un ordinateur ».
\item \eng{to project} → \eng{a projector} « un projecteur » ; \eng{to operate} → \eng{an operator} « un opérateur ».
\end{itemize}

Il n'existe pas de règle absolue pour choisir entre \emph{-er} et \emph{-or} : il faut apprendre chaque mot. Retenez : \eng{printer}, \eng{scanner}, \eng{computer} avec \emph{-er} ; \eng{projector}, \eng{operator} avec \emph{-or}.
\end{propriete}

\begin{exemplebox}[Exemples traduits]
\eng{The school bought a new printer and a scanner for the computer room.} « L'école a acheté une nouvelle imprimante et un scanner pour la salle informatique. »

\eng{Cybersecurity is a growing field with many job opportunities.} « La cybersécurité est un domaine en pleine croissance avec de nombreuses offres d'emploi. »

\eng{E-learning helped many students during the school closures.} « L'apprentissage en ligne a aidé de nombreux élèves pendant les fermetures d'écoles. »
\end{exemplebox}

\courssub{False friends: beware of traps!}

Les \cle{faux amis} sont des mots qui ressemblent à des mots français mais ont un sens différent. Ce sont les pièges préférés des correcteurs du Baccalauréat.

\begin{center}
\begin{tabularx}{\linewidth}{|X|X|X|}
\hline
\rowcolor{popblueL} \textbf{Mot anglais} & \textbf{Sens réel} & \textbf{Piège français} \\
\hline
digital & numérique & « digital » au sens d'empreinte = \eng{fingerprint} (empreinte digitale) \\
\hline
to assist (somebody) & aider quelqu'un & « assister à » = \eng{to attend} (assister à une réunion) \\
\hline
a library & une bibliothèque & « une librairie » = \eng{a bookshop} \\
\hline
actually & en fait, en réalité & « actuellement » = \eng{currently}, \eng{at present} \\
\hline
\end{tabularx}
\end{center}

\begin{attention}[Actually ne veut pas dire « actuellement »]
\eng{Actually} signifie \cle{en fait, en réalité} : \eng{I don't actually have a laptop; I use the school computers.} « En fait, je n'ai pas d'ordinateur portable ; j'utilise les ordinateurs de l'école. » Pour dire « actuellement », employez \eng{currently} ou \eng{at present} : \eng{I am currently preparing for the Baccalauréat.} « Je prépare actuellement le Baccalauréat. » C'est l'une des erreurs les plus fréquentes dans les copies du Bac — et l'une des plus sanctionnées.
\end{attention}

\begin{attention}[To assist et a library : deux autres pièges]
\begin{itemize}
\item \eng{To assist} = \cle{aider} : \eng{The technician assisted me with the installation.} « Le technicien m'a aidé pour l'installation. » « Assister à » une réunion se dit \eng{to attend a meeting}.
\item \eng{A library} = \cle{une bibliothèque} (où l'on emprunte des livres) ; le lieu où l'on achète des livres est \eng{a bookshop} (ou \eng{a bookstore} en anglais américain) : \eng{I bought this e-reader manual at a bookshop in Douala.}
\item \eng{Digital} se rapporte au \cle{numérique} (\eng{digital camera}, \eng{digital native}) ; pour les doigts et les empreintes, l'anglais utilise \eng{finger} : \eng{a fingerprint} « une empreinte digitale ».
\end{itemize}
\end{attention}

\begin{savaistu}[Le saviez-vous ? : Les natifs du numérique]
L'expression \eng{digital native} « natif du numérique » désigne une personne née dans un monde déjà équipé des technologies numériques : ordinateurs, téléphones portables, Internet. La plupart d'entre vous sont des \emph{digital natives}, contrairement à vos parents ou grands-parents, que l'on appelle parfois \eng{digital immigrants} — ils ont dû « immigrer » dans le monde numérique en l'apprenant adultes ! Au Cameroun, les cybercafés de Yaoundé, Douala ou Bamenda ont longtemps été la porte d'entrée de toute une génération vers le \eng{cyberspace}.
\end{savaistu}

\courssub{Practice: a solved exercise}

\begin{exoresolu}[Collocations, familles de mots et faux amis]
\textbf{A.} Complétez chaque phrase avec la bonne collocation : \emph{browse / log in / charge / post / download}.

\begin{enumerate}
\item Before you can read the lesson, you must \ldots\ with your password.
\item I always \ldots\ my phone at night.
\item Students can \ldots\ the file from the school platform.
\item Many teenagers \ldots\ the web instead of reading books.
\item Please \ldots\ a message in the group to confirm your presence.
\end{enumerate}

\textbf{B.} Donnez la forme correcte du mot entre parenthèses.

\begin{enumerate}
\item We need a good \ldots\ to attend the online lesson. (connect)
\item This website is very \ldots\ : I learnt a lot. (inform)
\item The school library is now fully \ldots\ . (computer)
\end{enumerate}

\textbf{C.} Traduisez en anglais : « En fait, j'assiste actuellement à un cours en ligne à la bibliothèque. »

\tcblower

\textbf{A.} 1. \eng{log in} (se connecter avec un mot de passe) — 2. \eng{charge} (charger la batterie) — 3. \eng{download} (télécharger depuis la plateforme) — 4. \eng{browse} (\eng{browse the web} : naviguer sur le web) — 5. \eng{post} (\eng{post a message} : publier un message).

\textbf{B.} 1. \eng{connection} (nom après l'adjectif \emph{good}) — 2. \eng{informative} (adjectif après \emph{is very}) — 3. \eng{computerized} (adjectif : « informatisée »).

\textbf{C.} \eng{Actually, I am currently attending an online course at the library.} « En fait » = \eng{actually} ; « actuellement » = \eng{currently} ; « assister à » = \eng{to attend} ; « bibliothèque » = \eng{library}. Quatre faux amis évités dans une seule phrase !
\end{exoresolu}

\begin{exercice}[À vous de jouer]
\textbf{1.} Trouvez l'intrus dans chaque série de collocations et justifiez : (a) \eng{to surf the Internet / to browse the web / to swim the Internet} ; (b) \eng{a powerful computer / a strong computer / a reliable connection}.

\textbf{2.} Complétez le tableau de famille de mots : \eng{to communicate} (verbe), \ldots\ (nom), \ldots\ (adjectif).

\textbf{3.} Corrigez les erreurs dans ces phrases d'élèves : (a) \emph{I need more informations about e-learning.} (b) \emph{Actually, I am studying in the librairie.} (c) \emph{She assisted the video conference yesterday.}

\textbf{4.} Rédigez un court paragraphe (cinq phrases) intitulé \eng{How I use ICTs to learn English}, en employant au moins quatre collocations de la leçon et un mot de chaque famille étudiée.
\end{exercice}

\begin{corrige}[Exercice « À vous de jouer »]
\textbf{1.} (a) \eng{to swim the Internet} est l'intrus : « surfer » sur Internet se dit \eng{to surf} ou \eng{to browse}, jamais \eng{to swim} (nager). (b) \eng{a strong computer} est l'intrus : la collocation correcte est \eng{a powerful computer}.

\textbf{2.} \eng{communication} (nom), \eng{communicative} (adjectif).

\textbf{3.} (a) \eng{I need more information about e-learning.} (\eng{information} est indénombrable, sans \emph{s}). (b) \eng{Currently, I am studying in the library.} (« actuellement » = \eng{currently} ; « bibliothèque » = \eng{library}). (c) \eng{She attended the video conference yesterday.} (« assister à » = \eng{to attend} ; \eng{to assist} = aider).

\textbf{4.} Modèle de production : \eng{Every evening, I log in to an e-learning platform and download audio files to improve my listening skills. I browse the web to read short articles in English, and I post messages in an English learners' group. My connection is not always reliable, but the cybercafe near my house has high-speed internet. Thanks to good connectivity, I can communicate with other learners and find informative websites. Computing has truly enhanced my learning.}
\end{corrige}

\begin{aretenir}[L'essentiel de la section]
\begin{itemize}
\item Les \cle{collocations} s'apprennent en blocs : \eng{to browse the web}, \eng{to log in}, \eng{to download a file}, \eng{a reliable connection}, \eng{a user-friendly app}, \eng{high-speed internet}.
\item Les \cle{familles de mots} permettent de passer du verbe au nom et à l'adjectif : \eng{connect → connection → connected → connectivity} ; \eng{compute → computer → computing → computerized}.
\item Les préfixes \cle{e-} (electronic) et \cle{cyber-} (computers and the Internet) et les suffixes \cle{-er}/\cle{-or} (\eng{printer}, \eng{scanner}, \eng{projector}) forment l'essentiel du vocabulaire des TIC.
\item Méfiez-vous des \cle{faux amis} : \eng{actually} = en fait (\eng{currently} = actuellement) ; \eng{to assist} = aider (\eng{to attend} = assister à) ; \eng{a library} = bibliothèque (\eng{a bookshop} = librairie) ; \eng{digital} = numérique (\eng{fingerprint} = empreinte digitale).
\item \eng{Information} est indénombrable : jamais de \emph{s}.
\end{itemize}
\end{aretenir}

\coursec{Using the vocabulary in context}

\courssub{Transition from the previous section}
\begin{textebox}[Transition]
In the previous section, we explored \cle{collocations} (natural word partnerships like \eng{download a file} or \eng{slow connection}), \cle{word families} (e.g. \eng{connect, connection, connected, connectivity}) and \cle{false friends} (e.g. \eng{actually} $\neq$ \fr{actuellement}). Now we will see how all this vocabulary works together in a coherent paragraph. Writing or speaking about ICTs requires not only knowing isolated words but also structuring ideas with appropriate \cle{connectors} (\eng{first, moreover, however, in conclusion}) and using the right collocations in context.
\end{textebox}

\courssub{Model paragraph: “How I Use ICTs to Prepare for the Bac”}
\begin{textebox}[Model paragraph — “How I Use ICTs to Prepare for the Bac”]
ICTs play an important role in my studies. \textbf{First}, I use my \textbf{smartphone} to \textbf{download past papers} and \textbf{watch video tutorials} when I do not understand a lesson. \textbf{Moreover}, our class \textbf{WhatsApp group} allows us to \textbf{share notes} and \textbf{audio explanations}. \textbf{However}, the \textbf{connection} in my area is sometimes \textbf{slow}, so I often go to a \textbf{cybercafé}. \textbf{In conclusion}, ICTs help me \textbf{learn at my own pace} and \textbf{prepare for the Bac} more \textbf{efficiently}.
\end{textebox}

\begin{definition}[Glossaire du paragraphe modèle]
\begin{itemize}
\item \eng{ICTs} (n. pl.) = \fr{TIC (Technologies de l'Information et de la Communication)}
\item \eng{play a role} (colloc.) = \fr{jouer un rôle}
\item \eng{past papers} (n. pl.) = \fr{annales (sujets d'examens précédents)}
\item \eng{video tutorial} (n.) = \fr{tutoriel vidéo}
\item \eng{share notes} (colloc.) = \fr{partager des notes / des résumés}
\item \eng{audio explanation} (n.) = \fr{explication audio}
\item \eng{connection} (n.) = \fr{connexion (Internet)}
\item \eng{slow} (adj.) = \fr{lent(e)}
\item \eng{cybercafé} (n.) = \fr{cybercafé}
\item \eng{learn at my own pace} (colloc.) = \fr{apprendre à mon propre rythme}
\item \eng{prepare for} (v. + prep.) = \fr{se préparer pour / à}
\item \eng{efficiently} (adv.) = \fr{efficacement}
\end{itemize}
\end{definition}

\courssub{Analyse du modèle : structure, champs lexicaux et connecteurs}
\begin{methode}[Comment analyser un paragraphe sur les TIC]
\begin{enumerate}
\item \textbf{Repérer la phrase d'introduction (topic sentence)} : elle présente l'idée générale. Ici : \eng{ICTs play an important role in my studies.}
\item \textbf{Identifier les connecteurs de structure} : \eng{First} (premier argument), \eng{Moreover} (argument additionnel), \eng{However} (concession / nuance), \eng{In conclusion} (synthèse + opinion).
\item \textbf{Surligner les collocations du thème} : \eng{download past papers}, \eng{watch video tutorials}, \eng{share notes}, \eng{audio explanations}, \eng{slow connection}, \eng{learn at my own pace}, \eng{prepare for the Bac}, \eng{work efficiently}.
\item \textbf{Noter les champs lexicaux} : \textit{Devices} (\eng{smartphone}), \textit{Internet / Access} (\eng{connection, cybercafé}), \textit{Learning actions} (\eng{download, watch, share, learn, prepare}).
\item \textbf{Observer la nuance} : \eng{However} introduit un inconvénient réel (\eng{slow connection}) et une solution (\eng{go to a cybercafé}), ce qui rend le paragraphe crédible.
\end{enumerate}
\end{methode}

\begin{popfigure}
\begin{tikzpicture}[
    node distance=1.2cm and 1.5cm,
    box/.style={rectangle, rounded corners, draw=popblue, thick, fill=popblueL, text width=3.2cm, align=center, font=\small},
    arrow/.style={-Stealth, thick, popdark}
]
\node[box, align=center] (intro){Introduction\\General statement\\\eng{ICTs play an important role in my studies.}};
\node[box, below=of intro, align=center] (body1){Body 1 — Example 1\\\eng{First} + collocations\\\eng{download past papers, watch video tutorials}};
\node[box, below=of body1, align=center] (body2){Body 2 — Example 2\\\eng{Moreover} + collocations\\\eng{share notes, audio explanations}};
\node[box, below=of body2, align=center] (nuance){Nuance / Concession\\\eng{However} + problem + solution\\\eng{slow connection $\rightarrow$ cybercafé}};
\node[box, below=of nuance, align=center] (concl){Conclusion\\\eng{In conclusion} + personal opinion\\\eng{learn at my own pace, prepare efficiently}};

\draw[arrow] (intro) -- (body1);
\draw[arrow] (body1) -- (body2);
\draw[arrow] (body2) -- (nuance);
\draw[arrow] (nuance) -- (concl);
\end{tikzpicture}
\legende{Structure du paragraphe modèle : Introduction $\rightarrow$ Corps (exemples + collocations) $\rightarrow$ Nuance $\rightarrow$ Conclusion}
\end{popfigure}

\begin{propriete}[Connecteurs utiles pour structurer un paragraphe sur les TIC]
\begin{center}
\begin{tabularx}{\linewidth}{|X|X|X|}
\hline
\rowcolor{popblueL} \textbf{Connecteur} & \textbf{Fonction} & \textbf{Exemple du modèle} \\ \hline
\eng{First} / \eng{Firstly} & Introduire le premier argument & \eng{First, I use my smartphone…} \\ \hline
\eng{Moreover} / \eng{Furthermore} / \eng{In addition} & Ajouter un argument & \eng{Moreover, our class WhatsApp group…} \\ \hline
\eng{However} / \eng{Nevertheless} / \eng{On the other hand} & Introduire une concession / un inconvénient & \eng{However, the connection… is sometimes slow} \\ \hline
\eng{Therefore} / \eng{Consequently} / \eng{So} & Exprimer la conséquence & \eng{…so I often go to a cybercafé.} \\ \hline
\eng{In conclusion} / \eng{To sum up} / \eng{Finally} & Conclure + opinion personnelle & \eng{In conclusion, ICTs help me learn at my own pace…} \\ \hline
\end{tabularx}
\end{center}
\end{propriete}

\courssub{Production guidée : phrases à trous (gap-fill)}
\begin{exercice}[Complétez les phrases avec le mot ou l'expression appropriée (choisissez dans la boîte)]
\begin{textebox}[Boîte à mots]
\textbf{download past papers} \quad \textbf{video tutorials} \quad \textbf{share notes} \quad \textbf{slow connection} \quad \textbf{cybercafé} \quad \textbf{learn at my own pace} \quad \textbf{prepare for} \quad \textbf{efficiently} \quad \textbf{WhatsApp group} \quad \textbf{smartphone}
\end{textebox}
\begin{enumerate}
\item I use my \_\_\_\_\_\_\_\_\_\_\_\_ to take photos of the blackboard after the lesson.
\item Our class \_\_\_\_\_\_\_\_\_\_\_\_ is very active; we \_\_\_\_\_\_\_\_\_\_\_\_ and audio explanations every evening.
\item When the \_\_\_\_\_\_\_\_\_\_\_\_ at home is too slow, I go to the nearest \_\_\_\_\_\_\_\_\_\_\_\_.
\item To \_\_\_\_\_\_\_\_\_\_\_\_ the Bac, I \_\_\_\_\_\_\_\_\_\_\_\_ from the official website.
\item If I don't understand a topic, I watch \_\_\_\_\_\_\_\_\_\_\_\_ on YouTube.
\item Thanks to ICTs, I can \_\_\_\_\_\_\_\_\_\_\_\_ and revise \_\_\_\_\_\_\_\_\_\_\_\_.
\end{enumerate}
\end{exercice}

\begin{corrige}[Exercice — Production guidée]
\begin{enumerate}
\item \textbf{smartphone} — \textit{Device personnel le plus courant.}
\item \textbf{WhatsApp group} / \textbf{share notes} — \textit{Collocations : \eng{share notes} (partager des notes), pas \eng{share lessons}.}
\item \textbf{slow connection} / \textbf{cybercafé} — \textit{\eng{Connection} (nom) $\neq$ \eng{connect} (verbe).}
\item \textbf{prepare for} / \textbf{download past papers} — \textit{\eng{Prepare for} + nom / gérondif. \eng{Past papers} = annales.}
\item \textbf{video tutorials} — \textit{Pluriel de \eng{tutorial}.}
\item \textbf{learn at my own pace} / \textbf{efficiently} — \textit{Collocation figée \eng{learn at one's own pace} ; adverbe \eng{efficiently} (pas \eng{efficient}).}
\end{enumerate}
\end{corrige}

\courssub{Production semi-libre : 5 phrases personnelles}
\begin{methode}[Rédiger 5 phrases vraies sur votre usage des TIC]
\begin{enumerate}
\item \textbf{Phrase 1 — Généralité} : Utilisez \eng{ICTs play an important role in…} ou \eng{I often use… for…}
\item \textbf{Phrase 2 — Exemple précis (learning)} : \eng{First, I…} + collocation (\eng{download past papers}, \eng{watch video tutorials}, \eng{use an app to revise vocabulary}).
\item \textbf{Phrase 3 — Exemple précis (communication / collaboration)} : \eng{Moreover, I…} + \eng{share notes on…}, \eng{join a study group on…}, \eng{ask questions on…}
\item \textbf{Phrase 4 — Nuance (avantage / inconvénient)} : \eng{However, …} + \eng{the connection is slow / data is expensive / I get distracted by social media}.
\item \textbf{Phrase 5 — Conclusion / Opinion} : \eng{In conclusion, ICTs help me…} + \eng{learn at my own pace / work more efficiently / stay organised}.
\end{enumerate}
\end{methode}

\begin{exemplebox}[Exemples d'élèves (modèles)]
\begin{enumerate}
\item ICTs play an important role in my daily life as a student in Douala.
\item First, I use my smartphone to \textbf{download past papers} and \textbf{watch video tutorials} on Khan Academy.
\item Moreover, our class \textbf{WhatsApp group} allows us to \textbf{share notes} and \textbf{ask questions} before exams.
\item However, mobile \textbf{data is expensive}, so I often wait for \textbf{free Wi-Fi} at school.
\item In conclusion, ICTs help me \textbf{prepare for the Bac} more \textbf{efficiently} and \textbf{learn at my own pace}.
\end{enumerate}
\end{exemplebox}

\begin{exercice}[Votre tour : écrivez vos 5 phrases]
Écrivez vos 5 phrases dans votre cahier en suivant le modèle ci-dessus. Utilisez au moins \textbf{6 collocations} vues dans la leçon (ex. \eng{slow connection, video tutorial, share notes, download past papers, learn at my own pace, prepare for, efficiently, cybercafé, WhatsApp group, smartphone}).
\end{exercice}

\courssub{Expression orale en binômes : Interview « How do ICTs help you learn or work? »}
\begin{experience}[Interview en binômes — Grille de questions]
\begin{enumerate}
\item \textbf{Student A (journalist)} : \eng{“Hello! Can you tell me how ICTs help you in your studies or work?”}
\item \textbf{Student B (interviewé)} : Répond en utilisant la structure du paragraphe modèle (Introduction → 2 exemples → Nuance → Conclusion).
\item \textbf{Questions de relance (Student A)} :
\begin{itemize}
\item \eng{“What device do you use most?”}
\item \eng{“Which apps or websites are the most useful?”}
\item \eng{“Do you ever have connection problems? What do you do then?”}
\item \eng{“Do you think ICTs make you more efficient? Why?”}
\end{itemize}
\item \textbf{Changement de rôle} : Student B devient journaliste, Student A répond.
\item \textbf{Retour en classe} : 2-3 binômes présentent leur interview devant la classe. Le professeur note les bonnes collocations et les erreurs au tableau pour la correction collective.
\end{enumerate}
\end{experience}

\begin{aretenir}[Check-list pour l'oral]
\begin{itemize}
\item \checkmark J'utilise \textbf{au moins 3 connecteurs} (\eng{First, Moreover, However, In conclusion}).
\item \checkmark J'utilise \textbf{au moins 4 collocations} correctes (\eng{download past papers, slow connection, share notes, learn at my own pace…}).
\item \checkmark Je prononce correctement : \eng{connection} « co-nèk-chn », \eng{efficiently} « i-fi-chène-li », \eng{tutorial} « tiou-to-ri-al ».
\item \checkmark Je ne fais pas de calque du français (voir boîte \textbf{Attention} ci-dessous).
\end{itemize}
\end{aretenir}

\courssub{Correction collective : calques du français et erreurs fréquentes}
\begin{attention}[Ne pas traduire mot à mot depuis le français — Pièges classiques]
\begin{center}
\begin{tabularx}{\linewidth}{|X|X|X|}
\hline
\rowcolor{poporangeL} \textbf{Calque incorrect (FR $\rightarrow$ EN littéral)} & \textbf{Pourquoi c'est faux} & \textbf{Forme correcte (collocation naturelle)} \\ \hline
\eng{I inform myself on the Internet.} & \fr{« Je m'informe » $\neq$ \eng{inform myself} (réfléchi rare)} & \eng{I look for information on the Internet.} / \eng{I get information online.} \\ \hline
\eng{I navigate on the Internet.} & \fr{\eng{Navigate} = naviguer (bateau, GPS) ; pour le web : \eng{browse / surf}} & \eng{I browse the Internet.} / \eng{I surf the web.} \\ \hline
\eng{I download a video on YouTube.} & \fr{YouTube = streaming ; on ne « télécharge » pas légalement depuis l'appli} & \eng{I watch a video on YouTube.} / \eng{I stream a video.} \\ \hline
\eng{The connexion is bad.} & \fr{Orthographe : \eng{connexion} (FR) $\rightarrow$ \eng{connection} (EN)} & \eng{The connection is slow / weak / unstable.} \\ \hline
\eng{I learn English by Internet.} & \fr{Préposition : \eng{by} + V-ing / \eng{via} / \eng{using} / \eng{on}} & \eng{I learn English \textbf{on} the Internet.} / \eng{I learn English \textbf{using} ICTs.} \\ \hline
\eng{I make my homework on the computer.} & \fr{\eng{Make} = fabriquer ; \eng{do} = faire (devoir, travail)} & \eng{I \textbf{do} my homework on the computer.} \\ \hline
\eng{I pass my time on social media.} & \fr{\eng{Pass time} = tuer le temps ; \eng{spend time} = passer du temps (neutre)} & \eng{I \textbf{spend} time on social media.} \\ \hline
\eng{The network doesn't take.} & \fr{Calque de « le réseau ne prend pas »} & \eng{There is \textbf{no signal}.} / \eng{I have \textbf{no reception}.} \\ \hline
\end{tabularx}
\end{center}
\end{attention}

\begin{exoresolu}[Correction collective d'exemples anonymes]
\textbf{Énoncé :} Corrigez les phrases suivantes écrites par des élèves (erreurs de calque, collocation, orthographe, préposition).
\begin{enumerate}
\item I inform myself on Google for my homework.
\item The connexion in my village is very bad.
\item I navigate on Facebook to find my friends.
\item I download videos on TikTok to watch them later.
\item I make my research on Wikipedia.
\item I pass three hours on my phone every day.
\item The network doesn't take in my bedroom.
\item I learn my lessons by Internet.
\end{enumerate}
\tcblower
\textbf{Solution détaillée :}
\begin{enumerate}
\item \eng{I \textbf{look for information} on Google for my homework.} (\eng{inform myself} $\times$)
\item \eng{The \textbf{connection} in my village is \textbf{very slow / weak}.} (\eng{connexion} $\times$, \eng{bad} vague $\rightarrow$ \eng{slow/weak})
\item \eng{I \textbf{browse / scroll} Facebook to find my friends.} (\eng{navigate} $\times$ pour les réseaux sociaux)
\item \eng{I \textbf{save / bookmark} videos on TikTok to watch them later.} (\eng{download} $\times$ hors appli ; \eng{save} = enregistrer dans l'appli)
\item \eng{I \textbf{do} my research on Wikipedia.} (\eng{make research} $\times$ ; \eng{do research} ou \eng{carry out research})
\item \eng{I \textbf{spend} three hours on my phone every day.} (\eng{pass time} $\times$ ; \eng{spend time})
\item \eng{There is \textbf{no signal} in my bedroom.} / \eng{I have \textbf{no reception} in my bedroom.} (\eng{network doesn't take} $\times$)
\item \eng{I learn my lessons \textbf{on} the Internet.} / \eng{I learn my lessons \textbf{using} the Internet.} (\eng{by Internet} $\times$)
\end{enumerate}
\end{exoresolu}

\courssub{Synthèse visuelle : Carte mentale du vocabulaire en contexte}
\begin{popfigure}
\begin{center}\tcbox[colback=popblue,colframe=popblue,coltext=white,arc=9pt,boxsep=4pt,left=6pt,right=6pt]{\bfseries\small ICTs in Context}\end{center}
\vspace{-2pt}
\begin{tcbraster}[raster columns=2,raster equal height=rows,raster column skip=6pt,raster row skip=6pt,enhanced,blankest]
\begin{tcolorbox}[enhanced,colback=white,colframe=poporange,boxrule=0.9pt,arc=3pt,title={Devices \& Access},fonttitle=\bfseries\scriptsize,coltitle=white,colbacktitle=poporange,left=4pt,right=4pt,top=2pt,bottom=2pt,boxsep=1pt,toptitle=1.5pt,bottomtitle=1.5pt,halign=left]
\begin{itemize}[nosep,leftmargin=*,itemsep=1pt,font=\scriptsize]
\item smartphone
\item laptop / computer
\item connection
\item cybercafé
\item Wi-Fi / data
\end{itemize}
\end{tcolorbox}
\begin{tcolorbox}[enhanced,colback=white,colframe=popgreen,boxrule=0.9pt,arc=3pt,title={Learning Actions},fonttitle=\bfseries\scriptsize,coltitle=white,colbacktitle=popgreen,left=4pt,right=4pt,top=2pt,bottom=2pt,boxsep=1pt,toptitle=1.5pt,bottomtitle=1.5pt,halign=left]
\begin{itemize}[nosep,leftmargin=*,itemsep=1pt,font=\scriptsize]
\item download past papers
\item watch video tutorials
\item share notes
\item learn at my own pace
\item prepare for exams
\end{itemize}
\end{tcolorbox}
\begin{tcolorbox}[enhanced,colback=white,colframe=poppurple,boxrule=0.9pt,arc=3pt,title={Work Actions},fonttitle=\bfseries\scriptsize,coltitle=white,colbacktitle=poppurple,left=4pt,right=4pt,top=2pt,bottom=2pt,boxsep=1pt,toptitle=1.5pt,bottomtitle=1.5pt,halign=left]
\begin{itemize}[nosep,leftmargin=*,itemsep=1pt,font=\scriptsize]
\item send emails
\item join video calls
\item share files / screens
\item work remotely
\item use cloud storage
\end{itemize}
\end{tcolorbox}
\begin{tcolorbox}[enhanced,colback=white,colframe=poppink,boxrule=0.9pt,arc=3pt,title={Connectors \& Structure},fonttitle=\bfseries\scriptsize,coltitle=white,colbacktitle=poppink,left=4pt,right=4pt,top=2pt,bottom=2pt,boxsep=1pt,toptitle=1.5pt,bottomtitle=1.5pt,halign=left]
\begin{itemize}[nosep,leftmargin=*,itemsep=1pt,font=\scriptsize]
\item First / Firstly
\item Moreover / Furthermore
\item However / Nevertheless
\item Therefore / So
\item In conclusion / To sum up
\end{itemize}
\end{tcolorbox}
\end{tcbraster}

\legende{Carte mentale récapitulative : vocabulaire, collocations et connecteurs pour parler des TIC}
\end{popfigure}

\begin{exercice}[Entraînement final — Mini-paragraphe écrit]
Rédigez un paragraphe de 8-10 lignes (environ 80-100 mots) intitulé \eng{« How ICTs Help Me in My Daily Life »}. Respectez la structure : \textbf{Introduction} $\rightarrow$ \textbf{2 exemples avec collocations} $\rightarrow$ \textbf{Nuance (However)} $\rightarrow$ \textbf{Conclusion (In conclusion)}. Utilisez au moins \textbf{8 mots/expressions} de la leçon et \textbf{3 connecteurs} différents. Échangez votre copie avec un camarade pour une relecture croisée (vérifiez les calques du français !).
\end{exercice}

\begin{corrige}[Exercice — Mini-paragraphe écrit (modèle de correction)]
\textbf{Exemple de production d'élève (niveau attendu) :}
\eng{ICTs are essential in my daily life as a Terminale student in Yaoundé. First, I use my laptop to do research for my projects and download past papers for the Bac. Moreover, I join online study groups on WhatsApp where we share notes and explain difficult lessons to each other. However, the Internet connection at home is often unstable, so I sometimes work offline or go to a cybercafé. In conclusion, ICTs allow me to learn at my own pace and prepare for my future more efficiently.}
\\
\textbf{Points vérifiés :} structure respectée, 9 collocations (\eng{do research, download past papers, join study groups, share notes, explain lessons, unstable connection, work offline, learn at my own pace, prepare for}), 4 connecteurs (\eng{First, Moreover, However, In conclusion}), aucun calque.
\end{corrige}

\newpage

\coursec{Activité d'intégration}

\courssub{Objectifs de l’activité}
\noindent À l’issue de cette activité d’intégration, l’élève sera capable de~:
\begin{itemize}
    \item \textbf{mobiliser} le vocabulaire des TIC (appareils, outils, Internet, apprentissage, travail) dans une production écrite et orale cohérente~;
    \item \textbf{employer} correctement au moins cinq \cle{collocations} (associations de mots figées) et trois \cle{mots de la même famille} (dérivation) étudiés lors de la leçon~;
    \item \textbf{identifier} et \textbf{éviter} les \cle{faux amis} et les calques syntaxiques du français (ex. \eng{download} vs \eng{upload}, \eng{reliable} vs \eng{fiable}, \eng{to attend a meeting} vs \eng{to assist})~;
    \item \textbf{structurer} un guide pratique clair, impératif et visuellement organisé, adapté à un public cible (élèves de Seconde)~;
    \item \textbf{présenter} son travail à l’oral en 2 minutes avec une prononciation soignée (accent tonique, sons \eng{/ð/}, \eng{/θ/}, \eng{/æ/}) et une intonation engageante.
\end{itemize}

\courssub{Situation}
\begin{textebox}{Contexte : Semaine d’intégration au Lycée Bilingue de Buea}
\eng{The administration of Lycée Bilingue de Buea is organising an ``Integration Week'' for the new Form 1 (Seconde) students arriving in September. The Head Teacher has launched a challenge for the Upper Sixth (Terminale) classes:}

\eng{``Design a one-page `Digital Survival Guide' to help the newcomers use ICTs effectively and safely for their studies. The best guide will be printed and displayed in the school library and computer lab.''}

\eng{Your group (4 students) has decided to take part. You have 45 minutes to create the guide and 2 minutes to pitch it to the class. The guide must contain exactly 5 numbered tips. Each tip must include at least one collocation from the lesson (e.g., \textbf{surf the web}, \textbf{back up files}, \textbf{submit an assignment online}) and one word from a word family studied (e.g., \textbf{connection} $\rightarrow$ \textbf{connect}, \textbf{connected}, \textbf{connectivity}). You must avoid false friends like using \eng{actually} for ``actuellement'' or \eng{assist} for ``assister à''.}
\end{textebox}

\courssub{Consignes}
\noindent \textbf{Étape 1 ~: Brainstorming et sélection lexicale (10 min)}
\begin{enumerate}
    \item En groupe, dressez la liste des 15 mots/expressions clés de la leçon (p. ex. \eng{cloud storage, broadband, search engine, plagiarism, shortcut, user-friendly, to log in, to crash, bandwidth, digital literacy}).
    \item Sélectionnez 5 \cle{collocations} et 3 \cle{familles de mots} que vous intégrerez obligatoirement. Cochez-les sur votre brouillon.
    \item Identifiez 3 \cle{faux amis} à surveiller absolument (ex. \eng{eventually} $\neq$ ``éventuellement'', \eng{library} $\neq$ ``librairie'', \eng{formidable} $\neq$ ``formidable'').
\end{enumerate}

\textbf{Étape 2 ~: Rédaction du « Digital Survival Guide » (20 min)}
\begin{enumerate}
    \setcounter{enumi}{3}
    \item Rédigez 5 conseils numérotés, courts (1--2 phrases), à l’impératif ou au modal \eng{should/must}.
    \item \textbf{Contrainte lexicale~:} Chaque conseil doit contenir \textbf{au moins une collocation} (surlignée en jaune sur la version finale) et \textbf{au moins un mot de famille} (souligné).
    \item \textbf{Contrainte langagière~:} Zéro faux ami. Vérifiez \eng{reliable} (fiable) vs \eng{reliable source}, \eng{download} (télécharger) vs \eng{upload} (téléverser), \eng{attend a webinar} (participer) vs \eng{assist} (aider).
    \item Mettez en page sur une feuille A4 (titre accrocheur, icônes, police lisible, numérotation claire).
\end{enumerate}

\textbf{Étape 3 ~: Révision par les pairs (5 min)}
\begin{enumerate}
    \setcounter{enumi}{7}
    \item Échangez votre guide avec un autre groupe. Vérifiez~: collocations présentes ? Mots de famille présents ? Faux amis absents ? Orthographe correcte (\eng{receive, necessary, separate, accommodation}) ? Clarté du message ?
    \item Annotez les corrections au crayon et rendez la feuille.
\end{enumerate}

\textbf{Étape 4 ~: Présentation orale et vote (10 min)}
\begin{enumerate}
    \setcounter{enumi}{9}
    \item Chaque groupe présente son guide en \textbf{2 minutes chrono} (un ou deux orateurs). Utilisez des connecteurs logiques (\eng{Firstly, Moreover, Finally}) et un ton convaincant.
    \item La classe vote à bulletin secret pour le guide le plus \eng{useful}, le plus \eng{well-designed} et le plus \eng{accurate linguistically}.
    \item Le guide gagnant est affiché au tableau de la salle / au CDI.
\end{enumerate}

\courssub{Ressources à mobiliser}
\begin{propriete}[Rappel : Collocations clés de la leçon (Verbe + Nom / Adjectif + Nom)]
\begin{center}
\begin{tabularx}{\linewidth}{|X|X|X|}
\hline
\rowcolor{popblueL} \textbf{Collocation} & \textbf{Traduction / Contexte} & \textbf{Famille de mots (Word Family)} \\
\hline
\eng{surf the web / browse the Internet} & naviguer sur le Web & \eng{browser, browsing, surfer} \\
\hline
\eng{download / upload files} & télécharger / téléverser des fichiers & \eng{download (n), upload (n), downloadable, uploadable} \\
\hline
\eng{back up data / back up notes} & sauvegarder des données / notes & \eng{backup (n), backing up} \\
\hline
\eng{submit an assignment online} & rendre un devoir en ligne & \eng{submission, submittable} \\
\hline
\eng{attend a webinar / virtual class} & assister à un webinaire / classe virtuelle & \eng{attendance, attendee, attending} \\
\hline
\eng{reliable source / trustworthy website} & source fiable / site de confiance & \eng{reliability, rely on, unreliable} \\
\hline
\eng{high-speed broadband / connectivity} & haut débit / connectivité & \eng{connect, connection, connected, disconnect} \\
\hline
\eng{enhance learning / improve skills} & améliorer l'apprentissage / compétences & \eng{enhancement, enhanced, skilful} \\
\hline
\eng{avoid plagiarism / cite sources} & éviter le plagiat / citer ses sources & \eng{plagiarize, plagiarist, citation, cite} \\
\hline
\eng{log in / log out / sign up} & se connecter / se déconnecter / s'inscrire & \eng{login (n), logout (n), sign-up (n)} \\
\hline
\end{tabularx}
\end{center}
\end{propriete}

\begin{popfigure}
\begin{tikzpicture}[
    node distance=1.2cm and 1.5cm,
    main/.style={draw=popblue, thick, fill=popblueL, rounded corners, minimum width=2.5cm, minimum height=1cm, align=center, font=\small},
    arrow/.style={-Stealth, thick, popdark}
]
\node[main, align=center] (start){Choisir une\\ \cle{collocation}};
\node[main, right=of start, align=center] (family){Ajouter un mot de\\ \cle{famille}};
\node[main, right=of family, align=center] (check){Vérifier les\\ \cle{faux amis}};
\node[main, below=of family, align=center] (write){Rédiger à l'\textbf{impératif}\\ + justification};
\node[main, left=of write, align=center] (tip){\textbf{Tip final}\\ clair \& correct};

\draw[arrow] (start) -- (family);
\draw[arrow] (family) -- (check);
\draw[arrow] (check) -- (write);
\draw[arrow] (write) -- (tip);
\draw[arrow] (family) -- (write);
\end{tikzpicture}
\legende{Processus de rédaction d'un conseil (Tip) efficace}
\end{popfigure}

\begin{attention}[Faux amis et pièges fréquents pour francophones]
\begin{itemize}
    \item \eng{Actually} = « en fait, en réalité » (pas « actuellement » $\rightarrow$ \eng{currently / at present}).
    \item \eng{Assist} = « aider, assister quelqu'un » (pas « assister à un événement » $\rightarrow$ \eng{attend}).
    \item \eng{Eventually} = « finalement, à la fin » (pas « éventuellement » $\rightarrow$ \eng{possibly / potentially}).
    \item \eng{Library} = « bibliothèque » (pas « librairie » $\rightarrow$ \eng{bookshop / bookstore}).
    \item \eng{Formidable} = « redoutable, impressionnant » (pas « excellent » $\rightarrow$ \eng{great, amazing, fantastic}).
    \item \eng{Sensible} = « sensé, raisonnable » (pas « sensible » au sens « qui sent » ou « émotif » $\rightarrow$ \eng{sensitive}).
    \item \eng{Download} (vers son appareil) vs \eng{Upload} (vers un serveur/le cloud). Ne pas confondre.
    \item \eng{Pass an exam} = « réussir un examen » (pas « passer » $\rightarrow$ \eng{take / sit an exam}).
    \item \eng{Programme} (UK spelling) vs \eng{Program} (code informatique). Attention à l'orthographe britannique.
    \item \eng{Fibre} (UK) vs \eng{Fiber} (US). Au Cameroun : \eng{fibre}.
\end{itemize}
\end{attention}

\begin{methode}[Structurer un conseil efficace (Tip Structure)]
\begin{enumerate}
    \item \textbf{Verbe d'action fort} (Impératif ou \eng{Should/Must}) : \eng{Back up}, \eng{Use}, \eng{Avoid}, \eng{Check}.
    \item \textbf{Collocation précise} : \eng{your files on cloud storage}, \eng{a reliable search engine}.
    \item \textbf{Justification courte} (cause/conséquence) : \eng{so that you don't lose them}, \eng{to find accurate information}.
    \item \textbf{Mot de la famille} intégré naturellement : \eng{connectivity}, \eng{submission}, \eng{reliability}.
\end{enumerate}
\eng{Example~: \textbf{Back up} your notes \textbf{on cloud storage} daily \textbf{to prevent data loss}; a stable \textbf{connection} ensures automatic \textbf{synchronisation}.} « Sauvegardez vos notes sur le cloud quotidiennement pour éviter la perte de données ; une connexion stable assure une synchronisation automatique. »
\end{methode}

\courssub{Proposition de production et corrigé commenté}
\begin{exoresolu}{Modèle de « Digital Survival Guide » + Commentaire linguistique}
\textbf{\eng{DIGITAL SURVIVAL GUIDE -- FORM 1 EDITION}}\\
\eng{\textbf{Welcome to LB Buea! Master your tech, ace your studies.}}\\
\vspace{0.3cm}
\noindent
\eng{1. \textbf{Back up} your \textbf{notes on cloud storage} (Google Drive, OneDrive) \textbf{daily}; a sudden \textbf{crash} or lost USB key shouldn't erase your term's work. \textit{(Collocation: back up notes on cloud storage ; Famille: crash $\rightarrow$ crash, crashing ; Grammaire: impératif + ; + proposition conséquence)}.}\\
\vspace{0.2cm}
\eng{2. \textbf{Use reliable sources} and \textbf{cite them correctly} to \textbf{avoid plagiarism}; \textbf{citation} tools like Zotero make \textbf{referencing} easy and honest. \textit{(Collocations: use reliable sources, cite correctly, avoid plagiarism ; Familles: cite $\rightarrow$ citation, referencing ; Faux ami évité: \textbf{plagiarism} $\neq$ ``plagiat'' orthographe fr., \textbf{reliable} $\neq$ ``réel'')}.}\\
\vspace{0.2cm}
\eng{3. \textbf{Attend virtual classes} and \textbf{webinars} on time; stable \textbf{connectivity} and a working \textbf{microphone} ensure active \textbf{participation}. \textit{(Collocations: attend virtual classes, stable connectivity ; Familles: connect $\rightarrow$ connectivity, participate $\rightarrow$ participation ; Faux ami évité: \textbf{attend} $\neq$ \textbf{assist})}.}\\
\vspace{0.2cm}
\eng{4. \textbf{Download past papers} only from \textbf{official school portals} or \textbf{trusted educational platforms}; \textbf{unverified} files may contain \textbf{malware}. \textit{(Collocations: download past papers, trusted platforms ; Familles: verify $\rightarrow$ unverified, trust $\rightarrow$ trusted ; Vocabulaire technique: \textbf{malware})}.}\\
\vspace{0.2cm}
\eng{5. \textbf{Log out} of all accounts on \textbf{shared computers} and \textbf{protect your passwords}; \textbf{digital security} starts with \textbf{awareness}. \textit{(Collocations: log out of, shared computers, protect passwords ; Familles: secure $\rightarrow$ security, aware $\rightarrow$ awareness ; Structure: impératif + ; + nominalisation)}.}\\
\vspace{0.5cm}
\eng{\textbf{STAY CONNECTED, STAY SAFE, SUCCEED!}}

\tcblower
\textbf{Commentaire détaillé des choix de langue~:}
\begin{itemize}
    \item \textbf{Collocations respectées (5/5)~:} \eng{back up ... on cloud storage} (Tip 1), \eng{use reliable sources / cite correctly / avoid plagiarism} (Tip 2), \eng{attend virtual classes / stable connectivity} (Tip 3), \eng{download past papers / trusted platforms} (Tip 4), \eng{log out of / shared computers / protect passwords} (Tip 5). Toutes sont des associations naturelles (corpus BNC/COCA).
    \item \textbf{Familles de mots intégrées (6 occurrences)~:} \eng{crash/crashing} (dérivation verbale), \eng{cite/citation/referencing} (nominalisation + préfixe re-), \eng{connect/connectivity, participate/participation} (suffixes -ity, -ation), \eng{verify/unverified, trust/trusted} (préfixes un-, participes passés adjectivaux), \eng{secure/security, aware/awareness} (suffixes -ity, -ness). Cela démontre la maîtrise de la \cle{dérivation}.
    \item \textbf{Faux amis évités~:} \eng{Attend} (et non \eng{assist}) pour ``assister à'' ; \eng{Reliable} (et non \eng{real/actual}) pour ``fiable'' ; \eng{Plagiarism} (orthographe anglaise conservée, pas ``plagiat'') ; \eng{Currently} non utilisé à tort pour ``actuellement'' (le contexte est général).
    \item \textbf{Grammaire \& Registre~:} Impératif pour l'injonction directe (adapté au guide de survie). Point-virgule (;) pour lier deux indépendantes proches sémantiquement (niveau B2/C1). Vocabulaire technique précis (\eng{malware, cloud storage, portal, citation tools}). Ton encourageant (\eng{Welcome, Master, Ace, Succeed}).
    \item \textbf{Prononciation ciblée~:} \eng{connectivity} (accent sur la 3e syllabe \textbf{tiv}), \eng{malware}, \eng{plagiarism}, \eng{reliable} (diphthongue /aI/), \eng{awareness} (schwa initial). \textit{(Note : la prononciation est notée en alphabet phonétique simplifié entre crochets pour l'enseignant ; l'élève utilise la notation française « kɒ-nek-ti-vi-ti »)}.
\end{itemize}
\end{exoresolu}

\courssub{Bilan de l’activité}
\begin{aretenir}[Auto-évaluation de l'élève (à coller dans le cahier)]
\begin{tabularx}{\linewidth}{|X|c|c|c|}
\hline
\rowcolor{popblueL} \textbf{Compétence} & \textbf{Acquis} & \textbf{En cours} & \textbf{À revoir} \\ \hline
J'utilise 5+ collocations TIC correctes. & $\square$ & $\square$ & $\square$ \\ \hline
J'intègre 3+ mots de familles différentes. & $\square$ & $\square$ & $\square$ \\ \hline
Je n'ai commis aucun faux ami (attend/assist, actually/actuellement...). & $\square$ & $\square$ & $\square$ \\ \hline
Mon guide est clair, visuel, impératif et adressé aux Seconde. & $\square$ & $\square$ & $\square$ \\ \hline
Ma présentation orale (2 min) était fluide, prononcée, convaincante. & $\square$ & $\square$ & $\square$ \\ \hline
\end{tabularx}\par\medskip
\noindent \textbf{Action corrective~:} Si « En cours » ou « À revoir » coché $\rightarrow$ Refaire 3 phrases avec la collocation manquante + fiche de révision « Faux amis TIC » (p. 142 du manuel).
\end{aretenir}

\begin{savaistu}[Le saviez-vous ? -- Le « Digital Divide » au Cameroun]
\eng{According to the International Telecommunication Union (ITU), in 2023, only about 38\% of the Cameroonian population used the Internet. The gap between urban areas (Yaoundé, Douala, Buea -- 4G/5G, fibre) and rural zones (limited bandwidth, expensive data) is a major challenge for \textbf{digital equity} in education. Initiatives like \textbf{``GovTech''} or \textbf{``Digital Cameroon 2025''} aim to equip schools with \textbf{computer labs} and train teachers in \textbf{ICT pedagogy}. Your ``Digital Survival Guide'' is a small but real step towards \textbf{digital literacy} for all!}
\end{savaistu}

\newpage

\coursec{Méthodes et exercices résolus}

\coursec{Lesson 54 — Vocabulary: Words and expressions related to using ICTs to enhance work or learning}

\courssub{Méthodes et exercices résolus}

\begin{methode}[Identifier et classer le vocabulaire technique des TIC (Dispositifs, Outils, Internet)]
\textbf{Objectif :} Repérer les noms d'appareils, de logiciels, de composants réseau et d'actions techniques dans un texte ou une conversation, puis les organiser par catégories pour les mémoriser durablement.

\textbf{Démarche en 4 étapes :}
\begin{enumerate}
    \item \textbf{Balayer le support (texte / audio) :} Souligner ou noter tout mot qui désigne un objet physique (\eng{laptop, smartphone, router, server, USB drive, headset}), un logiciel/service (\eng{app, browser, OS, cloud storage, VPN, LMS}) ou une action technique (\eng{to log in, to upload, to download, to sync, to troubleshoot, to back up}).
    \item \textbf{Classer par champ lexical :} Créer un tableau mental ou papier à 3 colonnes : \textit{Devices (Hardware)} / \textit{Software \& Platforms} / \textit{Actions \& Processes (Verbes + particules)}.
    \item \textbf{Vérifier la prononciation et l'orthographe :} Attention aux mots qui s'écrivent presque comme en français mais se prononcent différemment (\eng{router} [“rou-teur”], \eng{server} [“seur-veur”], \eng{software} [“softeur”], \eng{hardware} [“hard-veur”]). Noter la prononciation entre parenthèses.
    \item \textbf{Construire des collocations types :} Associer chaque nom aux verbes qui lui vont naturellement (\eng{to charge a laptop}, \eng{to update an app}, \eng{to configure a router}, \eng{to access the cloud}). Cela évite les calques (\eng{*to open the computer} au lieu de \eng{to turn on / boot up the laptop}).
\end{enumerate}

\textbf{Structures à retenir :}
\begin{itemize}
    \item \eng{to log in / log into} (se connecter à) $\neq$ \eng{to log out / log out of} (se déconnecter de).
    \item \eng{to sign up for} (s'inscrire à un service) vs \eng{to sign in to} (ouvrir une session).
    \item \eng{a stable / fast / slow connection}, \eng{high-speed broadband}, \eng{Wi-Fi signal}.
\end{itemize}
\end{methode}

\begin{exoresolu}[Analyse lexicale d'un témoignage d'élève à Yaoundé]
\textbf{Énoncé :} Lisez le texte ci-dessous et répondez aux questions.

\begin{textebox}[Témoignage d'une élève de Terminale A au Lycée de Nkolbisson]
My name is Mireille and I am in Upper Sixth at Lycée de Nkolbisson in Yaoundé. Since the COVID-19 pandemic, our school has adopted a \textbf{blended learning} approach. Last week, our Physics teacher asked us to \textbf{submit} our lab reports via the school \textbf{LMS} (Learning Management System) before midnight. I typed my report on my \textbf{laptop}, saved it on a \textbf{USB flash drive} as a backup, then \textbf{uploaded} the PDF file to the platform. Unfortunately, the \textbf{Wi-Fi signal} in my neighbourhood (Etoudi) was very \textbf{weak}, so the \textbf{upload} failed twice. I had to \textbf{troubleshoot} the connection: I \textbf{restarted} my \textbf{router}, moved closer to the \textbf{access point}, and finally \textbf{logged in} again. The file \textbf{went through} at 11:45 PM. What a relief!
\end{textebox}

\textbf{Questions :}
\begin{enumerate}
    \item Relevez dans le texte 5 noms (hardware/software) et 5 verbes (actions techniques) liés aux TIC. Classez-les dans un tableau.
    \item Expliquez en anglais la différence entre \eng{to upload} et \eng{to download}. Donnez un exemple pour chacun.
    \item Le texte mentionne \eng{blended learning}. Définissez ce terme en français.
    \item Traduisez la phrase : \eng{“I had to troubleshoot the connection: I restarted my router...”} en français naturel.
\end{enumerate}

\tcblower
\textbf{Solution détaillée :}

\textbf{1. Classification lexicale :}
\begin{center}
\begin{tabularx}{\linewidth}{|X|X|}
\hline
\rowcolor{popblueL} \textbf{Devices / Hardware / Platforms (Noms)} & \textbf{Actions / Processes (Verbes + particules)} \\
\hline
laptop / USB flash drive / router / access point / LMS (Learning Management System) / Wi-Fi signal / platform / PDF file & to submit / to save (as a backup) / to upload / to troubleshoot / to restart / to log in / to go through (phrasal verb : passer, réussir) / to adopt / to type \\
\hline
\end{tabularx}
\end{center}
\textit{Note :} \eng{Blended learning} est une expression nominale (approche pédagogique). \eng{Weak/strong} sont des adjectifs qualifiant le signal.

\textbf{2. Différence \eng{upload} vs \eng{download} :}
\begin{itemize}
    \item \eng{\textbf{To upload} :} Transférer un fichier \textbf{depuis} son appareil local \textbf{vers} un serveur distant ou le cloud (direction : local $\to$ distant). \eng{Ex: “She uploaded her photos to Google Drive.”}
    \item \eng{\textbf{To download} :} Télécharger un fichier \textbf{depuis} un serveur distant \textbf{vers} son appareil local (direction : distant $\to$ local). \eng{Ex: “He downloaded the new update for his antivirus.”}
\end{itemize}
\emph{Astuce mnémotechnique :} \eng{Up} = vers le haut (le cloud/serveur) ; \eng{Down} = vers le bas (mon ordi/téléphone).

\textbf{3. Définition de \eng{blended learning} :}
C'est une \textbf{approche hybride} combinant l'enseignement en présentiel (en classe avec le professeur) et l'apprentissage en ligne (via une plateforme numérique, ressources numériques, classes virtuelles). Au Cameroun, cela permet de continuer les cours même en cas de perturbation (grèves, intempéries, pandémie).

\textbf{4. Traduction :}
« J'ai dû \textbf{résoudre le problème de connexion} / \textbf{faire du dépannage sur la connexion} : j'ai \textbf{redémarré} ma \textbf{box internet} / mon \textbf{routeur}, je me suis rapproché du \textbf{point d'accès}, et je me suis \textbf{reconnecté} (j'ai rouvert ma session). »
\begin{itemize}
    \item \eng{Troubleshoot} (irrégulier : troubleshot, troubleshot) $\to$ dépanner, diagnostiquer/résoudre un problème technique.
    \item \eng{Restart} $\to$ redémarrer (synonyme : \eng{reboot}).
    \item \eng{Log in} $\to$ ouvrir une session / se connecter.
    \item \eng{Went through} (passé de \eng{go through}) $\to$ a abouti / a réussi / est passé.
\end{itemize}
\end{exoresolu}

\begin{methode}[Maîtriser le vocabulaire des TIC pour l'apprentissage (E-learning, Recherche, Collaboration)]
\textbf{Objectif :} Utiliser avec justesse le lexique spécifique aux environnements d'apprentissage numériques (LMS, MOOC, visioconférence, outils collaboratifs) à l'oral comme à l'écrit.

\textbf{Démarche en 3 étapes :}
\begin{enumerate}
    \item \textbf{Distinguisher les plateformes et leurs fonctions :}
    \begin{itemize}
        \item \eng{LMS / VLE (Virtual Learning Environment)} : \eng{Moodle, Google Classroom, Microsoft Teams (Education)}. Verbes : \eng{to enrol in a course, to submit an assignment, to track progress, to grade}.
        \item \eng{Video conferencing tools} : \eng{Zoom, Google Meet, Teams}. Verbes : \eng{to join a meeting, to host a webinar, to share screen, to mute/unmute, to raise hand, breakout rooms}.
        \item \eng{Collaborative suites} : \eng{Google Workspace (Docs, Slides, Sheets), Office 365}. Verbes : \eng{to co-edit in real time, to comment, to suggest edits, to share a link with edit/view rights}.
    \end{itemize}
    \item \textbf{Maîtriser les noms composés et acronymes clés :}
    \eng{E-learning, m-learning (mobile), MOOC (Massive Open Online Course), SPOC, OER (Open Educational Resources), PDF, URL, hyperlink, bandwidth, deadline, timestamp, plagiarism, citation, reference manager (Zotero, Mendeley)}.
    \item \textbf{Éviter les faux-amis et calques :}
    \begin{itemize}
        \item \eng{Assiduity} (noun) $\neq$ \eng{Assiduité} (rarely used). Use \eng{attendance} (presence) or \eng{regular participation}.
        \item \eng{To assist a class} $\to$ \textbf{WRONG} (assist = help). Say \eng{to attend a class / a lecture / a webinar}.
        \item \eng{To follow a course} $\to$ possible but \eng{to take a course} or \eng{to enrol in a course} is more precise for "suivre un cours" (enrol and complete).
        \item \eng{Homework} (uncountable) $\neq$ \eng{Homeworks}. \eng{An assignment, a task, a piece of homework}.
    \end{itemize}
\end{enumerate}
\end{methode}

\begin{exoresolu}[Rédaction d'un mail formel pour un problème technique sur une plateforme d'e-learning]
\textbf{Énoncé :} Vous êtes Paul, élève en Terminale A au Government Bilingual High School (GBHS) de Bamenda. Vous n'avez pas pu rendre votre devoir d'anglais (\eng{English assignment}) sur la plateforme \eng{Moodle} hier soir à cause d'une erreur \eng{"404 Page Not Found"} puis d'une coupure d'électricité (\eng{power outage}). Rédigez un mail formel (120-150 mots) à votre professeur, M. Nfor, pour expliquer la situation, joindre le devoir en pièce jointe et demander une dérogation. Utilisez au moins 8 mots/expressions du lexique TIC "Apprentissage" vus en cours.

\tcblower
\textbf{Solution détaillée :}

\textbf{Analyse de la tâche :}
\begin{itemize}
    \item \textbf{Format :} Email formel (objet, formules de politesse, structure claire).
    \item \textbf{Ton :} Respectueux, explicatif, pas plaintif.
    \item \textbf{Contraintes lexicales :} 8 termes TIC minimum (surlignés ci-dessous).
\end{itemize}

\textbf{Proposition de rédaction :}

\begin{textebox}[Email formel : Soumission tardive de devoir]
\textbf{To:} nfor.english@gbhs-bamenda.edu.cm \\
\textbf{Subject:} Late submission – English Assignment (Terminale A) – Technical Issues

Dear Mr Nfor,

I am writing to \textbf{submit} my English \textbf{assignment} (Essay on "The Impact of Social Media") which was due yesterday on the school \textbf{LMS} (Moodle).

Unfortunately, I encountered major \textbf{technical issues}. Around 8:00 PM, when I tried to \textbf{upload} the file, the platform returned a \textbf{"404 Error"} page. I \textbf{refreshed} the \textbf{browser} and \textbf{cleared the cache}, but the \textbf{dashboard} remained inaccessible. Shortly after, a sudden \textbf{power outage} cut off my \textbf{internet connection} and I could not \textbf{log back in} before the \textbf{deadline}.

I have \textbf{attached} the completed \textbf{PDF file} to this email. Could you please accept it despite the delay? I can also \textbf{share} it via \textbf{Google Drive} if the attachment is too large.

Thank you for your understanding.

Yours sincerely, \\
Paul Atanga \\
Terminale A -- GBHS Bamenda \\
\textbf{Attachment:} Atanga\_Paul\_English\_Essay.pdf
\end{textebox}

\textbf{Vocabulaire TIC utilisé (11 termes) :} \cle{submit}, \cle{assignment}, \cle{LMS}, \cle{upload}, \cle{error / 404 Error}, \cle{browser}, \cle{cache}, \cle{dashboard}, \cle{power outage}, \cle{internet connection}, \cle{log in / log back in}, \cle{deadline}, \cle{attach / attached}, \cle{PDF file}, \cle{share}, \cle{Google Drive}.

\textbf{Points de grammaire vérifiés :}
\begin{itemize}
    \item \eng{Subject} line concise (nominal style).
    \item \eng{I am writing to + V} (structure formelle standard).
    \item \eng{Encountered} (prétérit), \eng{tried to upload} (infinitif de tentative), \eng{could not} (incapacité passée).
    \item \eng{Attached} (participe passé adjectival : "I have attached" = présent perfect pour action passée avec conséquence présente).
    \item \eng{Could you please accept...} (politesse modale).
\end{itemize}
\end{exoresolu}

\begin{methode}[Maîtriser le vocabulaire des TIC pour le monde professionnel (Télétravail, Productivité, Sécurité)]
\textbf{Objectif :} Employer le lexique du travail numérique hybride/à distance (\eng{remote work, hybrid setup, productivity tools, cybersecurity}) dans des contextes professionnels simulés (entretien, réunion, rapport).

\textbf{Démarche en 3 étapes :}
\begin{enumerate}
    \item \textbf{Construire le triptyque : Matériel / Logiciel / Soft Skills numériques.}
    \begin{center}
    \begin{tabular}{|l|l|l|}
    \hline
    \rowcolor{popblueL} \textbf{Hardware / Setup} & \textbf{Software / SaaS (Tools)} & \textbf{Compétences / Actions (Verbes)} \\
    \hline
    \eng{workstation, dual monitor, docking station, ergonomic chair, webcam, noise-cancelling headset, VPN router, UPS (onduleur)} & \eng{CRM, ERP, Slack/Teams (instant messaging), Asana/Trello (project mgmt), cloud storage (OneDrive, Dropbox), VPN client, antivirus, firewall, 2FA (Two-Factor Authentication)} & \eng{to onboard (intégrer), to clock in/out, to sync files, to collaborate remotely, to troubleshoot, to encrypt, to back up data, to comply with (se conformer à), to flag an issue} \\
    \hline
    \end{tabular}
    \end{center}
    \item \textbf{Maîtriser les collocations "Business English" :}
    \eng{Stable internet connection, bandwidth-heavy tasks, screen real estate (espace écran), action items, deliverables, KPIs (Key Performance Indicators), SLA (Service Level Agreement), downtime, uptime, remote access, virtual private network}.
    \item \textbf{Attention aux faux-amis "Bureau/Office" :}
    \begin{itemize}
        \item \eng{Office} = Bureau (lieu/pièce) OU Suite logicielle (Microsoft Office). Pas "Officice".
        \item \eng{Desk} = Bureau (meuble). \eng{Desk job} = travail de bureau.
        \item \eng{Remote work / Work from home (WFH)} = Télétravail. \eng{Hybrid work} = Hybride.
        \item \eng{To work on a file} (travailler sur un fichier) vs \eng{To work out} (résoudre / faire du sport).
    \end{itemize}
\end{enumerate}
\end{methode}

\begin{exoresolu}[Compréhension orale / Lecture : Témoignage d'un jeune développeur à Douala]
\textbf{Énoncé :} Lisez le script d'une interview (simulée) de Christian, développeur Full Stack junior dans une startup fintech à Douala (Akwa). Répondez aux questions en anglais.

\begin{textebox}[Interview : "A Day in the Life of a Junior Dev in Douala"]
\textbf{Interviewer:} Christian, describe your typical \textbf{remote work} setup.
\textbf{Christian:} Sure. I work from a \textbf{co-working space} in Akwa most days because the \textbf{internet bandwidth} is more reliable than at home in Bonabéri. I have a \textbf{dual-monitor setup} connected to my \textbf{laptop} via a \textbf{docking station}. A good \textbf{noise-cancelling headset} is essential for our daily \textbf{stand-up meetings} on \textbf{Slack} huddles or \textbf{Google Meet}. We use \textbf{Jira} for \textbf{task tracking} and \textbf{GitHub} for \textbf{version control}. Security is strict: \textbf{2FA} is mandatory for \textbf{VPN access} to the \textbf{production server}, and we must \textbf{encrypt} sensitive data.

\textbf{Interviewer:} What about \textbf{productivity} challenges?
\textbf{Christian:} Power cuts (\textbf{load shedding}) are the main issue. The co-working space has a \textbf{generator} and \textbf{UPS units}, so zero \textbf{downtime}. But when I work from home, I rely on my phone's \textbf{hotspot} as a backup. Time management is key: I use the \textbf{Pomodoro technique} and block \textbf{deep work} slots in my \textbf{Google Calendar}. Communication must be \textbf{asynchronous} sometimes because our \textbf{CTO} is based in Lagos (different time zone, but only 1 hour gap).
\end{textebox}

\textbf{Questions :}
\begin{enumerate}
    \item Why does Christian prefer the co-working space over working from home? (2 reasons)
    \item List 5 specific software tools/platforms mentioned and their function (e.g., Slack = communication).
    \item Explain the meaning of \eng{"zero downtime"} and \eng{"load shedding"} in this context.
    \item Translate into French: \eng{"We use Jira for task tracking and GitHub for version control."}
    \item Find in the text the English equivalents for: \eng{Onduleur}, \eng{Partage de connexion (téléphone)}, \eng{Travail en profondeur / concentration}, \eng{Communication asynchrone}, \eng{Double authentification}.
\end{enumerate}

\tcblower
\textbf{Solution détaillée :}

\textbf{1. Raisons du choix du co-working space :}
\begin{itemize}
    \item \textbf{Reliable internet bandwidth} (bande passante fiable / connexion stable) — crucial pour un dev.
    \item \textbf{Infrastructure de secours :} \eng{Generator} (groupe électrogène) et \eng{UPS units} (onduleurs) garantissant \eng{zero downtime} (zéro temps d'arrêt / aucune coupure de service).
\end{itemize}

\textbf{2. Outils logiciels et fonctions :}
\begin{center}
\begin{tabularx}{\linewidth}{|X|X|}
\hline
\rowcolor{popblueL} \textbf{Outil / Plateforme} & \textbf{Fonction} \\
\hline
Slack (huddles) & Messagerie instantanée / Appels audio courts quotidiens (\eng{stand-up meetings}) \\
Google Meet & Visioconférence \\
Jira & Suivi des tâches / Gestion de projet (\eng{task tracking}) \\
GitHub & Contrôle de version / Dépôt de code (\eng{version control}) \\
VPN Client & Accès distant sécurisé au réseau de l'entreprise (\eng{VPN access}) \\
Google Calendar & Gestion du temps / Planification (\eng{block deep work slots}) \\
\hline
\end{tabularx}
\end{center}

\textbf{3. Définitions contextuelles :}
\begin{itemize}
    \item \eng{\textbf{Zero downtime} :} Disponibilité continue du système (100\% \eng{uptime}). Aucune interruption de service due à l'électricité ou au réseau. Ici, grâce aux onduleurs et générateurs.
    \item \eng{\textbf{Load shedding} :} \textbf{Délestage électrique} (coupures programmées ou tournantes par l'opérateur Eneo pour gérer la demande). Terme très courant en Afrique du Sud et au Cameroun.
\end{itemize}

\textbf{4. Traduction :}
« Nous utilisons \textbf{Jira} pour le \textbf{suivi des tâches} et \textbf{GitHub} pour le \textbf{contrôle de version} (gestion des versions du code). »

\textbf{5. Équivalents anglais trouvés dans le texte :}
\begin{itemize}
    \item Onduleur $\to$ \eng{\textbf{UPS (unit)} / \textbf{UPS units}} (Uninterruptible Power Supply).
    \item Partage de connexion (téléphone) $\to$ \eng{\textbf{hotspot}} (ou \eng{personal hotspot}).
    \item Travail en profondeur / concentration $\to$ \eng{\textbf{deep work}} (terme de Cal Newport, très utilisé en tech).
    \item Communication asynchrone $\to$ \eng{\textbf{asynchronous}} (communication) / \eng{async} (argot tech).
    \item Double authentification $\to$ \eng{\textbf{2FA}} (Two-Factor Authentication).
\end{itemize}
\end{exoresolu}

\begin{methode}[Collocations, familles de mots et faux-amis : Les pièges du Bac]
\textbf{Objectif :} Sécuriser l'orthographe, la dérivation et les associations de mots (collocations) pour ne pas perdre de points bêtes à l'écrit (composition, traduction, gap-fill).

\textbf{Démarche en 3 axes :}

\textbf{Axe 1 : Familles de mots (Word Formation) — Tableau de dérivation à connaître par cœur}
\begin{center}
\begin{tabularx}{\linewidth}{|X|X|X|X|}
\hline
\rowcolor{popblueL} \textbf{Verbe} & \textbf{Nom (Action/Concept)} & \textbf{Nom (Agent/Appareil)} & \textbf{Adjectif / Participe} \\
\hline
\eng{to connect} & \eng{connection / connectivity} & \eng{connector} & \eng{connected / wireless} \\
\eng{to compute} & \eng{computation / computing} & \eng{computer} & \eng{computational} \\
\eng{to program} & \eng{programming / program (US) / programme (UK)} & \eng{programmer} & \eng{programmable} \\
\eng{to develop} & \eng{development} & \eng{developer} & \eng{developmental} \\
\eng{to secure} & \eng{security} & \eng{—} & \eng{secure / insecure} \\
\eng{to encrypt} & \eng{encryption} & \eng{—} & \eng{encrypted} \\
\eng{to authenticate} & \eng{authentication} & \eng{—} & \eng{authenticated} \\
\eng{to back up} & \eng{a backup (n.m.)} & \eng{—} & \eng{backed up} \\
\hline
\end{tabularx}
\end{center}
\emph{Attention orthographe :} \eng{Programming} (double M UK/US), \eng{Connection} (pas *connexion), \eng{Developer} (pas *developper).

\textbf{Axe 2 : Collocations fortes (mariages de mots obligatoires)}
\begin{itemize}
    \item \eng{\textbf{High-speed} broadband / connection} (pas *fast speed).
    \item \eng{\textbf{Surf / Browse} the web / the internet} (pas *visit the internet).
    \item \eng{\textbf{Run} a program / an app / a script} (pas *make run).
    \item \eng{\textbf{Crash / Freeze / Hang} (verbes intransitifs pour l'ordi/logiciel)}. \eng{My laptop crashed.}
    \item \eng{\textbf{Install / Uninstall / Update / Upgrade} software}.
    \item \eng{\textbf{Grant / Deny} access / permission}.
    \item \eng{\textbf{Flag / Report / Fix / Resolve} a bug / an issue / an error}.
\end{itemize}

\textbf{Axe 3 : Top 5 Faux-amis / Calques "Français $\to$ Anglais" interdits}
\begin{center}
\begin{tabularx}{\linewidth}{|X|X|X|}
\hline
\rowcolor{poporangeL} \textbf{Français (Piège)} & \textbf{Anglais INCORRECT (Calque)} & \textbf{Anglais CORRECT} \\
\hline
Actuellement / En ce moment & \eng{\textbf{Actually}} (FAUX = en fait) & \eng{\textbf{Currently / At the moment / Nowadays}} \\
Une formation (cours) & \eng{A formation} (FAUX = géologie/militaire) & \eng{A course / Training / A workshop / A program} \\
Un stage (entreprise) & \eng{A stage} (FAUX = scène/théâtre) & \eng{An internship (US) / A work placement (UK)} \\
Un délai (date limite) & \eng{A delay} (FAUX = retard) & \eng{A deadline / A time limit} \\
Assister à (un cours/réunion) & \eng{To assist} (FAUX = aider) & \eng{To attend / To take part in} \\
\hline
\end{tabularx}
\end{center}
\end{methode}

\begin{exoresolu}[Exercice de transformation lexicale (Gap-fill / Word Formation / Correction d'erreurs)]
\textbf{Énoncé :} Complétez les phrases avec la forme correcte du mot entre parenthèses (dérivation), choisissez la bonne collocation ou corrigez l'erreur soulignée.

\textbf{Partie A : Word Formation (Dérivation)}
\begin{enumerate}
    \item The \_\_\_\_\_\_\_\_\_\_ (develop) of the new mobile \_\_\_\_\_\_\_\_\_\_ (apply) took six months.
    \item You need a strong \_\_\_\_\_\_\_\_\_\_ (authenticate) method, like 2FA, to \_\_\_\_\_\_\_\_\_\_ (secure) the server.
    \item The \_\_\_\_\_\_\_\_\_\_ (program) wrote a script to automate the \_\_\_\_\_\_\_\_\_\_ (back up) process.
    \item Poor \_\_\_\_\_\_\_\_\_\_ (connect) in rural areas limits access to \_\_\_\_\_\_\_\_\_\_ (educate) resources.
\end{enumerate}

\textbf{Partie B : Collocations (Choix multiple)}
\begin{enumerate}
    \setcounter{enumi}{4}
    \item The system \_\_\_\_\_\_\_\_\_\_ suddenly, so I lost my unsaved data. (a) crashed (b) fell (c) broke (d) stopped.
    \item We need to \_\_\_\_\_\_\_\_\_\_ the latest security patch immediately. (a) make (b) install (c) put (d) place.
    \item The IT manager \_\_\_\_\_\_\_\_\_\_ me access to the shared drive. (a) gave (b) \textbf{granted} (c) offered (d) provided.
    \item Don't forget to \_\_\_\_\_\_\_\_\_\_ your files to the cloud before the \_\_\_\_\_\_\_\_\_\_. (a) upload / deadline (b) download / delay (c) charge / date.
\end{enumerate}

\textbf{Partie C : Correction d'erreurs (Faux-amis / Calques)}
\begin{enumerate}
    \setcounter{enumi}{8}
    \item "I am doing a \textbf{formation} on cybersecurity next week." $\to$ Corriger.
    \item "The \textbf{delay} for the submission is Friday midnight." $\to$ Corriger.
    \item "Actually, I \textbf{assist} the online lecture every Tuesday." $\to$ Corriger (2 erreurs).
    \item "He is a good \textbf{informaticien}." $\to$ Corriger (métier).
\end{enumerate}

\tcblower
\textbf{Solution détaillée :}

\textbf{Partie A : Word Formation}
\begin{enumerate}
    \item \textbf{development} (nom, suffixe \textit{-ment}) / \textbf{application} ou \textbf{app} (nom, suffixe \textit{-tion}). \textit{Note : "App" est l'abréviation standard acceptée.}
    \item \textbf{authentication} (nom) / \textbf{secure} (verbe infinitif après \eng{to}). \textit{Piège : pas "securize" (franglais), pas "security" (nom).}
    \item \textbf{programmer} (agent, suffixe \textit{-er}) / \textbf{backup} (nom composé, s'écrit en un mot : \eng{a backup}). \textit{Verbe = two words \eng{back up}, Nom/Adj = one word \eng{backup}.}
    \item \textbf{connectivity} (nom qualité/état, plus précis que \eng{connection} pour l'infrastructure) / \textbf{educational} (adj) ou \textbf{education} (nom non comptable). \eng{Educational resources} (collocation standard).
\end{enumerate}

\textbf{Partie B : Collocations}
\begin{enumerate}
    \setcounter{enumi}{4}
    \item \textbf{(a) crashed.} \eng{Crash} = planter (système/logiciel). \eng{Freeze} = geler (écran figé). \eng{Fell/Broke} = physique.
    \item \textbf{(b) install.} \eng{Install software / an update / a patch}. \eng{Make/put/place} impossibles.
    \item \textbf{(b) granted.} \eng{Grant access / permission / a licence} (verbe formel/administratif). \eng{Gave} possible mais moins précis/professionnel.
    \item \textbf{(a) upload / deadline.} \eng{Upload} (vers cloud). \eng{Deadline} (date limite). \eng{Delay} = retard. \eng{Charge} = recharger (batterie).
\end{enumerate}

\textbf{Partie C : Correction d'erreurs (Explications grammaticales)}
\begin{enumerate}
    \setcounter{enumi}{8}
    \item \textbf{Erreur :} \eng{formation} (faux ami). \textbf{Correction :} "I am taking a \textbf{course} / \textbf{training} / \textbf{workshop} on cybersecurity next week."
    \item \textbf{Erreur :} \eng{delay} (faux ami = retard). \textbf{Correction :} "The \textbf{deadline} for submission is Friday midnight." (\eng{Submission deadline}).
    \item \textbf{Erreurs :} \eng{Actually} (faux ami = en fait) + \eng{assist} (faux ami = aider). \textbf{Correction :} "\textbf{Currently} / \textbf{At the moment}, I \textbf{attend} / \textbf{follow} the online lecture every Tuesday."
    \item \textbf{Erreur :} \eng{informaticien} (français). \textbf{Correction :} "He is a good \textbf{IT specialist} / \textbf{software engineer} / \textbf{developer} / \textbf{computer scientist} / \textbf{tech expert}." \textit{Il n'existe pas de mot unique "informatician" en anglais courant.}
\end{enumerate}
\end{exoresolu}

\begin{methode}[Réinvestir le lexique TIC à l'oral (Speaking) : Décrire, Dépanner, Débattre]
\textbf{Objectif :} Produire un discours cohérent, fluide et lexicalement riche lors d'activités orales (exposé, jeu de rôle, débat, description d'image) sur le thème des TIC.

\textbf{Démarche en 3 situations types :}

\textbf{Situation 1 : Décrire une configuration / un environnement (Photo / Schéma / Personnel)}
\begin{itemize}
    \item \textbf{Structure :} \eng{"In this picture / In my setup, I can see..."} $\to$ \textbf{Hardware} (\eng{dual monitor, ergonomic chair, webcam, router}) $\to$ \textbf{Logiciel visible/implicite} (\eng{running Windows 11, using VS Code, connected via VPN}) $\to$ \textbf{Contexte/Usage} (\eng{for remote work / for coding / for online classes}).
    \item \textbf{Connecteurs de description spatiale :} \eng{On the left/right, in the background, next to, connected to, plugged into}.
    \item \textbf{Adjectifs utiles :} \eng{Cluttered / tidy / minimalist / ergonomic / high-spec / outdated / laggy}.
\end{itemize}

\textbf{Situation 2 : Jeu de rôle "Support Technique / Helpdesk" (Problème $\to$ Solution)}
\begin{itemize}
    \item \textbf{Script type :}
    \begin{enumerate}
        \item \textbf{User :} \eng{"Hi, I'm having trouble with... / I can't seem to... / There's an issue with..."} (\eng{trouble with + V-ing / noun}).
        \item \textbf{Agent :} \eng{"Let me help you troubleshoot. First, have you tried...?" / "Can you check if...?" / "Please try to..."}
        \item \textbf{Actions guidées :} \eng{restart / reboot, unplug / plug back in, update the driver, clear cache / cookies, check cables, run the troubleshooter, disable / enable the adapter}.
        \item \textbf{Résolution :} \eng{"It's working now!" / "The issue persists." / "I'll escalate a ticket."}
    \end{enumerate}
\end{itemize}

\textbf{Situation 3 : Débat / Discussion "Pour ou Contre" (Ex: \eng{Remote work vs Office}, \eng{AI in education}, \eng{Phone ban in schools})}
\begin{itemize}
    \item \textbf{Vocabulaire d'opinion :} \eng{In my view, the main advantage/drawback is... / On the one hand... on the other hand... / However, we must consider... / For instance... / Statistics show that...}
    \item \textbf{Mots-clés arguments :} \eng{Productivity, work-life balance, commute, distraction, collaboration, isolation, cybersecurity risks, digital divide (fracture numérique), screen time, soft skills, critical thinking}.
\end{itemize}
\end{methode}

\begin{exoresolu}[Simulation orale : Jeu de rôle "Helpdesk Étudiant" + Grille d'auto-évaluation]
\textbf{Énoncé :} En binôme. L'élève A a un problème technique pour rendre un devoir sur Moodle. L'élève B est le support technique de l'école. Jouez le dialogue (2-3 min). Utilisez le maximum de vocabulaire ciblé.

\textbf{Rôle A (Étudiant - Problème) :} Vous êtes à la maison (quartier Makepe, Douala). Connexion 4G instable. Erreur \eng{"File size exceeds limit"} puis \eng{"Session expired"}. Vous êtes stressé, la deadline est dans 30 min.
\textbf{Rôle B (Support - Solution) :} Calme, procédure pas à pas. Vérifier taille fichier $\to$ compresser (zip) / convertir en PDF $\to$ vérifier connexion $\to$ vider cache navigateur $\to$ relancer session $\to$ soumettre.

\tcblower
\textbf{Proposition de dialogue modèle (Script) :}

\begin{textebox}[Role-play: IT Support / Student]
\textbf{Student:} Hello, this is Grace from Terminale A. I'm \textbf{struggling to submit} my History assignment on the \textbf{LMS}. I keep getting an error message: "\textbf{File size exceeds limit}".
\textbf{Support:} Hello Grace. Don't panic, we have time. What \textbf{file format} is it?
\textbf{Student:} It's a \textbf{Word document} (\textbf{.docx}), about 25 MB with many images.
\textbf{Support:} The \textbf{upload limit} is 10 MB. You need to \textbf{compress} it. Right-click the file, choose "\textbf{Send to} then \textbf{Compressed (zipped) folder}". Or better, \textbf{export / save as PDF} — it's usually lighter and \textbf{non-editable}.
\textbf{Student:} OK, I'll \textbf{convert to PDF}... Done. It's 3 MB now. But now I get "\textbf{Session expired}".
\textbf{Support:} Your \textbf{session timed out} due to inactivity. \textbf{Log in} again. Also, your 4G seems \textbf{unstable}. Try to \textbf{clear your browser cache} (Ctrl+Shift+Del) and \textbf{disable any VPN} if active.
\textbf{Student:} Logged in again. Cache cleared. \textbf{Uploading} now... \textbf{Success!} "File submitted at 23:45". Thank you so much!
\textbf{Support:} You're welcome. Remember to \textbf{keep a local backup} next time. Good luck!
\end{textebox}

\textbf{Grille d'auto-évaluation lexicale (à cocher après l'oral) :}
\begin{center}
\begin{tabular}{|l|c|c|}
\hline
\rowcolor{popgreenL} \textbf{Notion / Terme utilisé} & \textbf{Élève A} & \textbf{Élève B} \\
\hline
submit / submission / assignment / LMS / deadline & $\square$ & $\square$ \\
file size / limit / exceeds / format (.docx, .pdf) / MB & $\square$ & $\square$ \\
compress / zip / convert / export / save as & $\square$ & $\square$ \\
session expired / timed out / log in / cache / VPN / unstable & $\square$ & $\square$ \\
upload / success / backup / troubleshoot & $\square$ & $\square$ \\
\hline
\textbf{Connecteurs / Fluidité} (First, Then, Actually, The issue is..., Try to...) & $\square$ & $\square$ \\
\hline
\end{tabular}
\end{center}
\textbf{Conseil prof :} Enregistrez-vous (WhatsApp vocal) et réécoutez-vous pour vérifier la prononciation de \eng{exceeds} [ik-siidz], \eng{unstable} [un-stei-bul], \eng{cache} [kæʃ "cach"], \eng{VPN} [vi-pi-en].
\end{exoresolu}

\begin{savaistu}[Le Cameroun et la révolution numérique : "Silicon Mountain" et la fracture digitale]
Le Cameroun est un acteur majeur de la tech en Afrique centrale. \textbf{Buea} (région du Sud-Ouest) abrite le \textbf{"Silicon Mountain"}, un écosystème de startups (ex. \eng{Churchill Nanje} de \eng{Njorku}, \eng{Agro-Hub}) soutenu par des hubs comme \eng{ActivSpaces}. Le gouvernement a lancé \eng{"Cameroun Numérique 2020"} pour étendre la fibre optique et la 4G/5G. Pourtant, la \textbf{fracture numérique} (\eng{digital divide}) persiste : coût de la data, coupures d'électricité (\eng{load shedding}), couverture rurale insuffisante. L'anglais est la langue de codage et de documentation technique mondiale : maîtriser ce vocabulaire, c'est s'ouvrir des carrières de \eng{developer, data analyst, cybersecurity expert} au pays et à l'international.
\end{savaistu}

\begin{experience}[Projet de classe : "Mon Portfolio Numérique"]
\textbf{Objectif :} Créer une page web simple (HTML/CSS) ou un site \eng{Google Sites} / \eng{WordPress} présentant votre parcours, vos compétences TIC et un projet scolaire.
\textbf{Étapes :}
\begin{enumerate}
    \item \textbf{Planifier :} Arborescence (Accueil, À propos, Projets, Contact). Lexique : \eng{navigation bar, header, footer, hyperlink, responsive design}.
    \item \textbf{Produire :} Rédiger les contenus en anglais (bio, description de projets). Utiliser \eng{Canva} pour visuels.
    \item \textbf{Publier :} Héberger (\eng{GitHub Pages, Netlify, Google Sites}). Tester sur mobile (\eng{responsive}).
    \item \textbf{Présenter :} Soutenance orale de 3 min devant la classe (\eng{pitch}).
\end{enumerate}
\textbf{Évaluation :} Grille critériée (Vocabulaire TIC 30\%, Grammaire 20\%, Design/UX 20\%, Oral 30\%).
\end{experience}

\begin{attention}[Les 5 erreurs qui coûtent des points au Bac (Spécial TIC Vocabulaire)]
\begin{enumerate}
    \item \textbf{Faux-ami \eng{Actually} / \eng{Actuellement} :} JAMAIS écrire \eng{"Actually, I study online"} pour "Actuellement, j'étudie en ligne".
    \begin{center}
    \eng{Actually = En fait / En réalité (correction).} \quad \eng{Currently / At present / Nowadays = Actuellement.}
    \end{center}
    \item \textbf{Confusion \eng{Delay} / \eng{Deadline} :} \eng{Delay} = RETARD (négatif, imprévu). \eng{Deadline} = DATE LIMITE (fixée à l'avance).
    \begin{center}
    \eng{Wrong: "The delay is tomorrow."} \quad \eng{Right: "The \textbf{deadline} is tomorrow."}
    \end{center}
    \item \textbf{Calque \eng{Assist a class / meeting} :} \eng{Assist} = AIDER (help). Pour "assister à" (être présent) : \eng{ATTEND a class / a meeting / a webinar / a lecture}.
    \item \textbf{Ordre des mots / Particules (Phrasal Verbs) :} Séparer le verbe de sa particule avec un pronom.
    \begin{center}
    \eng{Wrong: "Turn on it."} \quad \eng{Right: "Turn \textbf{it} on."} (Pronom $\to$ AU MILIEU). \\
    \eng{Right: "Upload \textbf{the large PDF file} to the cloud." / "Upload \textbf{it} to the cloud."}
    \end{center}
    \item \textbf{Orthographe \eng{Connection} vs \eng{Connexion} :} En anglais, s'écrit \textbf{CONNECTION} (double N, TION). \eng{Connexion} n'existe pas. Idem : \eng{Programming} (double M), \eng{Developer} (simple P, double E), \eng{Keyboard} (pas *key board), \eng{Online} (standard, trait d'union possible \eng{on-line} mais déconseillé), \eng{Website} (un mot).
\end{enumerate}
\end{attention}

\newpage

\coursec{Exercices}

\courssub{Je vérifie mes connaissances}

\begin{exercice}[QCM : Vocabulaire de base des TIC]
Choisissez la bonne réponse parmi les quatre propositions.
\begin{enumerate}
\item The teacher asked the students to \eng{log in} to the school platform. What does \eng{log in} mean?
      \begin{itemize}
      \item[a)] To turn off the computer
      \item[b)] To enter a username and password to access a system
      \item[c)] To save a document on a USB drive
      \item[d)] To print a document
      \end{itemize}
\item Which word is a \textbf{synonym} of \eng{attachment} in an email context?
      \begin{itemize}
      \item[a)] Subject
      \item[b)] File added to the message
      \item[c)] Spam
      \item[d)] Draft
      \end{itemize}
\item A \eng{webinar} is:
      \begin{itemize}
      \item[a)] A virus that attacks websites
      \item[b)] A seminar conducted over the Internet
      \item[c)] A type of web browser
      \item[d)] A wireless connection
      \end{itemize}
\item If a website is \eng{user-friendly}, it is:
      \begin{itemize}
      \item[a)] Difficult to navigate
      \item[b)] Only for experts
      \item[c)] Easy to use and understand
      \item[d)] Not secure
      \end{itemize}
\item To \eng{back up} your files means to:
      \begin{itemize}
      \item[a)] Delete them permanently
      \item[b)] Make a copy in case the original is lost
      \item[c)] Send them by email
      \item[d)] Compress them into a zip folder
      \end{itemize}
\end{enumerate}
\end{exercice}

\begin{exercice}[Vrai ou Faux justifié : Faux amis et sens]
Décidez si chaque affirmation est vraie (True) ou fausse (False). Justifiez votre réponse en français en donnant le sens correct du mot souligné.
\begin{enumerate}
\item \eng{To assist a meeting} means to attend it. \hfill (True / False)
\item \eng{Actually}, I prefer working from home. \eng{Actually} means « actuellement ». \hfill (True / False)
\item The \eng{library} in a programming context is a place where you borrow books. \hfill (True / False)
\item \eng{To download} a file means to send it from your computer to a server. \hfill (True / False)
\item A \eng{smartphone} is a \eng{portable} computer. \hfill (True / False)
\end{enumerate}
\end{exercice}

\begin{exercice}[Vocabulaire : Familles de mots et collocations]
Complétez le tableau en donnant la forme manquante (nom, verbe, adjectif ou adverbe) et une collocation fréquente (verbe + nom ou adjectif + nom). Le premier est fait en exemple.
\begin{center}
\begin{tabularx}{\linewidth}{|X|X|X|X|}
\hline
\rowcolor{popblueL} \textbf{Verbe} & \textbf{Nom} & \textbf{Adjectif / Adverbe} & \textbf{Collocation} \\
\hline
to connect & connection / connectivity & connected / connectedly & \eng{a stable connection} \\
\hline
to inform &  &  &  \\
\hline
 & computation / computer &  &  \\
\hline
to digitalize / digitize &  & digital / digitally &  \\
\hline
to browse &  &  &  \\
\hline
to upload &  &  &  \\
\hline
to stream &  &  &  \\
\hline
\end{tabularx}
\end{center}
\end{exercice}

\begin{exercice}[Grammaire en contexte : Prépositions et particules]
Complétez les phrases avec la préposition ou la particule appropriée (\eng{in, on, up, down, into, out, off, for, to, with}).
\begin{enumerate}
\item Please log \_\_\_\_\_\_ using your student ID and password.
\item The file is too large to send \_\_\_\_\_\_ email; use a cloud link instead.
\item She is very good \_\_\_\_\_\_ troubleshooting network issues.
\item We need to back \_\_\_\_\_\_ the database before the update.
\item The system crashed, so we had to start \_\_\_\_\_\_ all over again.
\item Click \_\_\_\_\_\_ the icon to open the application.
\item Are you connected \_\_\_\_\_\_ the Wi-Fi network?
\item He spends hours browsing \_\_\_\_\_\_ the web for research.
\item The software is compatible \_\_\_\_\_\_ both Windows and Linux.
\item Don't forget to sign \_\_\_\_\_\_ when you leave the computer lab.
\end{enumerate}
\end{exercice}

\courssub{Je m'entraîne}

\begin{exercice}[Traduction : Éviter les calques du français]
Traduisez les phrases suivantes en anglais naturel. Attention aux faux amis et à la structure verbale.
\begin{enumerate}
\item J'assiste à une vidéoconférence avec mes collègues de Douala.
\item Actuellement, je télécharge le rapport sur le serveur.
\item Il faut sauvegarder tes données sur un disque dur externe.
\item Cette application est très conviviale, même pour les débutants.
\item Nous naviguons sur Internet pour trouver des informations fiables.
\item Le formateur a partagé son écran pendant le webinaire.
\item Je vais mettre à jour mon logiciel antivirus ce soir.
\item Le fichier joint est un PDF de 5 Mo.
\item Elle a imprimé le document sans le vérifier.
\item Notre connexion Internet est lente aujourd'hui.
\end{enumerate}
\end{exercice}

\begin{exercice}[Transformation : Reformulation avec le vocabulaire cible]
Réécrivez chaque phrase en utilisant le mot ou l'expression entre parenthèses sans en changer le sens.
\begin{enumerate}
\item She sent the document attached to the email. (\eng{attachment})
\item The online seminar starts at 3 p.m. (\eng{webinar})
\item You must enter your password to access the platform. (\eng{log in})
\item Make a copy of your project on Google Drive. (\eng{back up})
\item The interface is easy to understand. (\eng{user-friendly})
\item He transferred the photos from his phone to his laptop. (\eng{upload})
\item They are watching a live video on YouTube. (\eng{streaming})
\item I need to install the latest version of the browser. (\eng{update})
\item The file was too big, so she compressed it. (\eng{zip})
\item We are working from home this week. (\eng{teleworking / remote work})
\end{enumerate}
\end{exercice}

\begin{exercice}[Texte à trous : Le télétravail au Cameroun]
Complétez le texte suivant avec les mots de la banque ci-dessous. Chaque mot s'utilise une seule fois. Deux mots sont en trop.
\begin{textebox}[Gap-filling: Teleworking in Cameroon]
In major Cameroonian cities like Yaoundé and Douala, \_\_\_\_\_\_ has become more common since the pandemic. Many employees now work \_\_\_\_\_\_ from home two or three days a week. A stable Internet \_\_\_\_\_\_ is essential, but power cuts remain a challenge. Workers use video \_\_\_\_\_\_ tools like Zoom or Teams for meetings. They share files via cloud storage and send \_\_\_\_\_\_ by email. Companies provide \_\_\_\_\_\_ laptops with pre-installed security software. Employees must \_\_\_\_\_\_ their work regularly to avoid data loss. Some find it hard to \_\_\_\_\_\_ professional and personal life. However, most appreciate the flexibility and the reduction in daily \_\_\_\_\_\_ time.
\end{textebox}
\begin{center}
\begin{tabular}{|c|c|c|c|c|c|}
\hline
\eng{teleworking} & \eng{remotely} & \eng{connection} & \eng{conferencing} & \eng{attachments} & \eng{company} \\
\hline
\eng{back up} & \eng{balance} & \eng{commute} & \eng{download} & \eng{offline} & \eng{hardware} \\
\hline
\end{tabular}
\end{center}
\end{exercice}

\begin{exercice}[Appariement : Collocations verbales]
Associez chaque verbe de la colonne A à un nom de la colonne B pour former une collocation correcte. Écrivez les paires (ex. 1–f).
\begin{center}
\begin{tabular}{|c|l|c|l|}
\hline
\multicolumn{2}{|c|}{\textbf{Colonne A (Verbes)}} & \multicolumn{2}{c|}{\textbf{Colonne B (Noms)}} \\
\hline
1 & to browse & a & a file \\
2 & to download & b & the web \\
3 & to log & c & a meeting \\
4 & to attend & d & in / out \\
5 & to upload & e & a document \\
6 & to print & f & a password \\
7 & to reset & g & a screenshot \\
8 & to take & h & a program \\
9 & to install & i & a bug \\
10 & to fix & j & a website \\
\hline
\end{tabular}
\end{center}
\end{exercice}

\courssub{Je me prépare au Bac}

\begin{exercice}[Compréhension de texte et langue : Type Bac — Texte original]
Lisez attentivement le texte ci-dessous, puis répondez aux questions qui suivent.
\begin{textebox}[ICTs Transforming Education in Rural Cameroon]
In the heart of the Northwest Region, Government High School Buea has embraced digital learning. Thanks to a partnership with a local NGO, the school received fifty tablets and a solar-powered router. Students now access \eng{offline} educational resources — videos, interactive quizzes, and e-books — without needing a constant Internet \eng{connection}. Teachers \eng{upload} weekly lessons to a local server; learners \eng{download} them during computer lab sessions. “Before, we had only outdated textbooks,” says Form 5 student Nguemga. “Now I can \eng{stream} a physics experiment on my tablet.” The headteacher notes improved attendance and better results in science subjects. However, challenges persist: device \eng{maintenance}, limited \eng{bandwidth} for updates, and the need for continuous teacher \eng{training}. The Ministry of Secondary Education plans to \eng{scale up} the pilot to ten more schools by 2026.
\end{textebox}

\textbf{Comprehension questions (answer in English, using your own words as much as possible):}
\begin{enumerate}
\item What equipment did the school receive, and how is it powered?
\item How do students access the educational content without a permanent Internet connection?
\item Quote two benefits mentioned in the text resulting from the use of tablets.
\item Identify two challenges the school still faces.
\item What does the Ministry plan to do by 2026?
\end{enumerate}

\textbf{Language work:}
\begin{enumerate}
\setcounter{enumi}{5}
\item Find in the text the \textbf{noun} formed from the verb \eng{to maintain} (paragraph 2).
\item Find the \textbf{phrasal verb} meaning “to expand” or “to extend to a larger scale” (last sentence).
\item Give the \textbf{opposite} of \eng{outdated} (paragraph 2) as used in the text.
\item Rewrite the sentence “Teachers upload weekly lessons to a local server” in the \textbf{passive voice}.
\item Complete with the right \textbf{preposition}: Students are interested \_\_\_\_\_\_ digital learning. / The tablet is connected \_\_\_\_\_\_ the local server.
\end{enumerate}
\end{exercice}

\begin{exercice}[Traduction : Type Bac — Thème]
Traduisez le passage suivant en anglais. Veillez à employer le vocabulaire de la leçon (TIC, apprentissage, travail) et à soigner la structure des phrases.
\begin{quote}
« Dans les lycées camerounais, l'intégration des TIC transforme peu à peu les méthodes d'enseignement. Les tableaux interactifs remplacent les tableaux noirs traditionnels. Les élèves utilisent des tablettes pour consulter des manuels numériques et faire des exercices en ligne. Les professeurs, formés aux outils numériques, créent des contenus interactifs et organisent des classes virtuelles. Malgré les coupures d'électricité et la connexion parfois lente, l'enthousiasme reste fort. Le ministère prévoit d'équiper davantage d'établissements d'ici 2027. »
\end{quote}
\end{exercice}

\begin{exercice}[Expression écrite guidée : Type Bac — Composition]
\textbf{Sujet :} \eng{“How ICTs can improve learning in Cameroonian schools”}
Rédigez un paragraphe cohérent de \textbf{80 à 100 mots} (hors titre) en respectant les consignes suivantes :
\begin{itemize}
\item Utilisez \textbf{au moins 6 mots ou expressions} du vocabulaire de la leçon (ex. \eng{tablet, offline, download, upload, webinar, user-friendly, bandwidth, teleworking, back up, interactive, stream, scale up}).
\item Intégrez \textbf{au moins 3 connecteurs logiques} (ex. \eng{moreover, however, therefore, consequently, in addition, for instance}).
\item Structurez votre paragraphe : phrase d'introduction, arguments illustrés, conclusion.
\item Soignez l'orthographe, la grammaire et la ponctuation anglaises.
\end{itemize}
\begin{center}
\textit{Word count: \_\_\_\_\_\_ words}
\end{center}
\end{exercice}

\newpage

\coursec{Corrigés des exercices}

\begin{corrige}[Exercice 1 — QCM : Vocabulaire de base des TIC]
\begin{enumerate}
\item \textbf{b) To enter a username and password to access a system.} \\
      \eng{Log in} (ou \eng{sign in}) signifie « ouvrir une session / se connecter » en fournissant ses identifiants.
\item \textbf{b) File added to the message.} \\
      Un \eng{attachment} est une « pièce jointe » (fichier joint au courriel). \eng{Subject} = objet, \eng{Spam} = courriel indésirable, \eng{Draft} = brouillon.
\item \textbf{b) A seminar conducted over the Internet.} \\
      \eng{Webinar} = \eng{web} + \eng{seminar} (webinaire), conférence ou atelier en ligne.
\item \textbf{c) Easy to use and understand.} \\
      \eng{User-friendly} = « convivial, facile d'utilisation » (litt. « ami de l'utilisateur »).
\item \textbf{b) Make a copy in case the original is lost.} \\
      \eng{Back up} (verbe à particule) = « sauvegarder (faire une copie de secours) ». \eng{Backup} (nom) = « copie de sauvegarde ».
\end{enumerate}
\end{corrige}

\begin{corrige}[Exercice 2 — Vrai ou Faux justifié : Faux amis et sens]
\begin{enumerate}
\item \textbf{False.} \eng{To assist} signifie « aider, assister (qn) » (\textit{to help}). Pour « assister à » (être présent), on utilise \eng{to attend}. \eng{Attend a meeting} = assister à une réunion.
\item \textbf{False.} \eng{Actually} est un \textbf{faux ami} : il signifie « en fait, en réalité, effectivement » (\textit{in fact}), et non « actuellement ». « Actuellement » se dit \eng{currently} ou \eng{at the moment}.
\item \textbf{False.} En programmation, une \eng{library} (bibliothèque) est un ensemble de codes précompilés réutilisables (\textit{collection of pre-written code}), pas un lieu de prêt de livres (\eng{library} = bibliothèque au sens architectural).
\item \textbf{False.} \eng{To download} = « télécharger » (recevoir des données d'un serveur vers son appareil). L'inverse (envoyer vers un serveur) est \eng{to upload}.
\item \textbf{True.} \eng{Smartphone} = téléphone intelligent. \eng{Portable} est un adjectif signifiant « transportable ». \eng{A portable computer} = un ordinateur portable. L'affirmation est correcte.
\end{enumerate}
\end{corrige}

\begin{corrige}[Exercice 3 — Vocabulaire : Familles de mots et collocations]
\begin{center}
\begin{tabularx}{\linewidth}{|X|X|X|X|}
\hline
\rowcolor{popblueL} \textbf{Verbe} & \textbf{Nom} & \textbf{Adjectif / Adverbe} & \textbf{Collocation} \\
\hline
to connect & connection / connectivity & connected / connectedly & \eng{a stable connection} \\
\hline
to inform & \eng{information} & \eng{informative} / \eng{informed} & \eng{provide information}, \eng{an informed decision} \\
\hline
\eng{to compute} & computation / computer & \eng{computational} / \eng{computerized} & \eng{perform a computation}, \eng{a personal computer} \\
\hline
to digitalize / digitize & \eng{digitalization} / \eng{digitization} & digital / digitally & \eng{digital transformation}, \eng{digitally connected} \\
\hline
to browse & \eng{browser} / \eng{browsing} & \eng{browsable} & \eng{browse the web}, \eng{a web browser} \\
\hline
to upload & \eng{upload} / \eng{uploader} & \eng{uploadable} & \eng{upload a file}, \eng{file upload} \\
\hline
to stream & \eng{stream} / \eng{streaming} & \eng{streaming} (adj. inv.) & \eng{stream a video}, \eng{live streaming} \\
\hline
\end{tabularx}
\end{center}
\emph{Notes :} \eng{Compute} est le verbe de base (calculer) ; \eng{computer} est aussi un nom d'agent (celui qui calcule) devenu le nom de la machine. \eng{Streaming} s'utilise comme nom verbal (\eng{live streaming}) et comme adjectif (\eng{a streaming service}).
\end{corrige}

\begin{corrige}[Exercice 4 — Grammaire en contexte : Prépositions et particules]
\begin{enumerate}
\item Please log \textbf{in} using your student ID and password. (\eng{log in} = se connecter)
\item The file is too large to send \textbf{by} email; use a cloud link instead. (\eng{send by email} / \eng{via email})
\item She is very good \textbf{at} troubleshooting network issues. (\eng{good at} + V-ing)
\item We need to back \textbf{up} the database before the update. (\eng{back up} = sauvegarder)
\item The system crashed, so we had to start \textbf{over} all over again. (\eng{start over} = recommencer à zéro)
\item Click \textbf{on} the icon to open the application. (\eng{click on} / \eng{click})
\item Are you connected \textbf{to} the Wi-Fi network? (\eng{connected to})
\item He spends hours browsing \textbf{on} the web for research. (\eng{browse on the web} / \eng{surf on the web} ; \eng{browse the web} sans préposition aussi possible mais le trou attend une préposition de la liste)
\item The software is compatible \textbf{with} both Windows and Linux. (\eng{compatible with})
\item Don't forget to sign \textbf{out} when you leave the computer lab. (\eng{sign out} / \eng{log out} = fermer sa session)
\end{enumerate}
\end{corrige}

\begin{corrige}[Exercice 5 — Traduction : Éviter les calques du français]
\begin{enumerate}
\item \eng{I am attending a videoconference with my colleagues from Douala.} \\
      (Piège : \eng{assist} $\neq$ assister à ; \eng{attend} = assister à / participer à)
\item \eng{Currently, I am downloading the report to the server.} \\
      (Piège : \eng{actually} $\neq$ actuellement ; \eng{currently} / \eng{at the moment} = actuellement. \eng{Download} = télécharger \textit{vers} son appareil ; ici \eng{to the server} suggère un envoi, mais la phrase française dit « télécharge \textit{sur} le serveur » ce qui est un abus de langage fr pour « uploader ». Traduction littérale respectée : \eng{downloading... to the server}. Mieux : \eng{uploading the report to the server}.)
\item \eng{You must back up your data on an external hard drive.} \\
      (\eng{Save} = enregistrer ; \eng{back up} = faire une copie de secours / sauvegarder).
\item \eng{This app is very user-friendly, even for beginners.} \\
      (\eng{Convivial} $\rightarrow$ \eng{user-friendly} ; \eng{application} $\rightarrow$ \eng{app} / \eng{application}).
\item \eng{We are browsing the Internet to find reliable information.} \\
      (\eng{Naviguer sur Internet} = \eng{browse the Internet} / \eng{surf the web} / \eng{go online}).
\item \eng{The trainer shared his screen during the webinar.} \\
      (\eng{Partager son écran} = \eng{share one's screen} ; \eng{formateur} = \eng{trainer} / \eng{instructor}).
\item \eng{I am going to update my antivirus software tonight.} \\
      (\eng{Mettre à jour} = \eng{update} ; \eng{logiciel} = \eng{software} [invariable]).
\item \eng{The attachment is a 5 MB PDF.} \quad ou \quad \eng{The attached file is a 5 MB PDF.} \\
      (\eng{Fichier joint} = \eng{attachment} [nom] ou \eng{attached file}).
\item \eng{She printed the document without checking it.} \\
      (\eng{Sans le vérifier} = \eng{without checking it} [gérondif après préposition]).
\item \eng{Our Internet connection is slow today.} \\
      (\eng{Connexion Internet} = \eng{Internet connection} ; \eng{lent} = \eng{slow}).
\end{enumerate}
\end{corrige}

\begin{corrige}[Exercice 6 — Transformation : Reformulation avec le vocabulaire cible]
\begin{enumerate}
\item She sent the document \textbf{as an attachment}.
\item The \textbf{webinar} starts at 3 p.m.
\item You must \textbf{log in} to access the platform.
\item \textbf{Back up} your project on Google Drive.
\item The interface is \textbf{user-friendly}.
\item He \textbf{uploaded} the photos from his phone to his laptop.
\item They are \textbf{streaming} a live video on YouTube.
\item I need to \textbf{update} the browser (to the latest version).
\item The file was too big, so she \textbf{zipped} it. (\eng{zip} / \eng{compress})
\item We are \textbf{teleworking} / doing \textbf{remote work} this week.
\end{enumerate}
\end{corrige}

\begin{corrige}[Exercice 7 — Texte à trous : Le télétravail au Cameroun]
\begin{textebox}[Gap-filling: Teleworking in Cameroon — Corrigé]
In major Cameroonian cities like Yaoundé and Douala, \textbf{teleworking} has become more common since the pandemic. Many employees now work \textbf{remotely} from home two or three days a week. A stable Internet \textbf{connection} is essential, but power cuts remain a challenge. Workers use video \textbf{conferencing} tools like Zoom or Teams for meetings. They share files via cloud storage and send \textbf{attachments} by email. Companies provide \textbf{company} laptops with pre-installed security software. Employees must \textbf{back up} their work regularly to avoid data loss. Some find it hard to \textbf{balance} professional and personal life. However, most appreciate the flexibility and the reduction in daily \textbf{commute} time.
\end{textebox}
\textbf{Mots non utilisés :} \eng{download}, \eng{offline}, \eng{hardware}.
\end{corrige}

\begin{corrige}[Exercice 8 — Appariement : Collocations verbales]
\begin{center}
\begin{tabular}{|c|l|c|l|}
\hline
\multicolumn{2}{|c|}{\textbf{Colonne A (Verbes)}} & \multicolumn{2}{c|}{\textbf{Colonne B (Noms)}} \\
\hline
1 & to browse & b & the web \\
2 & to download & a & a file \\
3 & to log & d & in / out \\
4 & to attend & c & a meeting \\
5 & to upload & e & a document \\
6 & to print & g & a screenshot \\
7 & to reset & f & a password \\
8 & to take & g & a screenshot \\
9 & to install & h & a program \\
10 & to fix & i & a bug \\
\hline
\end{tabular}
\end{center}
\textbf{Paires :} 1–b, 2–a, 3–d, 4–c, 5–e, 6–g, 7–f, 8–g, 9–h, 10–i. \\
(Note : \eng{take a screenshot} et \eng{print a screenshot} sont tous deux corrects ; ici \eng{take} $\rightarrow$ \eng{screenshot}, \eng{print} $\rightarrow$ \eng{screenshot} aussi possible, mais \eng{print a document} n'est pas dispo. On associe donc \eng{print} à \eng{a screenshot} pour l'exercice).
\end{corrige}

\begin{corrige}[Exercice 9 — Compréhension de texte et langue : Type Bac]
\textbf{Comprehension questions (suggested answers):}
\begin{enumerate}
\item The school received fifty tablets and a solar-powered router (powered by solar energy).
\item Students access offline educational resources (videos, quizzes, e-books) stored on a local server; they download them during computer lab sessions.
\item Improved attendance and better results in science subjects. (Also accepted: access to updated resources, interactive learning).
\item Device maintenance, limited bandwidth for updates, and the need for continuous teacher training.
\item The Ministry plans to scale up the pilot project to ten more schools by 2026.
\end{enumerate}

\textbf{Language work:}
\begin{enumerate}
\setcounter{enumi}{5}
\item \textbf{maintenance} (paragraph 2, « device \eng{maintenance} »).
\item \textbf{scale up} (last sentence, « \eng{scale up} the pilot »).
\item \textbf{up-to-date} / \textbf{current} / \textbf{modern} (opposite of \eng{outdated}).
\item \textbf{Weekly lessons are uploaded to a local server by teachers.} (Passive: \eng{are uploaded}).
\item Students are interested \textbf{in} digital learning. / The tablet is connected \textbf{to} the local server.
\end{enumerate}

\textbf{Barème indicatif (sur 10 pts) :}
\begin{itemize}
\item Comprehension (Q1–Q5) : 5 pts (1 pt par réponse correcte et formulée en anglais).
\item Language work (Q6–Q10) : 5 pts (1 pt par item).
\end{itemize}
\textbf{Points de vigilance :} Réponses en anglais, phrases complètes, vocabulaire du texte réemployé, voix passive bien formée (\eng{be} + V-en), prépositions \eng{interested in}, \eng{connected to}.
\end{corrige}

\begin{corrige}[Exercice 10 — Traduction : Type Bac — Thème]
\textbf{Proposition de traduction :}
\begin{quote}
\eng{In Cameroonian high schools, the integration of ICTs is gradually transforming teaching methods. Interactive whiteboards are replacing traditional blackboards. Students use tablets to access digital textbooks and do online exercises. Teachers, trained in digital tools, create interactive content and organize virtual classes. Despite power cuts and sometimes slow connectivity, enthusiasm remains strong. The Ministry plans to equip more schools by 2027.}
\end{quote}

\textbf{Points de vocabulaire et structures attendus :}
\begin{itemize}
\item \eng{ICTs} (TIC), \eng{integration}, \eng{gradually} (peu à peu), \eng{teaching methods}.
\item \eng{Interactive whiteboards} (tableaux interactifs), \eng{replace} (remplacent), \eng{traditional blackboards}.
\item \eng{Tablets}, \eng{access} / \eng{consult} (consulter), \eng{digital textbooks} (manuels numériques), \eng{online exercises}.
\item \eng{Trained in} (formés à), \eng{digital tools}, \eng{create interactive content}, \eng{virtual classes} (classes virtuelles).
\item \eng{Power cuts} (coupures d'électricité), \eng{connectivity} / \eng{Internet connection} (connexion), \eng{slow}.
\item \eng{Enthusiasm remains strong}, \eng{The Ministry plans to equip} (prévoit d'équiper), \eng{more schools}, \eng{by 2027} (d'ici 2027).
\end{itemize}

\textbf{Barème indicatif (sur 10 pts) :}
\begin{itemize}
\item Respect du sens et du registre : 3 pts.
\item Vocabulaire de la leçon (TIC, apprentissage, travail) correctement employé : 3 pts.
\item Grammaire (temps, voix passive, prépositions, ordre des mots) : 2 pts.
\item Orthographe et ponctuation : 2 pts.
\end{itemize}
\end{corrige}

\begin{corrige}[Exercice 11 — Expression écrite guidée : Type Bac — Composition]
\textbf{Sujet :} \eng{“How ICTs can improve learning in Cameroonian schools”}

\textbf{Proposition de rédaction modèle (environ 92 mots) :}
\begin{quote}
\eng{ICTs can significantly improve learning in Cameroonian schools. First, tablets allow students to download educational resources and study offline, which is crucial in areas with low bandwidth. Moreover, interactive whiteboards make lessons more user-friendly and engaging. For instance, teachers can stream live experiments during a webinar. In addition, virtual classes enable remote learning during crises. However, schools must back up data regularly and ensure device maintenance. Therefore, if the Ministry scales up digital equipment and teacher training, ICTs will transform education nationwide.}
\end{quote}

\textbf{Vérification des consignes :}
\begin{itemize}
\item \textbf{Mots du vocabulaire (8) :} \eng{tablets, download, offline, bandwidth, user-friendly, stream, webinar, back up, maintenance, scale up, virtual classes, remote learning, interactive}. ($\ge$ 6 requis)
\item \textbf{Connecteurs logiques (5) :} \eng{First, Moreover, For instance, In addition, However, Therefore}. ($\ge$ 3 requis)
\item \textbf{Structure :} Phrase d'introduction (\eng{ICTs can significantly improve...}), arguments illustrés (\eng{First... Moreover... For instance... In addition...}), concession/contre-argument (\eng{However...}), conclusion avec conséquence (\eng{Therefore...}).
\item \textbf{Nombre de mots :} 92 mots (dans la fourchette 80–100).
\end{itemize}

\textbf{Barème indicatif (sur 10 pts) :}
\begin{itemize}
\item Respect du nombre de mots (80–100) : 1 pt.
\item Utilisation de $\ge$ 6 mots/expressions du vocabulaire cible : 2 pts.
\item Utilisation de $\ge$ 3 connecteurs logiques appropriés : 2 pts.
\item Structure cohérente (intro, arguments, conclusion) : 2 pts.
\item Correction grammaticale, orthographe, ponctuation : 3 pts.
\end{itemize}
\end{corrige}

\newpage

\coursec{Fiche bilan}

\coursec{Lesson 54 — Vocabulary: Words and Expressions Related to Using ICTs to Enhance Work or Learning}

\begin{popfigure}
\begin{center}\tcbox[colback=popink,colframe=popink,coltext=white,arc=9pt,boxsep=4pt,left=6pt,right=6pt]{\bfseries\small ICTs for Work \& Learning}\end{center}
\vspace{-2pt}
\begin{tcbraster}[raster columns=3,raster equal height=rows,raster column skip=6pt,raster row skip=6pt,enhanced,blankest]
\begin{tcolorbox}[enhanced,colback=white,colframe=popblue,boxrule=0.9pt,arc=3pt,title={Devices \& Hardware},fonttitle=\bfseries\scriptsize,coltitle=white,colbacktitle=popblue,left=4pt,right=4pt,top=2pt,bottom=2pt,boxsep=1pt,toptitle=1.5pt,bottomtitle=1.5pt,halign=left]
\begin{itemize}[nosep,leftmargin=*,itemsep=1pt,font=\scriptsize]
\item laptop, smartphone
\item tablet, desktop
\item router, modem
\item webcam, headset
\end{itemize}
\end{tcolorbox}
\begin{tcolorbox}[enhanced,colback=white,colframe=poporange,boxrule=0.9pt,arc=3pt,title={Software \& Apps},fonttitle=\bfseries\scriptsize,coltitle=white,colbacktitle=poporange,left=4pt,right=4pt,top=2pt,bottom=2pt,boxsep=1pt,toptitle=1.5pt,bottomtitle=1.5pt,halign=left]
\begin{itemize}[nosep,leftmargin=*,itemsep=1pt,font=\scriptsize]
\item OS, browser
\item word processor
\item spreadsheet
\item video conferencing
\item cloud storage
\end{itemize}
\end{tcolorbox}
\begin{tcolorbox}[enhanced,colback=white,colframe=popgreen,boxrule=0.9pt,arc=3pt,title={Internet \& Connectivity},fonttitle=\bfseries\scriptsize,coltitle=white,colbacktitle=popgreen,left=4pt,right=4pt,top=2pt,bottom=2pt,boxsep=1pt,toptitle=1.5pt,bottomtitle=1.5pt,halign=left]
\begin{itemize}[nosep,leftmargin=*,itemsep=1pt,font=\scriptsize]
\item broadband, Wi-Fi
\item bandwidth, 4G/5G
\item URL, hyperlink
\item search engine
\item firewall, VPN
\end{itemize}
\end{tcolorbox}
\begin{tcolorbox}[enhanced,colback=white,colframe=poppurple,boxrule=0.9pt,arc=3pt,title={E-Learning \& Education},fonttitle=\bfseries\scriptsize,coltitle=white,colbacktitle=poppurple,left=4pt,right=4pt,top=2pt,bottom=2pt,boxsep=1pt,toptitle=1.5pt,bottomtitle=1.5pt,halign=left]
\begin{itemize}[nosep,leftmargin=*,itemsep=1pt,font=\scriptsize]
\item LMS (Moodle)
\item MOOC, webinar
\item e-portfolio
\item digital literacy
\item blended learning
\end{itemize}
\end{tcolorbox}
\begin{tcolorbox}[enhanced,colback=white,colframe=popgold,boxrule=0.9pt,arc=3pt,title={Remote Work \& Productivity},fonttitle=\bfseries\scriptsize,coltitle=white,colbacktitle=popgold,left=4pt,right=4pt,top=2pt,bottom=2pt,boxsep=1pt,toptitle=1.5pt,bottomtitle=1.5pt,halign=left]
\begin{itemize}[nosep,leftmargin=*,itemsep=1pt,font=\scriptsize]
\item telecommuting
\item virtual meeting
\item project management
\item screen sharing
\item deadline, task
\end{itemize}
\end{tcolorbox}
\begin{tcolorbox}[enhanced,colback=white,colframe=poppink,boxrule=0.9pt,arc=3pt,title={Digital Skills \& Safety},fonttitle=\bfseries\scriptsize,coltitle=white,colbacktitle=poppink,left=4pt,right=4pt,top=2pt,bottom=2pt,boxsep=1pt,toptitle=1.5pt,bottomtitle=1.5pt,halign=left]
\begin{itemize}[nosep,leftmargin=*,itemsep=1pt,font=\scriptsize]
\item cybersecurity
\item phishing, malware
\item password manager
\item netiquette
\item data privacy
\end{itemize}
\end{tcolorbox}
\begin{tcolorbox}[enhanced,colback=white,colframe=popblue,boxrule=0.9pt,arc=3pt,title={Collocations \& Word Families},fonttitle=\bfseries\scriptsize,coltitle=white,colbacktitle=popblue,left=4pt,right=4pt,top=2pt,bottom=2pt,boxsep=1pt,toptitle=1.5pt,bottomtitle=1.5pt,halign=left]
\begin{itemize}[nosep,leftmargin=*,itemsep=1pt,font=\scriptsize]
\item log in / log out
\item sign up for
\item download / upload
\item troubleshoot
\item user-friendly
\end{itemize}
\end{tcolorbox}
\end{tcbraster}

\legende{Carte mentale : Vocabulaire des TIC pour le travail et l'apprentissage (Leçon 54)}
\end{popfigure}

\courssub{Discovering the Theme: ICTs Around Us}

\begin{textebox}[Authentic context — A student's digital day in Yaoundé]
Marie, a final-year student at a bilingual high school in Yaoundé, starts her day checking her \eng{smartphone} for notifications. She connects to the school \eng{LMS} (Moodle) to \eng{download} the latest assignment for her \eng{blended learning} course. Because the campus \eng{Wi-Fi} has low \eng{bandwidth}, she uses her mobile \eng{4G} hotspot. Later, she joins a \eng{webinar} on \eng{digital literacy} via a \eng{video conferencing tool}. For her group project, the team uses a \eng{project management tool} (Trello) and \eng{collaborative editing} on Google Docs. Before submitting, she runs an \eng{antivirus} scan, enables \eng{2FA} on her accounts, and backs up her \eng{e-portfolio} to \eng{cloud storage}. She is careful about \eng{phishing} emails and manages her \eng{digital footprint}.
\end{textebox}

\courssub{Lexical Field 1: Devices, Tools and the Internet}

\textbf{1. Devices \& Hardware (Appareils \& Matériel)}
\begin{center}
\begin{tabularx}{\linewidth}{|X|X|X|}\hline
\rowcolor{popblueL} \textbf{English} & \textbf{Français} & \textbf{Exemple / Contexte} \\ \hline
laptop / notebook & ordinateur portable & \eng{I carry my laptop to class every day.} « J'apporte mon portable en classe chaque jour. » \\ \hline
smartphone & smartphone & \eng{She checks her smartphone for emails.} « Elle consulte ses mails sur son smartphone. » \\ \hline
tablet & tablette tactile & \eng{The tablet is handy for reading PDFs.} « La tablette est pratique pour lire des PDF. » \\ \hline
desktop (computer) & ordinateur de bureau & \eng{The school lab has 20 desktops.} « Le labo de l'école a 20 fixes. » \\ \hline
router / modem & routeur / modem & \eng{Reset the router if the Wi-Fi fails.} « Redémarrez le routeur si le Wi-Fi coupe. » \\ \hline
webcam & webcam & \eng{Turn on your webcam for the meeting.} « Allumez votre webcam pour la réunion. » \\ \hline
headset / headphones & casque / écouteurs & \eng{A good headset reduces background noise.} « Un bon casque réduit le bruit ambiant. » \\ \hline
external drive / USB stick & disque dur externe / clé USB & \eng{Back up your files on an external drive.} « Sauvegardez vos fichiers sur un disque externe. » \\ \hline
\end{tabularx}
\end{center}

\textbf{2. Software, Apps \& Internet (Logiciels, Applis \& Internet)}
\begin{center}
\begin{tabularx}{\linewidth}{|X|X|X|}\hline
\rowcolor{poporangeL} \textbf{English} & \textbf{Français} & \textbf{Exemple / Contexte} \\ \hline
operating system (OS) & système d'exploitation & \eng{Update your OS regularly for security.} « Mettez à jour votre OS régulièrement pour la sécurité. » \\ \hline
web browser & navigateur web & \eng{Chrome and Firefox are popular browsers.} « Chrome et Firefox sont des navigateurs populaires. » \\ \hline
word processor & traitement de texte & \eng{Write the report in a word processor.} « Rédigez le rapport dans un traitement de texte. » \\ \hline
spreadsheet & tableur & \eng{Use a spreadsheet for the budget.} « Utilisez un tableur pour le budget. » \\ \hline
video conferencing tool & outil de visioconférence & \eng{We use Zoom for video conferencing.} « Nous utilisons Zoom pour la visio. » \\ \hline
cloud storage & stockage en ligne (cloud) & \eng{Save it to the cloud (Google Drive).} « Enregistrez-le dans le cloud (Google Drive). » \\ \hline
broadband / high-speed Internet & haut débit & \eng{Rural areas need broadband access.} « Les zones rurales ont besoin d'un accès au haut débit. » \\ \hline
bandwidth & bande passante & \eng{Low bandwidth slows the download.} « Une bande passante faible ralentit le téléchargement. » \\ \hline
search engine & moteur de recherche & \eng{Google is a popular search engine.} « Google est un moteur de recherche populaire. » \\ \hline
URL / link / hyperlink & URL / lien / lien hypertexte & \eng{Click the hyperlink to open the page.} « Cliquez sur le lien hypertexte pour ouvrir la page. » \\ \hline
firewall / antivirus & pare-feu / antivirus & \eng{Keep your antivirus updated.} « Maintenez votre antivirus à jour. » \\ \hline
VPN (Virtual Private Network) & VPN (Réseau privé virtuel) & \eng{Use a VPN on public Wi-Fi.} « Utilisez un VPN sur le Wi-Fi public. » \\ \hline
\end{tabularx}
\end{center}

\courssub{Lexical Field 2: ICTs for Learning (E-Education)}

\begin{center}
\begin{tabularx}{\linewidth}{|X|X|X|}\hline
\rowcolor{popgreenL} \textbf{English} & \textbf{Français} & \textbf{Exemple / Contexte} \\ \hline
LMS (Learning Management System) & plateforme d'apprentissage (ex. Moodle) & \eng{Submit your assignment on the LMS before midnight.} « Déposez votre devoir sur la plateforme avant minuit. » \\ \hline
MOOC (Massive Open Online Course) & CLOM / Cours en ligne ouvert et massif & \eng{She enrolled in a MOOC on Python programming.} « Elle s'est inscrite à un CLOM sur la programmation Python. » \\ \hline
webinar & webinaire (séminaire en ligne) & \eng{The webinar starts at 3 p.m. sharp.} « Le webinaire commence à 15 h précises. » \\ \hline
e-portfolio & portfolio numérique & \eng{Build your e-portfolio for university applications.} « Construisez votre portfolio numérique pour les admissions à l'université. » \\ \hline
digital literacy & culture / littératie numérique & \eng{Digital literacy is a key skill for the 21st century.} « La littératie numérique est une compétence clé du 21e siècle. » \\ \hline
blended learning & apprentissage hybride (mixte) & \eng{Our school uses a blended learning approach.} « Notre école utilise une approche d'apprentissage hybride. » \\ \hline
online assessment & évaluation en ligne & \eng{The quiz is an online assessment on the LMS.} « Le quiz est une évaluation en ligne sur la plateforme. » \\ \hline
virtual classroom & classe virtuelle & \eng{Join the virtual classroom on Microsoft Teams.} « Rejoignez la classe virtuelle sur Microsoft Teams. » \\ \hline
\end{tabularx}
\end{center}

\courssub{Lexical Field 3: ICTs for Work (Remote Work \& Productivity)}

\begin{center}
\begin{tabularx}{\linewidth}{|X|X|X|}\hline
\rowcolor{poppurpleL} \textbf{English} & \textbf{Français} & \textbf{Exemple / Contexte} \\ \hline
telecommuting / remote work & télétravail & \eng{Telecommuting has become common since 2020.} « Le télétravail est devenu courant depuis 2020. » \\ \hline
virtual meeting / call & réunion virtuelle / appel & \eng{Schedule a virtual meeting for Monday morning.} « Planifiez une réunion virtuelle pour lundi matin. » \\ \hline
project management tool & outil de gestion de projet & \eng{Trello is a popular project management tool.} « Trello est un outil de gestion de projet populaire. » \\ \hline
screen sharing & partage d'écran & \eng{Enable screen sharing to present the slides.} « Activez le partage d'écran pour présenter les diapositives. » \\ \hline
deadline & date limite / échéance & \eng{We must meet the deadline for the project.} « Nous devons respecter l'échéance du projet. » \\ \hline
to-do list & liste de tâches & \eng{Check off items on your to-do list daily.} « Cochez les éléments de votre liste de tâches chaque jour. » \\ \hline
workflow & flux de travail & \eng{Automate your workflow with scripts.} « Automatisez votre flux de travail avec des scripts. » \\ \hline
collaborative editing & édition collaborative & \eng{Google Docs allows real-time collaborative editing.} « Google Docs permet l'édition collaborative en temps réel. » \\ \hline
\end{tabularx}
\end{center}

\courssub{Digital Safety and Netiquette}

\begin{center}
\begin{tabularx}{\linewidth}{|X|X|X|}\hline
\rowcolor{poppinkL} \textbf{English} & \textbf{Français} & \textbf{Exemple / Contexte} \\ \hline
cybersecurity & cybersécurité & \eng{Cybersecurity awareness is vital for everyone.} « La sensibilisation à la cybersécurité est vitale pour tous. » \\ \hline
phishing & hameçonnage & \eng{Don't click suspicious links; it might be phishing.} « Ne cliquez pas sur des liens suspects ; c'est peut-être de l'hameçonnage. » \\ \hline
malware / virus & logiciel malveillant / virus & \eng{Scan your computer for malware regularly.} « Analysez votre ordinateur régulièrement pour détecter les malwares. » \\ \hline
password manager & gestionnaire de mots de passe & \eng{Use a password manager to generate strong passwords.} « Utilisez un gestionnaire de mots de passe pour générer des mots de passe robustes. » \\ \hline
two-factor authentication (2FA) & double authentification (2FA) & \eng{Enable 2FA on all your important accounts.} « Activez la 2FA sur tous vos comptes importants. » \\ \hline
data privacy / protection & confidentialité / protection des données & \eng{Read the data privacy policy before accepting.} « Lisez la politique de confidentialité des données avant d'accepter. » \\ \hline
netiquette & nétiquette (règles de conduite en ligne) & \eng{Respect netiquette: avoid typing in ALL CAPS.} « Respectez la nétiquette : évitez d'écrire EN MAJUSCULES. » \\ \hline
digital footprint & empreinte numérique & \eng{Be careful with your digital footprint on social media.} « Attention à votre empreinte numérique sur les réseaux sociaux. » \\ \hline
\end{tabularx}
\end{center}

\courssub{Collocations, Word Families and False Friends}

\begin{propriete}[Collocations usuelles (Verbe + Particule / Adjectif + Nom)]
\begin{itemize}
\item \eng{log in / log on} (se connecter à un système) $\neq$ \eng{log out / log off} (se déconnecter)
\item \eng{sign up for} (s'inscrire à un service/cours) / \eng{sign in} (se connecter / ouvrir une session)
\item \eng{download} (télécharger \emph{vers} son appareil local) / \eng{upload} (téléverser \emph{vers} un serveur/cloud)
\item \eng{install / uninstall} (installer / désinstaller) + \eng{update} (mettre à jour)
\item \eng{troubleshoot} (v., dépanner, résoudre un problème technique) $\rightarrow$ \eng{troubleshooting} (n., le dépannage)
\item \eng{back up} (v., sauvegarder) / \eng{a backup} (n., une sauvegarde) — \emph{Attention :} \eng{back up} (deux mots verbe), \eng{backup} (un mot nom/adjectif).
\item \eng{scroll up / down} (faire défiler vers le haut / bas)
\item \eng{plug in / unplug} (brancher / débrancher)
\item \eng{high-speed / broadband connection} (connexion haut débit)
\item \eng{user-friendly interface} (interface conviviale) / \eng{intuitive design} (conception intuitive)
\end{itemize}
\end{propriete}

\begin{propriete}[Familles de mots (Dérivation) — Orthographe US de référence]
\begin{center}
\begin{tabularx}{\linewidth}{|X|X|X|X|}\hline
\rowcolor{popblueL} \textbf{Verbe} & \textbf{Nom (chose/action)} & \textbf{Nom (agent)} & \textbf{Adjectif} \\ \hline
compute & computer, computation & computer scientist & computational \\ \hline
program & program, programming & programmer & programmable \\ \hline
digitize & digitization & -- & digital \\ \hline
connect & connection, connectivity & -- & connected / connective \\ \hline
secure & security & -- & secure \\ \hline
encrypt & encryption & -- & encrypted \\ \hline
automate & automation & -- & automatic / automated \\ \hline
innovate & innovation & innovator & innovative \\ \hline
\end{tabularx}
\end{center}
\end{propriete}

\begin{attention}[Faux amis \& Pièges fréquents (Francophones)]
\begin{itemize}
\item \cle{Ordinateur} = \eng{computer} (pas \eng{ordinateur} qui n'existe pas en anglais). \eng{Ordinator} n'existe pas.
\item \cle{Logiciel} = \eng{software} (indénombrable, pas \eng{a software} $\rightarrow$ \eng{a program / an application / an app}).
\item \cle{Télécharger} : distinction cruciale.
    \begin{itemize}
    \item Vers son appareil (réception) = \eng{download}
    \item Vers le net/serveur (envoi) = \eng{upload}
    \end{itemize}
\item \cle{Clé USB} = \eng{USB stick / flash drive / thumb drive} (pas \eng{USB key} : \eng{key} = clé de licence ou touche de clavier).
\item \cle{Écran} = \eng{screen / monitor / display} (pas \eng{écran}).
\item \cle{Souris} = \eng{mouse} (pluriel \eng{mice}) / \eng{trackpad} (pavé tactile).
\item \cle{Navigateur web} = \eng{browser} (pas \eng{navigator} qui désigne un navigateur humain ou un GPS).
\item \cle{Site web} = \eng{website / site} (pas \eng{web site} deux mots, forme démodée).
\item \cle{Adresse e-mail} = \eng{email address} (pas \eng{mail address} seul). \eng{Mail} = courrier postal (US) ou e-mail (UK) ; \eng{Email} (standard international).
\item \cle{Pièce jointe} = \eng{attachment} (pas \eng{attached file} seulement, ni \eng{join} $\rightarrow$ \eng{attach a file}).
\item \cle{Mot de passe} = \eng{password} (un seul mot, pas \eng{pass word}).
\item \cle{Se connecter} = \eng{log in / sign in} (pas \eng{connect oneself}).
\item \cle{Supprimer (fichier)} = \eng{delete / remove / erase} (pas \eng{suppress} qui signifie réprimer, censurer, étouffer).
\item \cle{Enregistrer (fichier)} = \eng{save} (pas \eng{record} = enregistrer audio/vidéo/son).
\item \cle{Imprimer} = \eng{print} (pas \eng{impress} = impressionner).
\end{itemize}
\end{attention}

\courssub{Using the Vocabulary in Context}

\begin{exercice}[Entraînement Bac — Vocabulaire ICT en contexte]
\textbf{A. Complétez les phrases avec le mot juste (forme correcte).}
\begin{enumerate}
\item The \_\_\_\_\_\_\_\_\_\_ (bandwidth / broadband) in this rural area is too low for video calls.
\item Please \_\_\_\_\_\_\_\_\_\_ (upload / download) the completed form to the server.
\item She forgot her \_\_\_\_\_\_\_\_\_\_ (password / pass word) and couldn't \_\_\_\_\_\_\_\_\_\_ (log in / connect).
\item The \_\_\_\_\_\_\_\_\_\_ (software / program) requires a system \_\_\_\_\_\_\_\_\_\_ (update / up date).
\item Be careful of \_\_\_\_\_\_\_\_\_\_ (phishing / fishing) emails asking for bank details.
\end{enumerate}

\textbf{B. Remplacez le verbe générique par un verbe précis du domaine ICT.}
\begin{enumerate}
\item The technician \textbf{repaired} the network issue. $\rightarrow$ The technician \_\_\_\_\_\_\_\_\_\_ the network issue.
\item We need to \textbf{make} a backup of the database. $\rightarrow$ We need to \_\_\_\_\_\_\_\_\_\_ the database.
\item The student \textbf{put} the file on the cloud. $\rightarrow$ The student \_\_\_\_\_\_\_\_\_\_ the file to the cloud.
\item Please \textbf{look at} the attachment. $\rightarrow$ Please \_\_\_\_\_\_\_\_\_\_ the attachment.
\item Hackers \textbf{entered} the system illegally. $\rightarrow$ Hackers \_\_\_\_\_\_\_\_\_\_ the system.
\end{enumerate}

\textbf{C. Expression écrite (50--80 mots).}
Vous êtes délégué(e) de classe. Rédigez un email formel au proviseur pour demander l'amélioration de la connexion Internet au CDI. Utilisez au moins 5 termes du vocabulaire de la leçon (ex. \eng{bandwidth, broadband, router, LMS, download, deadline}).
\end{exercice}

\begin{corrige}[Exercice Entraînement Bac]
\textbf{A.}
\begin{enumerate}
\item \eng{bandwidth} (bande passante = mesure technique ; \eng{broadband} = type d'accès haut débit).
\item \eng{upload} (envoi vers serveur).
\item \eng{password} (un mot) / \eng{log in} (verbe à particule).
\item \eng{software} (indénombrable) / \eng{update} (nom/verbe un mot).
\item \eng{phishing} (hameçonnage ; \eng{fishing} = pêche).
\end{enumerate}
\textbf{B.}
\begin{enumerate}
\item \eng{troubleshot} (prétérit de \eng{troubleshoot}) ou \eng{resolved} / \eng{fixed}.
\item \eng{back up} (verbe à deux mots).
\item \eng{uploaded} (vers le cloud/serveur).
\item \eng{check} / \eng{review} / \eng{open}.
\item \eng{breached} / \eng{hacked} / \eng{compromised}.
\end{enumerate}
\textbf{C. Modèle de réponse :}
\eng{Subject: Request for Improved Internet Connectivity in the Library}
\eng{Dear Principal,}
\eng{I am writing on behalf of my classmates regarding the poor Internet connection in the school library. The current \textbf{bandwidth} is insufficient; students cannot \textbf{download} resources from the \textbf{LMS} or attend \textbf{webinars} before \textbf{deadlines}. The \textbf{router} often disconnects. We urgently need a \textbf{broadband} upgrade to support our \textbf{digital literacy} and research work.}
\eng{Yours faithfully,}
\eng{[Name], Class Representative}
\end{corrige}

\courssub{Methods and Exam Tips}

\begin{methode}[Apprendre et réemployer le vocabulaire ICT efficacement]
\begin{enumerate}
\item \textbf{Contexte d'abord :} Ne pas apprendre des listes isolées. Lire un article tech (BBC Tech, The Verge, \eng{TechCrunch}) ou regarder un tutoriel YouTube en anglais (ex. \eng{"How to set up a VPN"}).
\item \textbf{Cartes mentales (Mindmaps) :} Regrouper par \emph{champ lexical} (Hardware, Software, Safety, Verbs). Utiliser des couleurs (cf. carte en début de leçon).
\item \textbf{Collocations :} Apprendre les \textbf{verbes + prépositions/particules} (\eng{log \textbf{in}}, \eng{sign \textbf{up} \textbf{for}}, \eng{back \textbf{up}}).
\item \textbf{Familles de mots :} Dériver systématiquement (V $\rightarrow$ N $\rightarrow$ Adj) : \eng{secure $\rightarrow$ security $\rightarrow$ secure}.
\item \textbf{Production active :} Écrire 5 phrases personnelles par semaine utilisant 3 nouveaux mots. S'enregistrer à l'oral (Speaking).
\item \textbf{Révision espacée (SRS) :} Utiliser Anki / Quizlet avec des phrases complètes (\eng{cloze deletion}).
\end{enumerate}
\end{methode}

\begin{methode}[Réussir l'épreuve de Vocabulaire au Bac (Compréhension / Expression)]
\begin{enumerate}
\item \textbf{Reading/Listening :} Repérer les \emph{indices de contexte} (définition, exemple, contraste, synonyme) pour deviner le sens d'un mot ICT inconnu.
\item \textbf{Writing (Essay/Email/Report) :} Varier le lexique : éviter \eng{use / make / do}. Préférer \eng{operate, utilize, leverage, conduct, perform, execute, troubleshoot, implement, deploy, configure}.
\item \textbf{Grammar check -- Indénombrables :} \eng{Software, hardware, equipment, information, advice, bandwidth, malware, phishing} sont \textbf{indénombrables} (pas de \eng{a}, pas de \eng{-s} au pluriel). \eng{Data} est souvent traité comme singulier aujourd'hui (\eng{The data shows...}) mais pluriel formel (\eng{The data show...}).
\item \textbf{Orthographe US/UK :} Cohérence obligatoire. \eng{Program (US/UK tech)} vs \eng{Programme (UK général)}. \eng{Analyze (US)} / \eng{Analyse (UK)}. \eng{Digitize (US)} / \eng{Digitise (UK)}. Choisir une norme et la garder dans toute la copie.
\end{enumerate}
\end{methode}

\begin{aretenir}[Checklist avant le Bac — Leçon 54 : ICT Vocabulary]
$\square$ Je connais le vocabulaire de base : \textbf{hardware} (laptop, router, headset...), \textbf{software} (OS, browser, spreadsheet, cloud...), \textbf{connectivity} (broadband, bandwidth, URL, VPN...).
$\square$ Je maîtrise le lexique de l'\textbf{e-learning} : LMS, MOOC, webinar, blended learning, digital literacy, virtual classroom.
$\square$ Je maîtrise le lexique du \textbf{télétravail} : telecommuting, virtual meeting, project management tool, screen sharing, deadline, collaborative editing.
$\square$ Je connais le vocabulaire de la \textbf{cybersécurité} : phishing, malware, 2FA, password manager, data privacy, digital footprint, netiquette.
$\square$ J'utilise correctement les \textbf{collocations clés} : \eng{log in/out, sign up for, download/upload, back up, troubleshoot, scroll up/down}.
$\square$ Je forme les \textbf{familles de mots} sans erreur : \eng{compute/computer/computational, secure/security/secure, automate/automation/automatic, innovate/innovation/innovative}.
$\square$ J'évite les \textbf{faux amis} majeurs : \eng{software} (indénombrable), \eng{download/upload} (distinction), \eng{browser} (pas navigator), \eng{attachment} (pas join), \eng{save} (pas record), \eng{delete} (pas suppress).
$\square$ Je respecte la \textbf{nature grammaticale} : \eng{equipment, software, hardware, information, bandwidth} = singulier (pas de \eng{a}, pas de \eng{s}).
$\square$ Je sais \textbf{réécrire une phrase} en remplaçant un verbe générique (\eng{do, make, use}) par un verbe précis ICT (\eng{execute, implement, troubleshoot, configure, deploy}).
$\square$ Je suis capable de \textbf{rédiger un email formel} ou un \textbf{court rapport} intégrant 5 à 7 termes techniques du thème dans un contexte camerounais (ex. \eng{internet access in rural areas, using Moodle at university, cybersecurity in mobile money}).
$\square$ Je peux \textbf{présenter oralement} (2 min) un outil numérique que j'utilise pour mes études, en structurant : Name $\rightarrow$ Purpose $\rightarrow$ Key features $\rightarrow$ Pros/Cons $\rightarrow$ Recommendation.
$\square$ Je connais la \textbf{prononciation approximative} des mots pièges : \eng{router} (« raou-teur » US / « rou-teur » UK), \eng{schedule} (« che-dule » UK / « ske-dule » US), \eng{data} (« déi-ta »), \eng{malware} (« mal-oueur »), \eng{phishing} (« fi-ching »).
\end{aretenir}

\findecours

\end{document}
